\part{Investigations}\label{cr:part:investigations}

\chapter{Investigations}\label{cr:ch:investigations}

This chapter holds the investigation reports. Each report keeps its original text.

% Converted from docs/ASAP7_CPU_SYNLIG_MA7.md by tools/docs/md2tex.py; edit this file, not the original.
\section{\texorpdfstring{ASAP7 CPU macro netlist vs RTL differential
(beads \code{ma7} / \code{u99} /
\code{dud4})}{ASAP7 CPU macro netlist vs RTL differential (beads  /  / )}}\label{cr:defect:asap7-cpu-synlig-ma7}

\begin{record}

Date 2026-10-10. Question: is the ASAP7 \code{rv32i_cpu_top} miscompile
proven in \code{u99}/\code{ma7} (\code{DBG_GPR[4]} = \code{0x30} instead
of \code{0x0c}; \texttt{BEQ\ x2,\allowbreak{}x3} with 9 == 9 not taken)
the same defect that \code{dud4} root-caused on Sky130 (Synlig
elaborates two hazard-unit port connections as \texttt{5'x}), or does
ASAP7 have a second cause (ma7\textquotesingle s "\code{OPT_MUXTREE}
mis-elaborates the runtime-indexed regfile read mux")? Test only: no PD
run, no config change, no macro view touched. Statements are tagged
MEASURED or INFERRED.

\end{record}

\subsection{Verdict}\label{cr:defect:asap7-cpu-synlig-ma7:verdict}

\begin{record}

\textbf{ma7 = dud4 (same defect), for the programs and netlists
covered.} The two netlists that carry the two \texttt{5'x} hazard-unit
connections (the committed macro and a fresh Synlig synthesis of main
RTL) reproduce the u99/ma7 failure exactly, in two independent
simulators. The two netlists that do not (Synlig with the two
part-selects routed through a named wire; and the sv2v frontend) are
identical to RTL on all 6 programs and pass all 124 cells of the
31-register x (rs1, rs2, both, debug) matrix. \textbf{No second cause
was found and the ma7 regfile read-mux / \code{OPT_MUXTREE} hypothesis
is not supported}: the regfile read path is correct in both clean
netlists, and \code{OPT_MUXTREE} removal counts and the regfile
\texttt{\$mem} structure are identical across all three synthesised arms
(PD agent\textquotesingle s \code{NOTES.md}).

What this does NOT establish: formal equivalence (6 programs, not
exhaustive); that the committed macro is bit-identical in behaviour to
arm S on programs not run; the GPU macro; anything about P\&R-introduced
faults beyond "the committed post-P\&R netlist behaves like the
synthesis-only Synlig netlist on all 6 programs".

\end{record}

\subsection{Provenance (MEASURED)}\label{cr:defect:asap7-cpu-synlig-ma7:provenance-measured}

\begin{record}

Netlists from the PD agent\textquotesingle s synthesis-only experiment,
\code{/nobackup/claude_sim_build/dud4/asap7/} (method, configs, census
in its \code{NOTES.md}). sha256 re-verified before use:

{\footnotesize\setlength{\tabcolsep}{3pt}

{
\begin{longtable}[]{@{}
  >{\raggedright\arraybackslash}p{\dimexpr(\linewidth-6\tabcolsep)*55/1000\relax}
  >{\raggedright\arraybackslash}p{\dimexpr(\linewidth-6\tabcolsep)*308/1000\relax}
  >{\raggedright\arraybackslash}p{\dimexpr(\linewidth-6\tabcolsep)*308/1000\relax}
  >{\raggedright\arraybackslash}p{\dimexpr(\linewidth-6\tabcolsep)*308/1000\relax}@{}}
\toprule\noalign{}
\begin{minipage}[b]{\linewidth}\raggedright
arm
\end{minipage} & \begin{minipage}[b]{\linewidth}\raggedright
netlist
\end{minipage} & \begin{minipage}[b]{\linewidth}\raggedright
frontend
\end{minipage} & \begin{minipage}[b]{\linewidth}\raggedright
sha256
\end{minipage} \\
\midrule\noalign{}
\endhead
\bottomrule\noalign{}
\endlastfoot
committed & gunzip of \code{pnr/asap7/soc/macro/rv32i_cpu_top.nl.v.gz}
(post-P\&R, built at RTL \code{999e44a}) & Synlig (config
\code{USE_SYNLIG:true}; "built by Synlig" is inferred from its flop
count 4863 and no \code{fwd_b_*}, its run log is gone) &
\code{a416d23f038065433eafa81553e85b9d599892dc9e6da64ca632563b6d1da0b2} \\
S & Synlig synthesis, unmodified main RTL & Synlig &
\code{a2abef9c8f5d024bc512a8bb09ab5fdcb39ed86c9b14d01e03b72b22b5efe046} \\
W & Synlig + the two hazard-unit part-selects routed through a named
full-width wire & Synlig &
\code{505b939796da1977ec8ab5243008b4eea66ed4f0175b2471fe9e7fb6d20d3eee} \\
V & sv2v frontend & sv2v &
\code{c2d7f3f37d7d08bece3fd0f8517eef4b682a19af9ee0268353d19a92f6216e24} \\
\end{longtable}
}

}

Flop counts (PD agent): committed 4863, S 4863, W 4868, V 4868
(\code{fwd_b_*} bank present only in W and V).

RTL arms (this test): \textbf{main (CPU RTL identical to
\code{5eb5f63})} for S/W/V; \textbf{\code{999e44a}} (the commit the
committed macro was built at, \ghref{96} run 23) for the committed macro.
The two RTL arms were compared with each other: IDENTICAL on all 6
programs, cycle-exact on every output port, so the committed-vs-main RTL
cleanups (declare-before-use, \code{default_nettype}, empty blocks;
hazard-unit lines and regfile unchanged) do not change observable
behaviour and the choice of RTL reference does not matter for the
verdicts below (committed vs \code{999e44a} and vs main give the same
result, MEASURED).

Committed-netlist physical cells: it instantiates
\code{DECAPx1/2/4/6/10}, \code{FILLER}, \code{FILLERxp5}, \code{TAPCELL}
(40,359 lines). None has a port connection (the netlist has no power
pins), so they were \textbf{removed} with \code{asap7/strip_physical.sh}
(function-neutral by construction; stub modules would be equivalent).
Every other cell type is covered (checked per netlist) by the
Liberty-derived combinational models plus
\code{asap7_seq_cell_models.v}; the 10 SRAM macros use
\code{asap7_sram_1rw_256x32_model.v}. The RTL arm uses the same SRAM and
ICG models (\code{SRAM_ASAP7}, \code{USE_ICG_CELL}).

\end{record}

\subsection{Method}\label{cr:defect:asap7-cpu-synlig-ma7:method}

\begin{record}

\begin{itemize}
\tightlist
\item
  \textbf{Simulator 1: yosys 0.62 \code{sim}}, the original u99/ma7
  harness unmodified (\code{run_cpu_macro_check.sh},
  \code{run_cpu_macro_check_branch.sh}, \code{tb_cpu_macro_check.v};
  Liberty
  \texttt{read\_\allowbreak{}liberty\ -\allowbreak{}ignore\_\allowbreak{}miss\_\allowbreak{}func}
  models). Programs: \code{ma7_straight} (u99) and \code{ma7_branch}
  (ma7 step 1). \code{DBG_GPR[4]} is the testbench\textquotesingle s
  captured APB read, the commit stream comes from the VCD. Wrapper:
  \code{asap7/run_yosys_gate.sh}.
\item
  \textbf{Simulator 2: Verilator 5.048}, cycle-exact. One testbench
  (\code{asap7/tb_asap7_cpu_check.sv}, derived from the dud4
  \code{tb_sky130_cpu_check.sv}) for both arms; gate arm = netlist +
  Liberty-derived cell models (\code{asap7/gen_cell_lib.sh}), RTL arm =
  RTL tree with the same SRAM/ICG models. All 6 programs of the dud4 set
  (\code{trivial}, \code{ma7_straight}, \code{ma7_branch},
  \code{sweepA/B/C}); compared: commit stream, AXI writes, 32 debug GPR
  reads and a per-cycle trace of every output port
  (\texttt{sky130/\allowbreak{}compare\_\allowbreak{}arms.\allowbreak{}py\ diff}),
  plus the 31-register port matrix (\code{sky130/port_matrix.py}).
\item
  Every yosys/\allowbreak{}Verilator job ran under
  \texttt{systemd-\allowbreak{}run\ -\allowbreak{}-\allowbreak{}user\ -\allowbreak{}-\allowbreak{}scope\ -\allowbreak{}p\ MemoryMax=\allowbreak{}5G\ -\allowbreak{}p\ MemorySwapMax=\allowbreak{}0}.
\end{itemize}

\end{record}

\subsubsection{Harness fix found on the way (MEASURED)}\label{cr:defect:asap7-cpu-synlig-ma7:harness-fix-found-on-the-way-measured}

\begin{record}

The dud4 TB\textquotesingle s APB sequencer never completes on
main-vintage RTL (0 debug reads, TIMEOUT even in the RTL arm). Cause:
the TB drove APB inputs with non-blocking assignments at the clock edge;
under Verilator \code{--timing} a TB process resumed by a posedge runs
before the DUT\textquotesingle s clocked logic of the same edge, so the
DUT already saw the new inputs at that edge and the registered
\code{pready} came one cycle earlier than the sequencer expected.
\code{tb_asap7_cpu_check.sv} drives APB inputs 1 ns after the edge with
blocking assignments; then all 32 debug reads work on the RTL arms
(\code{trivial}: \texttt{x1\ =\allowbreak{}\ 5}, rest 0) and the port
matrix and debug-GPR comparison are available for main-vintage RTL (they
were not in the dud4 main-RTL pair). The other TB inputs (AXI
responders) still use NBAs at the edge; they are identical in both arms
and were unaffected (cycle-exact agreement of the clean arms).

\end{record}

\subsection{\texorpdfstring{Result 1: the two ma7 programs, yosys
\code{sim} (original harness), four netlists
(MEASURED)}{Result 1: the two ma7 programs, yosys  (original harness), four netlists (MEASURED)}}\label{cr:defect:asap7-cpu-synlig-ma7:result-1-the-two-ma7-programs-yosys-sim-original-harness-four-netlists-measured}

\begin{record}

{\small\setlength{\tabcolsep}{4pt}

{
\begin{longtable}[]{@{}
  >{\raggedright\arraybackslash}p{\dimexpr(\linewidth-6\tabcolsep)*143/1000\relax}
  >{\raggedright\arraybackslash}p{\dimexpr(\linewidth-6\tabcolsep)*252/1000\relax}
  >{\raggedright\arraybackslash}p{\dimexpr(\linewidth-6\tabcolsep)*344/1000\relax}
  >{\raggedright\arraybackslash}p{\dimexpr(\linewidth-6\tabcolsep)*241/1000\relax}@{}}
\toprule\noalign{}
\begin{minipage}[b]{\linewidth}\raggedright
netlist
\end{minipage} & \begin{minipage}[b]{\linewidth}\raggedright
\code{ma7_straight}: \code{DBG_GPR[4]} (RTL \code{0x0c})
\end{minipage} & \begin{minipage}[b]{\linewidth}\raggedright
\code{ma7_branch}: \texttt{BEQ\ x2,\allowbreak{}x3} (x2 = x3 = 9, must
be TAKEN)
\end{minipage} & \begin{minipage}[b]{\linewidth}\raggedright
\code{ma7_branch}: \code{DBG_GPR[4]} (RTL \code{0x0e})
\end{minipage} \\
\midrule\noalign{}
\endhead
\bottomrule\noalign{}
\endlastfoot
committed (post-P\&R) & \textbf{\code{0x30}} FAIL & \textbf{not taken},
retires WRONG marker \code{0x14}, parks at \code{0x18} FAIL &
\textbf{\code{0x2e}} FAIL \\
S (Synlig, main RTL) & \textbf{\code{0x30}} FAIL & \textbf{not taken},
parks at \code{0x18} FAIL & \textbf{\code{0x2e}} FAIL \\
W (Synlig + workaround) & \code{0x0c} PASS & taken, retires
\code{0x20},\code{0x24}, parks at PASS \code{0x28} & \code{0x0e} PASS \\
V (sv2v) & \code{0x0c} PASS & taken, parks at \code{0x28} & \code{0x0e}
PASS \\
\end{longtable}
}

}

Commit streams: committed and S identical (73 commit edges); W and V
identical (70).

\end{record}

\subsection{Result 2: six programs, Verilator, cycle-exact vs RTL
(MEASURED)}\label{cr:defect:asap7-cpu-synlig-ma7:result-2-six-programs-verilator-cycle-exact-vs-rtl-measured}

\begin{record}

"IDENTICAL" = commit stream, AXI writes, 32 debug GPR reads and the
per-cycle trace of every output port all equal.

\begin{landscape}
\small\setlength{\tabcolsep}{4pt}

{
\begin{longtable}[]{@{}
  >{\raggedright\arraybackslash}p{\dimexpr(\linewidth-14\tabcolsep)*206/1000\relax}
  >{\raggedright\arraybackslash}p{\dimexpr(\linewidth-14\tabcolsep)*77/1000\relax}
  >{\raggedright\arraybackslash}p{\dimexpr(\linewidth-14\tabcolsep)*120/1000\relax}
  >{\raggedright\arraybackslash}p{\dimexpr(\linewidth-14\tabcolsep)*103/1000\relax}
  >{\raggedright\arraybackslash}p{\dimexpr(\linewidth-14\tabcolsep)*86/1000\relax}
  >{\raggedright\arraybackslash}p{\dimexpr(\linewidth-14\tabcolsep)*86/1000\relax}
  >{\raggedright\arraybackslash}p{\dimexpr(\linewidth-14\tabcolsep)*86/1000\relax}
  >{\raggedright\arraybackslash}p{\dimexpr(\linewidth-14\tabcolsep)*215/1000\relax}@{}}
\toprule\noalign{}
\begin{minipage}[b]{\linewidth}\raggedright
gate netlist (RTL ref)
\end{minipage} & \begin{minipage}[b]{\linewidth}\raggedright
trivial
\end{minipage} & \begin{minipage}[b]{\linewidth}\raggedright
ma7\_\allowbreak{}straight
\end{minipage} & \begin{minipage}[b]{\linewidth}\raggedright
ma7\_\allowbreak{}branch
\end{minipage} & \begin{minipage}[b]{\linewidth}\raggedright
sweepA
\end{minipage} & \begin{minipage}[b]{\linewidth}\raggedright
sweepB
\end{minipage} & \begin{minipage}[b]{\linewidth}\raggedright
sweepC
\end{minipage} & \begin{minipage}[b]{\linewidth}\raggedright
port matrix (124 cells)
\end{minipage} \\
\midrule\noalign{}
\endhead
\bottomrule\noalign{}
\endlastfoot
committed (\code{999e44a}) & IDENT & DIVERGES & DIVERGES & DIVERGES &
DIVERGES & DIVERGES & 96 FAIL \\
S (main) & IDENT & DIVERGES & DIVERGES & DIVERGES & DIVERGES & DIVERGES
& 96 FAIL \\
W (main) & IDENT & IDENT & IDENT & IDENT & IDENT & IDENT & \textbf{124 /
124 PASS} \\
V (main) & IDENT & IDENT & IDENT & IDENT & IDENT & IDENT & \textbf{124 /
124 PASS} \\
\end{longtable}
}

\end{landscape}

(committed vs main RTL gives the same divergences as vs \code{999e44a};
\code{rtl_999} vs \code{rtl_main} is IDENTICAL on all 6.) Failing-arm
signatures, identical on committed and S: \code{ma7_straight} x2/x3/x4 =
\code{0x0c/0x18/0x30} (RTL \code{0x07/0x0c/0x0c}); \code{ma7_branch}
first commit divergence \#5 (BEQ not taken, pc \code{0x14} instead of
\code{0x20}); sweeps A/B/C all diverge at commit \#69 (first MMIO store
at pc \code{0x114}, restart at pc 0, 1392 commits and 0 AXI writes
against 916 commits and 116 writes). This is the Sky130 dud4 signature,
number for number.

Caveat on the failing matrix: on S and committed the sweeps derail at
the first store, so the 96 FAIL cells are "the program never reached the
store", not 96 independently faulty register read cells. The clean
verdict on the regfile read path comes from W and V
(124/\allowbreak{}124). The debug-read column passes for 28 of 31
registers on S (x2, x30 and one more fail), i.e. most of the registers
hold the expected final value; the failing ones are those written by
wrong results.

\end{record}

\subsection{Controls (MEASURED unless marked)}\label{cr:defect:asap7-cpu-synlig-ma7:controls-measured-unless-marked}

\begin{record}

\begin{itemize}
\tightlist
\item
  \textbf{(a) positive:} \code{trivial} is IDENTICAL on all four
  netlists and both simulators agree on
  \code{ma7_straight}/\code{ma7_branch} verdicts on all four (yosys
  \code{sim} and Verilator give the same \code{x4} and the same branch
  outcome).
\item
  \textbf{(b) the gate arm really executes:} 74 / 96 / 196 commits on
  trivial / straight / branch, PCs advance through the expected
  sequences, the clean arms retire 916 commits and perform 116 AXI
  stores in each sweep; no stuck-in-reset and no X (all debug reads
  defined).
\item
  \textbf{(c) negative (fault injection into arm V, clean baseline
  IDENTICAL):} stuck-at-0 on the sinks of a single net,
  \code{asap7/inject_stuck.py}. \code{u_core.mem_wb_reg[160]} (INFERRED
  from the struct packing: instruction bit 21): detected in \textbf{5 of
  6} programs (all but \code{trivial}); \code{u_core.mem_wb_reg[150]}
  (INFERRED: instruction bit 11 = \code{rd[4]}): detected in \textbf{4
  of 6} (\code{ma7_branch}, \code{sweepA/B/C}; not \code{trivial}, nor
  \code{ma7_straight}, which only writes x1..x4 so \code{rd[4]} is
  always 0 and the stuck-at is invisible). \textbf{Limitation:} only
  nets that keep a name after ASAP7 mapping can be targeted
  (\code{mem_wb_reg} pc/\allowbreak{}instruction bits); the regfile
  storage, read mux and hazard-compare nets are anonymous, so no fault
  could be injected \emph{into the regfile read mux} as the Sky130
  control did. The injected faults are on the commit-observation path.
  The sensitivity of the differential to a datapath fault is instead
  shown by arms S and committed, which are real datapath faults and are
  detected in 5 of 6 programs.
\item
  \textbf{(d) X / initialisation:} random power-up state
  (\code{+verilator+rand+reset+2}, seeds 7 and 23) on the RTL arms and
  all four gate arms: every arm is IDENTICAL to its own zero-initialised
  run on all 6 programs, and gate-vs-RTL verdicts are unchanged (W, V:
  6/6 identical; S, committed: 1/6, only \code{trivial}). The yosys runs
  use \code{-zinit}. INFERRED: that the randomisation actually changed
  the initial state was not independently checked (same mechanism as the
  dud4 control).
\end{itemize}

\end{record}

\subsection{First divergence and mechanism}\label{cr:defect:asap7-cpu-synlig-ma7:first-divergence-and-mechanism}

\begin{record}

MEASURED (per-cycle boundary trace, \code{asap7/first_divergence.py}),
arm S vs RTL main:

\begin{itemize}
\tightlist
\item
  \code{ma7_straight}: the commit stream and every output port are
  identical for the whole run; the first differing port value is the
  debug read of x2 (cycle 439, \code{prdata} \code{0x7} vs \code{0xc}).
  Same on committed vs \code{999e44a}. I.e. the instruction stream is
  fetched and retired correctly and only the \emph{data} computed is
  wrong.
\item
  \code{ma7_branch}: first differing port cycle is 39,
  \code{debug_pc_src} / \code{debug_take_branch_jump} asserted in RTL
  (BEQ taken) and not in S; the commit stream diverges at commit \#5.
\end{itemize}

Arithmetic mechanism (the failing values are reproduced exactly;
INFERRED internals, MEASURED values). Model: with the forwarding selects
independent of the consumer\textquotesingle s \code{rs1}/\code{rs2},
each ALU operand is the \emph{previous instruction\textquotesingle s
result}.

{
\begin{longtable}[]{@{}
  >{\raggedright\arraybackslash}p{\dimexpr(\linewidth-6\tabcolsep)*134/1000\relax}
  >{\raggedright\arraybackslash}p{\dimexpr(\linewidth-6\tabcolsep)*255/1000\relax}
  >{\raggedright\arraybackslash}p{\dimexpr(\linewidth-6\tabcolsep)*215/1000\relax}
  >{\raggedright\arraybackslash}p{\dimexpr(\linewidth-6\tabcolsep)*376/1000\relax}@{}}
\toprule\noalign{}
\begin{minipage}[b]{\linewidth}\raggedright
program
\end{minipage} & \begin{minipage}[b]{\linewidth}\raggedright
instr
\end{minipage} & \begin{minipage}[b]{\linewidth}\raggedright
model
\end{minipage} & \begin{minipage}[b]{\linewidth}\raggedright
S / committed (debug read)
\end{minipage} \\
\midrule\noalign{}
\endhead
\bottomrule\noalign{}
\endlastfoot
straight & \texttt{ADDI\ x1,\allowbreak{}x0,\allowbreak{}5} & 0+5 = 5 &
x1 = 5 (RTL 5) \\
straight & \texttt{ADDI\ x2,\allowbreak{}x0,\allowbreak{}7} & prev(5)+7
= \textbf{12} & x2 = \code{0x0c} (RTL 7) \\
straight & \texttt{ADD\ x3,\allowbreak{}x1,\allowbreak{}x2} & 12+12 =
\textbf{24} & x3 = \code{0x18} (RTL 12) \\
straight & \texttt{ADD\ x4,\allowbreak{}x3,\allowbreak{}x0} & 24+24 =
\textbf{48} & x4 = \code{0x30} (RTL 12) \\
branch & \texttt{ADDI\ x2,\allowbreak{}x0,\allowbreak{}9} & prev(5)+9 =
\textbf{14} & x2 = \code{0x0e} (RTL 9) \\
branch & \texttt{ADDI\ x3,\allowbreak{}x0,\allowbreak{}9} & 14+9 =
\textbf{23} & x3 = \code{0x17} (RTL 9) \\
branch & \texttt{ADD\ x4,\allowbreak{}x1,\allowbreak{}x2} & 23+23 =
\textbf{46} & x4 = \code{0x2e} (RTL 14) \\
\end{longtable}
}

The BEQ then compares stored x2 = 14 with stored x3 = 23 and is not
taken: the ma7 "corruption reaches the execute datapath" observation is
real, but it is wrong \emph{data produced by forwarding}, not a bad
regfile read (INFERRED that the BEQ reads those stored values).

Causal isolation (MEASURED): arm W differs from arm S in exactly the two
hazard-unit port-connection lines (\code{rv32i_core.sv:436-437}), same
Synlig frontend, same flow, same config; W is IDENTICAL to RTL and S is
not. Arm V (a different frontend, which elaborates those lines
correctly: no \texttt{5'x}) is also identical. Synthesis evidence (PD
agent, MEASURED): Synlig logs
\texttt{Range\ select\ [639:\allowbreak{}608]\ /\allowbreak{}\ [799:\allowbreak{}768]\ out\ of\ bounds\ on\ signal\ \textbackslash{}if\_\allowbreak{}id\_\allowbreak{}reg\ .\allowbreak{}.\allowbreak{}.\allowbreak{}\ undef},
the raw RTLIL has
\texttt{connect\ \textbackslash{}if\_\allowbreak{}id\_\allowbreak{}rs1\_\allowbreak{}addr\ 5'x},
the S/\allowbreak{}committed netlists have 4863 flops and no
\code{fwd_b_*} bank while W/V have 4868.

Not done: direct observation of the forwarding-select flops. In the
ASAP7 netlists the flops are anonymous (the hierarchy and
\code{fwd_*}/\code{id_ex_reg}/regfile storage names are lost in
ABC/\allowbreak{}DFFHQN mapping), so the Sky130-style probes of the
select flops and regfile storage are not available. The mechanism rests
on the arithmetic reproduction above and the W/S isolation, not on a
waveform of the select flops.

\end{record}

\subsection{Consequences}\label{cr:defect:asap7-cpu-synlig-ma7:consequences}

\begin{record}

\begin{itemize}
\tightlist
\item
  The ASAP7 CPU macro views (\code{pnr/asap7/soc/macro/rv32i_cpu_top.*})
  are functionally wrong for the reason Sky130\textquotesingle s are
  (MEASURED on the netlist; the LEF/\allowbreak{}Liberty views are
  derived from the same netlist). The config switch to
  \code{USE_SYNLIG:false} (sv2v) from \code{ma7} fixes it, as does an
  RTL change that routes the two part-selects through a full-width wire
  (arm W).
\item
  \textbf{Option for the blocked \code{lxv}/\code{ma7} sv2v re-harden
  (not a decision):} if the RTL fix (\prref{259}) is merged, a Synlig build
  is a second route to a correct ASAP7 CPU macro (arm W is the evidence,
  at synthesis level; place-and-route of a W-style netlist was not run).
  The Synlig route would also avoid the sv2v-specific routing regression
  recorded in \code{lxv} only if that regression is frontend-dependent,
  which this test does not show (\code{e69}\textquotesingle s 2,045
  violations came from the Synlig-era macro; the 108,716 from sv2v
  builds).
\item
  The same \code{struct.member[a:b]} in a port connection pattern may
  exist elsewhere (the GPU\textquotesingle s
  \code{warp_scheduler.sv:152} was flagged in \code{ma7}); only the
  \code{rv32i_core} cone was checked for the Synlig
  \texttt{out\ of\ bounds} warning (PD agent), the GPU was not
  simulated.
\item
  No step gates on Synlig\textquotesingle s
  \texttt{Range\ select\ .\allowbreak{}.\allowbreak{}.\allowbreak{}\ undef}
  warning (see \cref{cr:defect:sky130-cpu-synlig-dud4}).
\end{itemize}

\end{record}

\subsection{Not done}\label{cr:defect:asap7-cpu-synlig-ma7:not-done}

\begin{record}

Formal equivalence of S/W/V vs RTL; waveform-level view of the
forwarding-select flops (anonymous after mapping); a fault injected into
the regfile read mux itself; P\&R of an arm-W netlist; the GPU macro;
programs beyond the 6 (no loads, no CSR, no interrupts, no
cache-miss-heavy streams beyond what the sweeps do); yosys \code{sim}
runs of the sweeps (the sweeps ran in Verilator only).

\end{record}

\subsection{Reproduce}\label{cr:defect:asap7-cpu-synlig-ma7:reproduce}

\begin{record}

Scratch: \code{/nobackup/claude_sim_build/dud4/asap7_sim/}.
\texttt{asap7/\allowbreak{}gen\_\allowbreak{}cell\_\allowbreak{}lib.\allowbreak{}sh\ <out.\allowbreak{}v>}
(cell models), \code{asap7/strip_physical.sh} (committed netlist),
\texttt{asap7/\allowbreak{}build\_\allowbreak{}arm.\allowbreak{}sh\ rtl|gate}
(Verilator; \code{RTL_ROOT=} selects the RTL tree, \code{CELLLIB=} the
cell models), \code{sky130/gen_all_progs.sh}, \code{sky130/run_arm.sh},
\texttt{sky130/\allowbreak{}compare\_\allowbreak{}arms.\allowbreak{}py\ diff},
\code{sky130/port_matrix.py}, \code{asap7/run_yosys_gate.sh},
\code{asap7/first_divergence.py}, \code{asap7/inject_stuck.py}. Netlists
are not committed.

\end{record}

\subsection{\texorpdfstring{Update 2026-10-10 (bead \code{ainf}): is
"Synlig + PR \#259" safe for the remaining struct-member
selects?}{Update 2026-10-10 (bead ): is "Synlig + PR \#259" safe for the remaining struct-member selects?}}\label{cr:defect:asap7-cpu-synlig-ma7:update-2026-10-10-bead-ainf-is-synlig-pr-259-safe-for-the-remaining-struct-member-selects}

\begin{record}

Full evidence: \cref{cr:defect:synlig-struct-select-ainf}. Verdict, stated
no stronger than measured:

\begin{itemize}
\tightlist
\item
  \textbf{CPU (\src{rtl/cpu}): no second Synlig miscompile of this class
  was found.} Synlig has three struct-select defects: D1 (\code{s.m[..]}
  on a struct type with exactly two packed-range members), D2 (LHS
  part-select of a member) and D3 (element select of a multi-dimensional
  packed-array member). In \src{rtl/cpu} there are 15 remaining member
  selects (\code{ex:109}, \code{mem:103,104,111,112}, \code{wb:75},
  \code{ex1c:122-150,197}), all on struct types with 7..19 ranged
  members; the only CPU R == 2 type (\code{if_id_reg_t}) has no select
  since \prref{259}; there is no LHS member select and no multi-dimensional
  array member in any CPU struct. MEASURED: Synlig-vs-sv2v miter + SAT
  (flop-exposed, with a NOT-EQUIV negative control) proves
  \code{rv32i_pipeline_ex}, \code{_mem}, \code{_wb}, \code{_ex1c} EQUIV,
  in addition to the earlier hazard, forwarding, ALU, branch-compare,
  decode and immgen proofs. So the byte/\allowbreak{}halfword
  load-extraction and store-lane paths are covered by a proof, not by a
  simulation. NOT covered: whole-core equivalence (\code{rv32i_core},
  regfile, CSR file, caches were not miter-checked), and defects of
  other classes (the \code{ma7} regfile \code{OPT_MUXTREE} class is a
  separate question).
\item
  \textbf{GPU (\src{rtl/gpu}): NOT safe under Synlig.}
  \code{gpu_compute_unit} reads \code{id_ex_q.rs1_data[l]} /
  \code{rs2_data[l]}
  (\texttt{logic\ [N\_\allowbreak{}LANES-\allowbreak{}1:\allowbreak{}0][31:\allowbreak{}0]}
  members): Synlig elaborates the 32-bit element as ONE bit (D3). The
  VLD/\allowbreak{}VST/\allowbreak{}VLDS/\allowbreak{}VSTS addresses and
  store data are wrong in a Synlig netlist; the committed Synlig-built
  \code{gpu_top} macro views are affected.
\item
  \textbf{Fragility.} Whether a CPU struct is exposed to D1 depends on
  how many members carry a packed range: \code{if_id_reg_t} is R == 2
  today (safe only because it has no selects); the next-closest struct
  types are R \textgreater= 7 in the CPU and \code{dma_engine.desc_t} R
  == 3 (one removal from R == 2). Adding a member select to
  \code{if_id_reg_t}, \code{warp_state_t} or \code{div_entry_t}, or
  removing members from a 3-ranged-member struct, would silently re-arm
  D1. Any new GPU/CPU use of a multi-dimensional packed member would
  re-arm D3. The conclusion above is a property of the current source,
  not of Synlig. Recommendation:
  \src{tools/verif/check_struct_member_select.py} (baseline of
  today\textquotesingle s hits) in CI, rewrite the remaining sites
  through named wires, and for the GPU use sv2v or the named-wire
  rewrite before any Synlig GPU build.
\end{itemize}

\end{record}

% Converted from docs/SKY130_CPU_SYNLIG_DUD4.md by tools/docs/md2tex.py; edit this file, not the original.
\section{\texorpdfstring{Sky130 CPU macro netlist vs RTL differential
(bead
\code{dud4})}{Sky130 CPU macro netlist vs RTL differential (bead )}}\label{cr:defect:sky130-cpu-synlig-dud4}

\begin{record}

Date 2026-10-10. Question: is the Sky130 \code{rv32i_cpu_top} hard macro
netlist (built with
\texttt{USE\_\allowbreak{}SYNLIG:\allowbreak{}\ true}) functionally
corrupt, the way \code{u99}/\code{ma7} showed for ASAP7? Test only: no
PD run, no config change, no re-harden. Statements are tagged MEASURED,
INFERRED or HYPOTHESIS.

\end{record}

\subsection{Verdict}\label{cr:defect:sky130-cpu-synlig-dud4:verdict}

\begin{record}

\textbf{CORRUPT, and the cause is localised to the Synlig frontend.} The
committed Sky130 CPU macro netlist computes wrong results. Synlig
elaborates two port connections of the hazard unit as undefined; a
synthesis of the same pinned RTL through sv2v is functionally identical
to RTL on every program tried, and a Synlig synthesis of the same RTL
reproduces the committed netlist\textquotesingle s failure signature.
The mechanism is \textbf{hazard/\allowbreak{}forwarding compare against
undefined \code{rs1}/\code{rs2}}, not the regfile read mux that
\code{ma7}/\code{u99} blame.

\end{record}

\subsection{Root cause (MEASURED unless marked)}\label{cr:defect:sky130-cpu-synlig-dud4:root-cause-measured-unless-marked}

\begin{record}

\code{rv32i_core.sv} connects the hazard unit as
\code{.if_id_rs1_addr(if_id_reg.instruction[19:15])},
\code{.if_id_rs2_addr(if_id_reg.instruction[24:20])} (lines 432/433 at
\code{5c49ddf}; 436/437 on main). Synlig/\allowbreak{}Surelog logs
during \code{Yosys.Synthesis} and JsonHeader:

\begin{lstlisting}[language={}]
rv32i_core.sv:432: Warning: Range select [639:608] out of bounds on signal `\if_id_reg': Setting all 32 result bits to undef.
rv32i_core.sv:433: Warning: Range select [799:768] out of bounds on signal `\if_id_reg': Setting all 32 result bits to undef.
\end{lstlisting}

The raw RTLIL straight from Synlig has
\texttt{connect\ \textbackslash{}if\_\allowbreak{}id\_\allowbreak{}rs1\_\allowbreak{}addr\ 5'x}
and
\texttt{connect\ \textbackslash{}if\_\allowbreak{}id\_\allowbreak{}rs2\_\allowbreak{}addr\ 5'x}
(\code{localise/raw_synlig.il}, lines 22700-22701). A part-select of a
packed-struct field inside a port connection is what triggers it. sv2v
elaborates the same lines correctly. With the compare address undefined,
every
\texttt{rd\_\allowbreak{}addr\ =\allowbreak{}=\allowbreak{}\ if\_\allowbreak{}id\_\allowbreak{}rs?\_\allowbreak{}addr}
in \code{rv32i_hazard_unit} (pre-decoded forwarding selects, load-use
detect) is a don\textquotesingle t-care, so Yosys merges the B-side
registered selects into the A-side ones.

Flop census (post-synthesis, pinned RTL \code{5c49ddf}): Synlig 4,829
\code{dfxtp} with \textbf{0} \code{fwd_b_*} flops; sv2v 4,834 with 5
(\code{fwd_b_sel_r} x3, \code{fwd_b_ex1c_r}, \code{fwd_b_ex1b2_r}). The
committed post-PnR netlist has \textbf{4,829} \code{dfxtp} and no
\code{fwd_b_*} nets, i.e. exactly the Synlig count.
\texttt{Removed\ N\ multiplexer\ ports} counts are identical in both
arms (\code{OPT_MUXTREE} is not the difference). Source of these
numbers: \code{/nobackup/claude_sim_build/dud4/synth/NOTES.md} (PD-agent
synthesis-only experiment; commands, hashes and tool versions there:
LibreLane 2.4.13, Yosys 0.46, Surelog 1.82, sv2v 0.0.13.1).

\end{record}

\subsection{Simulation result: three netlists, same testbench, same
programs (MEASURED)}\label{cr:defect:sky130-cpu-synlig-dud4:simulation-result-three-netlists-same-testbench-same-programs-measured}

\begin{record}

RTL arm: RTL at \code{5c49ddf} (+\code{SRAM_SKY130}). Gate arms:
synthesis-only netlists (\code{Yosys.Synthesis} output, same PDK,
config, library, strategy per pair; single variable = frontend), PDK
functional cell models, project SRAM model. Programs: \code{trivial},
\code{ma7_straight} (u99), \code{ma7_branch}, \code{sweepA/B/C}
(31-register sweeps). Comparison: commit stream, AXI writes, 32 debug
GPR reads and a per-cycle trace of every output port.

{\footnotesize\setlength{\tabcolsep}{3pt}

{
\begin{longtable}[]{@{}
  >{\raggedright\arraybackslash}p{\dimexpr(\linewidth-2\tabcolsep)*490/1000\relax}
  >{\raggedright\arraybackslash}p{\dimexpr(\linewidth-2\tabcolsep)*490/1000\relax}@{}}
\toprule\noalign{}
\begin{minipage}[b]{\linewidth}\raggedright
gate netlist (pinned RTL \code{5c49ddf})
\end{minipage} & \begin{minipage}[b]{\linewidth}\raggedright
result vs RTL
\end{minipage} \\
\midrule\noalign{}
\endhead
\bottomrule\noalign{}
\endlastfoot
committed post-PnR netlist (\code{995c487}) & x4 = \code{0x30}, BEQ
x2,x3 not taken, sweeps trap at the first MMIO store (earlier section of
this PR) \\
\code{pinned/synlig} synth netlist & \textbf{5 of 6 programs diverge}
with the same signature: \code{ma7_straight} x4 = \code{0x30} (RTL
\code{0x0c}); \code{ma7_branch} first commit divergence \#5 (pc 0x14
instead of 0x20, BEQ not taken); sweeps A/B/C diverge at commit \#69
(first MMIO store, pc 0x114 -\textgreater{} restart at pc 0).
\code{trivial} identical. Port matrix: 96 of 124 exercised cells FAIL \\
\code{pinned/sv2v} synth netlist & \textbf{IDENTICAL to RTL on all 6
programs}, cycle-exact on every output port, all 32 debug reads. Port
matrix: \textbf{124 of 124 cells PASS} (x1..x31 x \{rs1, rs2, both,
debug\}) \\
\code{wa/synlig} (main RTL + scratch workaround: the two selects routed
through a full-width wire) & see "Main-RTL pair" below \\
\end{longtable}
}

}

So: the fault exists at synthesis output (P\&R, CTS, resizer and ECO are
exonerated for this signature), the fault comes from the frontend, and
the sv2v netlist of the same RTL is functionally correct for what was
tested.

Ruled out as harness effects (MEASURED): initial state/X (random-reset
runs, two seeds, identical to zero-init per arm); Verilator scheduling
(original \code{ma7} yosys-\code{sim} harness also gives \code{0x30} and
the not-taken branch on the committed netlist, and \code{0xc} on the
pinned sv2v synth netlist); SRAM model and latency (same file, both
arms); RTL commit (pinned; the RTL arm matches the by-construction
expectation 116/116 stores and final GPRs); cell models (\code{trivial}
is cycle-exact identical on every netlist).

\end{record}

\subsubsection{Harness problem found and fixed on the way (MEASURED)}\label{cr:defect:sky130-cpu-synlig-dud4:harness-problem-found-and-fixed-on-the-way-measured}

\begin{record}

On the first attempt the synthesis-output netlists produced no commits
at all in Verilator, for \textbf{both} frontends. Cause: the PDK
\code{sky130_fd_sc_hd__conb_1} model uses \code{pullup}/\code{pulldown}
primitives; where the Yosys netlist drives output
ports/\allowbreak{}nets straight from tie cells, Verilator resolves them
as pulled tri-states, warns
\texttt{Circular\ combinational\ logic\ .\allowbreak{}.\allowbreak{}.\allowbreak{}\ \_\allowbreak{}\_\allowbreak{}out\_\allowbreak{}\_\allowbreak{}strong\_\allowbreak{}\_\allowbreak{}out}
and returns 0 for \code{HI}, so the CPU never starts. Fix:
\code{CONB_FIX=1} in \code{build_arm.sh} swaps the cell type for a plain
\texttt{HI=\allowbreak{}1,\allowbreak{}\ LO=\allowbreak{}0} module
(\code{conb_model.v}). The yosys-\code{sim} Liberty run is unaffected
and independently gave the same answers. The committed post-PnR netlist
was not affected (its ties go through \code{assign} nets).

\end{record}

\subsection{Corrected mechanism}\label{cr:defect:sky130-cpu-synlig-dud4:corrected-mechanism}

\begin{record}

\begin{itemize}
\tightlist
\item
  MEASURED (committed netlist, \code{ma7_straight} probes): regfile read
  data into ID/EX equals RTL in every cycle except the final
  \texttt{SW\ x4}; stored values are wrong from \code{x2} on
  (\code{0x0c} instead of \code{0x07} at cycle 26); the registered
  forwarding-select flops assert in a rolling pattern (cycles 25-31)
  where RTL asserts once.
\item
  INFERRED from the arithmetic and the root cause: with
  \code{if_id_rs1/rs2} undefined the forwarding selects depend on the
  producer alone, so each ALU operand is the previous
  instruction\textquotesingle s result (\code{7+5}, \code{12+12},
  \code{24+24}), including for \code{x0}. That also explains why the
  sweeps trap at the first store: the address operand is forwarded
  garbage.
\item
  The debug read mux is faithful (storage == debug read for all 32
  registers). The regfile read mux is \textbf{not} the fault; it passes
  the full 31-register x port matrix on the sv2v netlist, and the Synlig
  netlist\textquotesingle s regfile read results are not what fails.
\end{itemize}

\end{record}

\subsection{Why nothing caught it}\label{cr:defect:sky130-cpu-synlig-dud4:why-nothing-caught-it}

\begin{record}

\begin{itemize}
\tightlist
\item
  Port-connection fault \emph{between} modules: per-module
  Synlig-vs-sv2v SAT equivalence (hazard, forwarding, ALU, branch
  comparator, decode, immediate generator: all equivalent in isolation;
  \code{equiv_all_comb.sh}) cannot see it because the broken wires are
  the instance\textquotesingle s port connection inside
  \code{rv32i_core}, outside every isolated module.
\item
  It is a logged \textbf{Warning} that no step gates on.
  \code{Checker.YosysSynthChecks}, lint (Verilator on the sources),
  unmapped-cell, NetlistAssign and inferred-latch checks all pass
  (\code{synthesis__check_error__count} 0 in every arm).
\item
  Netgen LVS compares layout to the same wrong netlist; STA and P\&R are
  happy with a smaller, wrong design (Synlig netlist is 333-957 cells
  smaller).
\item
  The earlier \code{OPT_MUXTREE} hypothesis fits the symptom only by
  coincidence (identical removal counts in both arms).
\end{itemize}

\end{record}

\subsection{Consequences}\label{cr:defect:sky130-cpu-synlig-dud4:consequences}

\begin{record}

\begin{itemize}
\tightlist
\item
  The Sky130 SoC consumes the wrong CPU macro
  (\src{pnr/sky130/soc/config.json} MACROS). Every Sky130 SoC result
  that depends on CPU behaviour at gate level rests on it; RTL-level
  cocotb results are unaffected.
\item
  ASAP7 (INFERRED, untested): the ASAP7 CPU config is also
  \code{USE_SYNLIG:true} on RTL with the same two lines, and its
  \code{u99}/\code{ma7} signature (x4 = \code{0x30}, BEQ not taken) is
  identical to the one reproduced here by the Synlig netlist. So the
  ASAP7 failure is very likely this defect, and the
  \code{u99}/\code{ma7} "regfile read-mux OPT\_\allowbreak{}MUXTREE"
  diagnosis is in doubt. Not verified by building or simulating any
  ASAP7 netlist here.
\item
  Remedies (owner\textquotesingle s decision, both need a re-harden):
  switch the Sky130 CPU config to sv2v (as the SoC and, since
  \code{ma7}, ASAP7 already are) and/or rewrite the two part-selects
  through a full-width wire; gate on the Synlig
  \texttt{Range\ select\ .\allowbreak{}.\allowbreak{}.\allowbreak{}\ out\ of\ bounds}
  / \code{undef} warning.
\item
  Surelog shared-cache caveat (MEASURED, from the synthesis run): the
  Makefile runs LibreLane with cwd = the librelane checkout, so
  Surelog\textquotesingle s \code{slpp_all} cache is shared across every
  project and worktree; one run logged source paths from an older
  worktree whose \code{rv32i_core.sv}, \code{rv32i_hazard_unit.sv} and
  \code{rv32i_forwarding_unit.sv} differ. The resulting netlist was
  byte-identical to the clean-cwd run, so no stale content leaked in
  that instance, but it is a latent hazard and cannot be excluded for
  the committed netlist (its run directory is gone).
\end{itemize}

\end{record}

\subsection{Negative control, on a clean baseline (MEASURED)}\label{cr:defect:sky130-cpu-synlig-dud4:negative-control-on-a-clean-baseline-measured}

\begin{record}

Baseline = \code{pinned/sv2v} netlist (IDENTICAL to RTL on all 6
programs). Fault (scratch copy,
\texttt{inject\_\allowbreak{}fault.\allowbreak{}py\ tie1}): data input
A1 tied to 0 on the 63 first-level read-mux cells that select x1 among
x0..x3 in every read tree (\code{find_x1_read_cells.py}). Detected in
\textbf{4 of 6 programs} (\code{ma7_branch}, \code{sweepA},
\code{sweepB}, \code{sweepC}; first AXI write divergence: expected
\code{0x91b7584a}, got \code{0x00001000}). Not detected in
\code{trivial} and \code{ma7_straight} (neither reads x1 through a
settled read-mux path that matters there: \code{x1} is consumed via
forwarding or only through the debug path). So the differential detects
a regfile read-path fault when the program reads settled registers,
which the sweeps do. Earlier controls against the faulty committed
netlist were weak (an S0/S1 swap was invisible because x1 and x2 share
bit 0).

\end{record}

\subsection{\texorpdfstring{Main-RTL pair (\code{origin/main} =
\code{5eb5f63}, MEASURED)}{Main-RTL pair ( = , MEASURED)}}\label{cr:defect:sky130-cpu-synlig-dud4:main-rtl-pair-originmain-5eb5f63-measured}

\begin{record}

RTL arm \code{RTL_ROOT=.../synth/rtl_main}; gate arms \code{synlig/},
\code{sv2v/} and \code{wa/synlig} (main + the scratch workaround).
Caveat: \code{999e44a} registered the APB outputs, and
\code{tb_sky130_cpu_check.sv}\textquotesingle s APB sequencer hangs on
the main RTL itself (0 debug reads, TIMEOUT in the RTL arm too), so for
this pair \textbf{debug-port reads and the port matrix are unavailable};
the sequencer was not adapted. Everything else (commit stream functional
and cycle-exact, AXI write stream functional and cycle-exact, per-cycle
trace of every output port) was compared, and the only failing check
reported in every program is the 0-vs-32 debug-read count, equal on both
sides. The RTL arm\textquotesingle s 116/116 sweep stores match the
by-construction expectation.

{\footnotesize\setlength{\tabcolsep}{3pt}

{
\begin{longtable}[]{@{}
  >{\raggedright\arraybackslash}p{\dimexpr(\linewidth-2\tabcolsep)*490/1000\relax}
  >{\raggedright\arraybackslash}p{\dimexpr(\linewidth-2\tabcolsep)*490/1000\relax}@{}}
\toprule\noalign{}
\begin{minipage}[b]{\linewidth}\raggedright
gate netlist (main RTL)
\end{minipage} & \begin{minipage}[b]{\linewidth}\raggedright
result vs main RTL
\end{minipage} \\
\midrule\noalign{}
\endhead
\bottomrule\noalign{}
\endlastfoot
\code{sv2v/} & commit, AXI and per-cycle trace identical on all 6
programs \\
\code{wa/synlig} (workaround: two selects via a full-width wire) &
commit, AXI and per-cycle trace identical on all 6 programs, i.e. the
same Synlig frontend is fine once the two lines are routed through a
full-width wire. \textbf{This isolates the defect to those two lines
under Synlig.} \\
\code{synlig/} (unmodified) & fails: \code{ma7_branch} first commit
divergence \#5 (BEQ not taken), sweeps A/B/C divergence at commit \#69,
same as pinned \\
\end{longtable}
}

}

The defect is therefore still present on current \code{main}.

\end{record}

\subsection{Not done / limits}\label{cr:defect:sky130-cpu-synlig-dud4:not-done-limits}

\begin{record}

\begin{itemize}
\tightlist
\item
  Whole-CPU formal equivalence between netlists and RTL; ASAP7 netlist
  simulation; the sv2v netlist was shown correct only for the programs
  above (all 31 registers on both ID ports, plus both ma7 programs), not
  exhaustively.
\item
  The committed netlist\textquotesingle s own run log is gone, so "built
  by Synlig" for that exact file is inferred from its signature (4,829
  flops, no \code{fwd_b_*}) plus the config, not read from a log.
\item
  \code{rv32i_core.sv} line-number mapping is per RTL version (432/433
  pinned, 436/437 main).
\end{itemize}

\end{record}

\subsection{Reproduce}\label{cr:defect:sky130-cpu-synlig-dud4:reproduce}

\begin{record}

\code{gen_all_progs.sh},
\texttt{build\_\allowbreak{}arm.\allowbreak{}sh\ rtl|gate}
(\code{CONB_FIX=1} for Yosys synthesis-output netlists; \code{RTL_ROOT=}
selects the RTL tree), \code{run_arm.sh},
\texttt{compare\_\allowbreak{}arms.\allowbreak{}py\ diff},
\code{port_matrix.py}, \code{run_and_compare.sh},
\code{run_yosys_gate_check.sh}, \code{inject_fault.py},
\code{find_x1_read_cells.py}, \code{equiv_all_comb.sh}; probes
\code{build_probe_arms.sh}, \code{run_probe.sh}, \code{decode_trace.py}.
Scratch under \code{/nobackup/claude_sim_build/dud4}. Netlists are not
committed.

\end{record}

\subsection{Gate (bead gc0y)}\label{cr:defect:sky130-cpu-synlig-dud4:gate-bead-gc0y}

\begin{record}

The defect class is now gated in every \src{pnr/Makefile} LibreLane run
(\src{docs/SYNTH_UNDEF_GATE.md} (\cref{tm:flow:synth-undef-gate})): a log scanner plus a structural check
of the elaborated header JSON, via the
\src{pnr/plugins/librelane_plugin_cvt_synthgate} plugin. Negative
control: this note\textquotesingle s reverted \code{rv32i_core.sv}
hazard lines make
\texttt{make\ librelane-\allowbreak{}sky130-\allowbreak{}cpu} abort at
the JsonHeader step.

Testbench fix: \src{tools/verif/gls/sky130/tb_sky130_cpu_check.sv} now
drives APB inputs 1 ns after the clock edge with blocking assignments
(back-ported from the ASAP7 testbench). Before, its NBA-at-the-edge
sequencer raced the registered \code{pready} on main RTL and debug reads
timed out. Measured on main RTL (RTL arm, Verilator 5.048): all six
programs complete with 32 APB reads each; \code{ma7_straight} reads
\code{x3=0x0c}, \code{x4=0x0c}.

\end{record}

\subsection{Re-harden on sv2v (2026-10-11) --- remedy applied}\label{cr:defect:sky130-cpu-synlig-dud4:re-harden-on-sv2v-2026-10-11-remedy-applied}

\begin{record}

User-approved remedy: sv2v switch for \src{pnr/sky130/cpu} (this change)
plus the RTL part-select rewrite (\prref{259}). Correctness only; no DFT
pins.

\begin{itemize}
\tightlist
\item
  \textbf{S1}
  \texttt{make\ -\allowbreak{}C\ pnr\ sky130-\allowbreak{}cpu-\allowbreak{}sv2v}
  (26 SV files, defines
  \texttt{PDK\_\allowbreak{}sky130A\ SCL\_\allowbreak{}sky130\_\allowbreak{}fd\_\allowbreak{}sc\_\allowbreak{}hd\ \_\allowbreak{}\_\allowbreak{}librelane\_\allowbreak{}\_\allowbreak{}\ \_\allowbreak{}\_\allowbreak{}pnr\_\allowbreak{}\_\allowbreak{}\ SRAM\_\allowbreak{}SKY130}
  = the arm-V recipe above); \code{librelane-sky130-cpu} depends on it;
  config \code{USE_SYNLIG:false}, \code{VERILOG_FILES} = SRAM stub +
  generated file; nothing else changed. Real yosys: 22/22 modules
  defined, only the SRAM stub is a blackbox. \code{FP_PIN_ORDER_CFG} was
  already wired (\code{pin_order.cfg}), so no pin file was generated.
\item
  \textbf{S2} synthesis only (\code{dud4_synth},
  origin/\allowbreak{}main \code{4d8e543} with \prref{259} and \#268): gate
  PASS (\code{fail=0}), census 4 868 flops (\code{fwd_b_sel_r},
  \code{fwd_b_ex1c_r}, \code{fwd_b_ex1b2_r} present), 28 102 cells,
  netlist sha256 \code{2acc6f6a...921d}. Differential: 6/6 programs
  IDENTICAL, matrix 124/124 (RTL arm also 124/124, debug reads work on
  main RTL after the testbench fix).
\item
  \textbf{S3} full P\&R \code{RUN_2026-10-11_06-00-12}, 75 MHz
  unchanged: LVS match, KLayout 0, routing DRC 0, Magic 27 733 913
  (99.9869 \% inside SRAM footprints; 3 626 \code{nwell.4}), setup clean
  at six tt/ff corners (+3.44 ns nom\_tt), \textbf{ss reg-to-reg setup
  fails} (nom\_ss -3.198 ns/854, max\_ss -3.704 ns/996; icache tag-SRAM
  half-cycle paths; reported, not tuned), hold reg-to-reg clean /
  I/O-only violators (nom\_tt -0.379 ns), antenna 96/96, 77.8 mW, 40.83
  \% util, 4 868 flops. Old ss CPU-standalone numbers are not
  recoverable.
\item
  \textbf{S4} post-fill netlist (the one committed): census passes; 6/6
  IDENTICAL; 124/124; negative control (63 x1 read-mux cells tied)
  diverges on 4/6 programs, 90 matrix cells FAIL. \code{pnl.v} equals
  \code{nl.v} modulo the power ports. Re-validated on the committed
  \code{.nl.v.gz} with
  \src{tools/verif/gls/sky130/run_netlist_proof.sh}.
\item
  \textbf{S5} views replaced (sha256, 12 chars, old -\textgreater{}
  new): lef \texttt{018792de0046\ -\allowbreak{}>\ dfbfb5eba536};
  nl.v.gz \texttt{d905d526e8b8\ -\allowbreak{}>\ 9e45b339ffad}; gds
  \texttt{474056b81d76\ -\allowbreak{}>\ 7923de1f2d34}; spice
  \texttt{9aa1afb320e2\ -\allowbreak{}>\ 2cd659eac361}; nom\_tt lib
  \texttt{3b9410015273\ -\allowbreak{}>\ c777011d5b9f}, max\_ss
  \texttt{738952c5cd95\ -\allowbreak{}>\ 69c843f19d23} (all nine are
  pairwise distinct with correct per-corner temperature and voltage).
  403/403 LEF pins at identical position and layer; the Liberty has 0
  combinational in-\textgreater out arcs (old 462, the old macro
  predates registered APB outputs).
\item
  \textbf{S6} SoC \code{RUN_2026-10-11_08-39-41}: see
  \cref{cr:defect:sky130-slew-cap-e45j:8-new-baseline-on-the-sv2v-rebuilt-cpu-macro-2026-10-11-bead-dud4-stale-waiver-3-unwaived-ss-corner-nets} and
  \src{docs/SKY130_REAL_DRC_LVS_EVALUATION.md} (\cref{tm:pd:sky130-real-drc-lvs}). Not a single-variable
  comparison with the 2026-10-09 baseline.
\end{itemize}

Open (owner): the stale \code{axi_rdata_i[12]} max-cap waiver and 3
unwaived ss-corner nets; the \code{gc0y} structural \code{x-driven-net}
check false-positived on 509 sv2v function-inlining temporaries on the
SoC (that run used \code{CVT_SYNTH_GATE=warn}, shell-only) -
\textbf{fixed by bead \bead{j1b7}}: dead function temporaries are exempt
on a pre-proc RTLIL proof and the SoC gate runs blocking
(\src{docs/SYNTH_UNDEF_GATE.md} (\cref{tm:flow:synth-undef-gate})); ASAP7 macros are still the corrupt
views (shelved).

\end{record}

% Converted from docs/SKY130_SLEW_CAP_E45J.md by tools/docs/md2tex.py; edit this file, not the original.
\section{\texorpdfstring{Sky130 SoC max-slew / max-cap: root cause and
the RC-calibrated repair (bead
\code{e45j})}{Sky130 SoC max-slew / max-cap: root cause and the RC-calibrated repair (bead )}}\label{cr:defect:sky130-slew-cap-e45j}

\begin{record}

\textbf{Status 2026-10-09 (updated).} Mechanism found, fix measured on
four runs. Sections 5-6 record why adoption was first refused (one
max-cap net, cap 1). \textbf{Section 7 supersedes that: by user decision
2026-10-09 the RC-calibrated flow (\code{config_rccal.json} + the
project plugin) is now the DEFAULT Sky130 SoC flow and max-cap is gated
at all nine corners with ONE named single-net waiver.} Bead left
\textbf{open} (max-slew residual 102-423 is still skipped and tracked).
Everything here is measured from run artifacts unless marked
\emph{inference}.

\end{record}

\subsection{1. The earlier root cause was only half right}\label{cr:defect:sky130-slew-cap-e45j:1-the-earlier-root-cause-was-only-half-right}

\begin{record}

\code{e45j} attributed the violations to "repair runs on GRT-estimated
parasitics, there is no post-RCX repair stage". Two corrections, both
measured on the Gate B baseline \code{RUN_2026-10-05_06-50-11}:

\begin{enumerate}
\def\labelenumi{\arabic{enumi}.}
\tightlist
\item
  \textbf{The "0/0/0 at steps 33/40/41" figures were nom\_tt only.}
  Mid-PnR STA (\code{OpenROAD.STAMidPNR}) runs a single corner
  (\texttt{DEFAULT\_\allowbreak{}CORNER\ =\allowbreak{}\ nom\_\allowbreak{}tt});
  the other eight corners\textquotesingle{} values in each
  step\textquotesingle s \code{state_out.json} are stale copies from
  step 9. Step 51 is the first all-corner measurement, so the headline
  "0 -\textgreater{} 8357" compared one corner with nine.
\item
  \textbf{The repair steps are corner-aware; the estimator is the
  problem.} Loading the step-40 ODB, re-running global routing
  in-session and reporting per corner gave \textbf{max\_ss slew 307 (277
  on ANTENNA pins), cap 0, and 0 at every tt/ff corner} -\/- versus 7622
  / 140 after RCX. So at the GRT estimate the design was already almost
  clean at the worst corner; the post-RCX gap is the estimated RC being
  too optimistic.
\end{enumerate}

Why: the sky130 tech LEF gives ground-only per-layer C (met1 0.085 fF/um
... met5 0.088 fF/um, no coupling or fringe).
\texttt{estimate\_\allowbreak{}parasitics\ -\allowbreak{}spef\_\allowbreak{}file}
on the same ODB compared with the OpenRCX SPEF of the same nets (17 525
nets with C \textgreater{} 5 fF), total RCX / estimate:

{
\begin{longtable}[]{@{}lllll@{}}
\toprule\noalign{}
RC corner & sum C & sum R & per-net median C & per-net median R \\
\midrule\noalign{}
\endhead
\bottomrule\noalign{}
\endlastfoot
min & 1.69 & 0.62 & 1.58 & 0.58 \\
nom & 1.88 & 1.02 & 1.76 & 1.00 \\
max & 2.04 & 1.91 & 1.93 & 1.93 \\
\end{longtable}
}

Every repair step therefore sized buffers against roughly half the
capacitance that routing later delivers. (A post-RCX repair stage would
also work, but it is not needed to reach this result and it is the
option that needs the shared-LibreLane edit evaluated in
\code{LIBRELANE_PATCHING_EVALUATION.md} (\cref{tm:pd:librelane-patching}).)

Two more measured facts:

\begin{itemize}
\tightlist
\item
  \textbf{ANTENNA diode pins are 26 \% (max\_ss) to 76 \% (nom\_tt) of
  violating pins} in the baseline, but they sit on the same nets as
  ordinary sinks, so excluding them shrinks the count, not the set of
  bad nets. On max\_ss the 5 639 non-diode violating pins lie on
  \textbf{314 distinct nets}.
\item
  \textbf{\code{set_max_fanout} really is absent}, and the unconstrained
  distribution is large: baseline netlist has 5 181 nets with fanout
  \textgreater{} 10, 1 258 \textgreater{} 32, 51 \textgreater{} 100 (max
  217).
\end{itemize}

\end{record}

\subsection{2. The fix: project-local plugin, no edit to the shared
LibreLane install}\label{cr:defect:sky130-slew-cap-e45j:2-the-fix-project-local-plugin-no-edit-to-the-shared-librelane-install}

\begin{record}

\src{pnr/sky130/soc/plugin/librelane_plugin_cvt_sky130/} (LibreLane
imports any \code{librelane_plugin_*} on \code{PYTHONPATH}; the Makefile
adds the directory). It defines \code{CVT.RepairDesignPostGPL},
\code{CVT.ResizerTimingPostCTS}, \code{CVT.RepairDesignPostGRT},
\code{CVT.ResizerTimingPostGRT}, subclasses of the stock steps that run
the \textbf{stock script unchanged} through \code{tcl/cvt_wrap.tcl}. The
wrapper renames Tcl \code{source} so \code{tcl/rc_calibration.tcl} runs
right after \code{common/set_rc.tcl} and rescales the per-layer RC per
STA corner (\code{CVT_RC_CALIBRATION}, e.g.
\texttt{nom=\allowbreak{}1.\allowbreak{}88,\allowbreak{}1.\allowbreak{}02\ min=\allowbreak{}1.\allowbreak{}69,\allowbreak{}0.\allowbreak{}62\ max=\allowbreak{}2.\allowbreak{}04,\allowbreak{}1.\allowbreak{}91},
the table above). Verified by dumping the estimate
before/\allowbreak{}after: total estimated C scales 1.87 / 2.02 and R
1.02 / 1.91 for nom / max as set. Empty variable = byte-identical
pass-through.

Run it:
\texttt{make\ -\allowbreak{}C\ pnr\ librelane-\allowbreak{}sky130-\allowbreak{}soc-\allowbreak{}rccal-\allowbreak{}noklayout}
(selects \code{config_rccal.json}, which is \code{config.json} plus
\code{meta.substituting_steps} and \code{CVT_RC_CALIBRATION};
\code{config.json} is untouched).

Two plugin gotchas found: \code{TclStep} moves non-allowlisted env into
the \code{_TCL_ENV_IN} file, so a wrapper must \code{source} it before
reading its own variables; and \code{set_layer_rc} wants kohm/um and
pF/um (a first attempt with C off by 1e12 collapsed the estimate 67x,
caught by the dump).

\end{record}

\subsection{3. Result}\label{cr:defect:sky130-slew-cap-e45j:3-result}

\begin{record}

Baseline \code{RUN_2026-10-05_06-50-11} (the Gate B run; its config keys
and numbers match f7vs.15\textquotesingle s close reason) vs
\code{RUN_2026-10-06_22-43-34} (\code{config_rccal.json}). \textbf{RTL
caveat:} \code{dma_engine.sv} was pinned to its Gate B version
(\code{49ace76}) for the run, because \code{857750c} landed afterwards;
the run therefore changes exactly one thing. All 62 steps completed; the
only deferred error is the 8 984 Magic DRC of bead \bead{45a}.

{
\begin{longtable}[]{@{}
  >{\raggedright\arraybackslash}p{\dimexpr(\linewidth-4\tabcolsep)*405/1000\relax}
  >{\raggedright\arraybackslash}p{\dimexpr(\linewidth-4\tabcolsep)*342/1000\relax}
  >{\raggedright\arraybackslash}p{\dimexpr(\linewidth-4\tabcolsep)*234/1000\relax}@{}}
\toprule\noalign{}
\begin{minipage}[b]{\linewidth}\raggedright
\end{minipage} & \begin{minipage}[b]{\linewidth}\raggedright
baseline
\end{minipage} & \begin{minipage}[b]{\linewidth}\raggedright
RC-calibrated
\end{minipage} \\
\midrule\noalign{}
\endhead
\bottomrule\noalign{}
\endlastfoot
max\_ss slew / cap violating pins & 7622 / 140 & \textbf{295 / 0} \\
nom\_ss & 2648 / 78 & 143 / 0 \\
min\_ss & 1669 / 33 & 5 / 0 \\
max\_tt & 2331 / 3 & 209 / 0 \\
nom\_tt & 1016 / 1 & 64 / 0 \\
max\_ff / nom\_ff / min\_tt / min\_ff & 1661 / 944 / 876 / 753 (cap
2/0/0/0) & 122 / 53 / 0 / 0 (cap 0) \\
max\_ss slew, non-diode pins (distinct nets) & 5639 (314) & \textbf{97
(33)} \\
max\_ss slew composition ANT / other / SRAM & 1983 / 5609 / 30 & 198 /
81 / 16 \\
antenna nets / pins & 210 / 305 & \textbf{184 / 259} \\
setup worst slack (all corners) & +0.579 ns & \textbf{+3.276 ns} \\
hold worst slack & +0.279 ns & +0.275 ns \\
setup / hold violators & 0 / 0 & 0 / 0 \\
Netgen LVS & PASSED & PASSED \\
routing DRC & 0 & 0 \\
Magic DRC (bead \bead{45a}) & 8984 & 8984 \\
stdcell area / utilisation & 1 703 620 um2 / 44.84 \% & 1 778 150 um2 /
45.20 \% \\
power & 81.0 mW & 80.9 mW \\
fanout \textgreater{} 10 / \textgreater{} 32 / \textgreater{} 100 nets &
5181 / 1258 / 51 & 5839 / 694 / 26 \\
\end{longtable}
}

Max-cap is \textbf{0 at all nine corners}, which is the gate
\code{Checker.MaxCapViolations} would apply
(\texttt{TIMING\_\allowbreak{}VIOLATION\_\allowbreak{}CORNERS\ =\allowbreak{}\ ['*']}).
It is still \code{--skip}ped: the checker was not executed on this run,
only the metric read. Max-slew is not 0, so
\code{Checker.MaxSlewViolations} stays skipped.

Unlike the margin-raise (\code{e45j}, 2380 slew, +50 antenna nets) and
the antenna-effort raise (\code{58q}) that traded against each other,
this lever \textbf{improved both axes at once} and bought +2.7 ns of
setup margin on the thin max\_ss path. The antenna improvement is
\emph{inference-level} explained: fewer late-added repair buffers
because the repair was sized correctly the first time.

\end{record}

\subsection{4. What is left (max\_ss 295)}\label{cr:defect:sky130-slew-cap-e45j:4-what-is-left-max-ss-295}

\begin{record}

\begin{itemize}
\tightlist
\item
  16 pins on the SRAM macro inputs: checked against the relaxed 0.5 ns
  limit of \code{sky130_soc.sdc:120-124}; inherits bead \bead{o1i}, not
  a buffering problem.
\item
  198 ANTENNA diode pins at max\_ss (\emph{inference}: they sit on the
  same nets as real sinks, in line with the baseline pattern, and are
  the \code{58q} mirror effect; not re-mapped to nets in this run).
\item
  97 non-diode pins (81 non-SRAM) on 33 distinct nets; the worst have 8
  violating pins each (e.g. \code{net3131}, \code{net3132},
  \code{net23978}, \code{cpu_bridge_s_wdata[19]}). The nets are
  resizer-generated names and were not traced back to RTL signals.
\end{itemize}

Next experiments, each a single full run (\textasciitilde4 h on this
host): (a) raise the C multipliers a further \textasciitilde10 \% (the
median ratios are below the sum ratios used, but long nets dominate
slew); (b) set \code{GRT_DESIGN_REPAIR_MAX_WIRE_LENGTH} (untried; the
resizer auto-picks 6 335 um); (c) a real \code{set_max_fanout} in the
PnR SDC (changes repair behaviour: 5 839 nets exceed 10). A post-RCX ECO
pass stays an option only if (a)-(b) do not reach 0.

Single-run caveat: one full run per arm; no repeat was made, so
run-to-run noise on the slew counts is not characterised (the antenna
and setup differences are large against the historical drift on this
design).

\end{record}

\subsection{5. Confirmation run on current RTL (2026-10-08): NOT
adopted}\label{cr:defect:sky130-slew-cap-e45j:5-confirmation-run-on-current-rtl-2026-10-08-not-adopted}

\begin{record}

User decision 2026-10-08: adopt the RC-calibrated flow as the default
Sky130 SoC flow and un-skip \code{Checker.MaxCapViolations} \textbf{if}
a confirmation run on current \code{main} RTL matches section 3. Run:
\code{RUN_2026-10-08_05-21-39}
(\texttt{make\ -\allowbreak{}C\ pnr\ librelane-\allowbreak{}sky130-\allowbreak{}soc-\allowbreak{}rccal-\allowbreak{}noklayout},
\code{origin/main} at \code{3ad22f8}, no \code{dma_engine.sv} pin, so it
carries the \code{wdmo} DMA fix and the \code{3xtv}
\code{axi_lite_interconnect} fix; all 62 steps ran, sole deferred error
is the 8 984 Magic DRC of bead \bead{45a}). \textbf{Acceptance failed on
max-cap}, so no default was changed and \code{Checker.MaxCapViolations}
stays skipped.

{\small\setlength{\tabcolsep}{4pt}

{
\begin{longtable}[]{@{}
  >{\raggedright\arraybackslash}p{\dimexpr(\linewidth-6\tabcolsep)*175/1000\relax}
  >{\raggedright\arraybackslash}p{\dimexpr(\linewidth-6\tabcolsep)*208/1000\relax}
  >{\raggedright\arraybackslash}p{\dimexpr(\linewidth-6\tabcolsep)*274/1000\relax}
  >{\raggedright\arraybackslash}p{\dimexpr(\linewidth-6\tabcolsep)*323/1000\relax}@{}}
\toprule\noalign{}
\begin{minipage}[b]{\linewidth}\raggedright
(violating pins unless noted)
\end{minipage} & \begin{minipage}[b]{\linewidth}\raggedright
Gate B \code{RUN_2026-10-05_06-50-11}
\end{minipage} & \begin{minipage}[b]{\linewidth}\raggedright
rccal \code{RUN_2026-10-06_22-43-34} (pinned DMA)
\end{minipage} & \begin{minipage}[b]{\linewidth}\raggedright
rccal confirm \code{RUN_2026-10-08_05-21-39} (current RTL)
\end{minipage} \\
\midrule\noalign{}
\endhead
\bottomrule\noalign{}
\endlastfoot
max\_ss slew / cap & 7622 / 140 & 295 / 0 & \textbf{423 / 1} \\
nom\_ss & 2648 / 78 & 143 / 0 & 189 / \textbf{1} \\
min\_ss & 1669 / 33 & 5 / 0 & 16 / \textbf{1} \\
max\_tt & 2331 / 3 & 209 / 0 & 257 / \textbf{1} \\
nom\_tt & 1016 / 1 & 64 / 0 & 95 / \textbf{1} \\
min\_tt & 876 / 0 & 0 / 0 & 12 / \textbf{1} \\
max\_ff / nom\_ff / min\_ff slew & 1661 / 944 / 753 & 122 / 53 / 0 & 254
/ 46 / 11 (cap 0) \\
max\_ss slew ANT / other / SRAM & 1983 / 5609 / 30 & 198 / 81 / 16 & 304
/ 102 / 18 \\
antenna nets / pins & 210 / 305 & 184 / 259 & 205 / 273 \\
setup worst slack / violators & +0.579 ns / 0 & +3.276 ns / 0 & +3.329
ns / 0 \\
hold worst slack / violators & +0.279 ns / 0 & +0.275 ns / 0 & +0.282 ns
/ 0 \\
Netgen LVS / routing DRC & PASSED / 0 & PASSED / 0 & PASSED / 0 \\
Magic DRC (bead \bead{45a}) & 8984 & 8984 & 8984 \\
stdcell area / utilisation & 1 703 620 um2 / 44.84 \% & 1 778 150 um2 /
45.20 \% & 1 780 450 um2 / 45.21 \% \\
instances & 315 038 & 314 625 & 312 320 \\
power & 81.0 mW & 80.9 mW & 84.9 mW \\
\end{longtable}
}

}

\textbf{The single cap violation.} All six tt/ss corners fail on the
same driver, \code{_078722_/X}, a \code{sky130_fd_sc_hd__o22a_4} that
drives \code{cpu_bridge_s_rdata[12]} (the \code{cdc_gray_fifo} R-channel
read mux, \code{cdc_gray_fifo.sv} \code{mem_q[rd_bin_q]}, bead
\bead{7ohm}): 0.614 pF against a 0.530 pF limit at nom\_tt (+16 \%) and
0.643 pF against 0.333 pF at max\_ss; ff corners have no cap violation
(the same driver also appears in the nom\_tt slew list, 1.63 ns against
1.50 ns). \code{RepairDesignPostGRT} (step 37) logged "Found 296 slew
violations / 18 capacitance violations ... Inserted 34 buffers in 148
nets" and the estimate-side \code{nom_tt} metric was 0 afterwards, so
this net was inside the repair target on the calibrated GRT estimate and
exceeded it only after RCX (\emph{inference}: a detour or an unbuffered
long wire that the calibrated estimate under-counts on this one net; not
traced to a route length in this run).

\textbf{Reading.} Max-cap 0 at nine corners was \emph{not} a stable
property: it held on the pinned-DMA run and failed on one net on current
RTL. Un-skipping would therefore turn \code{librelane-sky130-soc} into a
failing flow (from \code{checker.py}
\code{TimingViolations.check_timing_violations}: any violating corner
matched by
\texttt{TIMING\_\allowbreak{}VIOLATION\_\allowbreak{}CORNERS\ =\allowbreak{}\ ['*']}
raises a \code{DeferredStepError}; read from source, the checker was not
executed; \textbf{this reading is wrong, see 7.1}). Slew and antenna
moved the wrong way against the previous rccal run, by 43 \% (max\_ss
295 -\textgreater{} 423), 2.1x (max\_ff 122 -\textgreater{} 254), +11 \%
antenna nets, but remain about 18x better than Gate B on max\_ss. Noise
reasoning: the netlist is different (312 320 vs 314 625 instances, -0.7
\%, different DMA and interconnect RTL), so placement and routing are a
different random draw; one run per arm cannot separate that from a real
effect. The deltas in slew and antenna are therefore uncharacterised,
not evidence of a regression in the flow; the cap violation is a hard
fail of the stated acceptance gate regardless of cause.

\textbf{Not done.} No default, \code{Makefile}, \code{config.json} or
checker skip was changed. The recommended next step is one more
single-variable run toward cap 0 (the +10 \% C multipliers, or
\code{GRT_DESIGN_REPAIR_MAX_WIRE_LENGTH}), and a repeat of the
\emph{same} config on the same RTL to measure run-to-run noise before
adopting anything.

\end{record}

\subsection{\texorpdfstring{6. Noise repeat and the wire-length
experiment (2026-10-08/\allowbreak{}09): still NOT
adopted}{6. Noise repeat and the wire-length experiment (2026-10-08/09): still NOT adopted}}\label{cr:defect:sky130-slew-cap-e45j:6-noise-repeat-and-the-wire-length-experiment-2026-10-0809-still-not-adopted}

\begin{record}

User decision 2026-10-08: experiment first, adopt \code{config_rccal} as
default (and un-skip \code{Checker.MaxCapViolations}) only if max-cap is
0 at all nine corners with LVS PASSED, routing DRC 0,
setup/\allowbreak{}hold 0 violators, and slew/\allowbreak{}antenna no
worse than the confirmation run within measured noise. Both runs below
used \code{origin/main} \code{944d8c5}, whose \src{rtl/} and \src{pnr/}
are identical to the confirmation run\textquotesingle s \code{3ad22f8}
(\texttt{git\ diff\ -\allowbreak{}-\allowbreak{}stat\ 3ad22f8\ 944d8c5\ -\allowbreak{}-\allowbreak{}\ rtl\ pnr}
is empty; the intervening commits are tests and docs).

\end{record}

\subsubsection{6.1 Noise: the flow is bit-deterministic here (measured)}\label{cr:defect:sky130-slew-cap-e45j:61-noise-the-flow-is-bit-deterministic-here-measured}

\begin{record}

Run 1 (\code{RUN_2026-10-08_21-07-20}, exact \code{config_rccal.json},
same RTL, a different worktree path) was compared with the confirmation
run \code{RUN_2026-10-08_05-21-39}. After normalising the worktree path
embedded in net names, the synthesis netlist and the DEFs after global
placement, CTS, global routing, step 37, step 40 and detailed routing
have identical checksums; the step-37 repair logs match line for line
(296 slew / 18 cap violations, 34 buffers in 148 nets, 40 resized); and
every step-51 metric is equal: slew at all nine corners (max\_ss 423 ...
min\_ff 11), cap (1 at the six tt/ss corners), antenna 205 nets / 273
pins, 312 320 instances, 45.2134 \% utilisation, 84.89 mW, setup +3.3286
/ hold +0.2822 ns, route wirelength 7 502 911. So the run-to-run noise
on slew, cap and antenna is \textbf{zero for identical inputs}: every
difference in a changed-input run is a real response to the input,
though the response of a chaotic flow to a small input change is still
one sample per arm. Run 1 was stopped at step 52 (IR-drop report) once
the post-RCX metrics proved identical; steps 53-62 (Magic, LVS, DRC)
were \textbf{not} repeated (LVS PASSED / Magic 8 984 on the confirmation
run carry over by determinism: \emph{inference}).

\end{record}

\subsubsection{6.2 Why the one cap violation survives (measured on the
confirmation run)}\label{cr:defect:sky130-slew-cap-e45j:62-why-the-one-cap-violation-survives-measured-on-the-confirmation-run}

\begin{record}

The single violating driver \code{_078722_}
(\code{sky130_fd_sc_hd__o22a_4}, max\_\allowbreak{}capacitance 0.5301
pF) drives \code{cpu_bridge_s_rdata[12]}. Everything about this net is
an outlier:

\begin{itemize}
\tightlist
\item
  \textbf{Load}: one functional sink, the CPU macro pin
  \code{axi_rdata_i[12]}, plus 12 \code{diode_2} antenna cells (0.0009
  pF each). The macro\textquotesingle s Liberty gives that pin
  \textbf{0.2464 pF}, against a median of 0.0083 pF over its 92 input
  pins (next largest 0.1305 pF \code{apb_paddr_i[4]}, then 0.0658 pF).
  The post-RCX SPEF wire cap is 0.358 pF. 0.358 + 0.2464

  \begin{itemize}
  \tightlist
  \item
    12 x 0.0009 = 0.615 pF, matching the 0.614 pF reported.
  \end{itemize}
\item
  \textbf{Length}: 2 196.5 um routed, rank 54 of 117 008 nets (p99 = 854
  um, p99.9 = 1 871 um). The global-routing guide for the net at step 37
  already spans 1 518 x 669 um, so the calibrated estimate saw the same
  length: this is not an under-counted detour.
\item
  \textbf{Repair never touched it}: the net is
  \texttt{o22a\_\allowbreak{}4\ -\allowbreak{}>\ macro\ pin} unchanged
  at step 33 and 37; only the antenna step added the 11-12 diodes (they
  add \textasciitilde0.01 pF, not the cause). The driver is already the
  largest drive in its family.
\end{itemize}

So the cause is a long wire into a macro input whose characterised pin
capacitance is \textasciitilde30x the median, not an RC-estimator error:
at \textasciitilde0.16 fF/um (the calibrated nom value) the estimate for
2 190 um is \textasciitilde0.35 pF, which together with the 0.246 pF pin
already exceeds 0.53 pF. The earlier "inference" in section 5 (a detour
that the estimate under-counts) is superseded.

\end{record}

\subsubsection{\texorpdfstring{6.3 The single-variable run:
\texttt{GRT\_\allowbreak{}DESIGN\_\allowbreak{}REPAIR\_\allowbreak{}MAX\_\allowbreak{}WIRE\_\allowbreak{}LENGTH\ =\allowbreak{}\ 1500}}{6.3 The single-variable run: }}\label{cr:defect:sky130-slew-cap-e45j:63-the-single-variable-run-grt-design-repair-max-wire-length-1500}

\begin{record}

Choice (justified from 6.2): the failing net is a long point-to-point
wire, not high fanout, and its estimate is already calibrated, so a
further +10 \% on the C multipliers would barely move a net that is 15
\% over the limit by a different mechanism, and would perturb every net.
A wire-length cap targets exactly long wires. The stock auto value is 3
205 um (\code{RSZ-0058} in the step-37 log), which never touches a 2 196
um net; 1 500 um is the smallest perturbation that splits it.
\code{config_rccal_wl1500.json} is \code{config_rccal.json} plus that
one key (verified by \code{diff} and by \code{resolved.json});
\code{config.json} and \code{config_rccal.json} are unchanged. Run:
\code{RUN_2026-10-08_23-12-19}, all 62 steps, same host,
\texttt{make\ -\allowbreak{}C\ pnr\ librelane-\allowbreak{}sky130-\allowbreak{}soc\ SKY130\_\allowbreak{}SOC\_\allowbreak{}CONFIG=\allowbreak{}config\_\allowbreak{}rccal\_\allowbreak{}wl1500.\allowbreak{}json\ SKY130\_\allowbreak{}SOC\_\allowbreak{}EXTRA\_\allowbreak{}SKIP=\allowbreak{}"-\allowbreak{}-\allowbreak{}skip\ KLayout.\allowbreak{}DRC\ -\allowbreak{}-\allowbreak{}skip\ Checker.\allowbreak{}KLayoutDRC"}.
(Passing \code{SKY130_SOC_CONFIG=} to the \code{-rccal-noklayout} target
does \textbf{not} work: that recipe re-assigns it, and a first launch
silently ran the plain rccal config; it was caught from
\code{resolved.json}, killed after minutes, and relaunched.)

{\footnotesize\setlength{\tabcolsep}{3pt}

{
\begin{longtable}[]{@{}
  >{\raggedright\arraybackslash}p{\dimexpr(\linewidth-8\tabcolsep)*183/1000\relax}
  >{\raggedright\arraybackslash}p{\dimexpr(\linewidth-8\tabcolsep)*162/1000\relax}
  >{\raggedright\arraybackslash}p{\dimexpr(\linewidth-8\tabcolsep)*209/1000\relax}
  >{\raggedright\arraybackslash}p{\dimexpr(\linewidth-8\tabcolsep)*264/1000\relax}
  >{\raggedright\arraybackslash}p{\dimexpr(\linewidth-8\tabcolsep)*162/1000\relax}@{}}
\toprule\noalign{}
\begin{minipage}[b]{\linewidth}\raggedright
(violating pins unless noted)
\end{minipage} & \begin{minipage}[b]{\linewidth}\raggedright
Gate B \code{RUN_2026-10-05_06-50-11}
\end{minipage} & \begin{minipage}[b]{\linewidth}\raggedright
rccal, pinned DMA \code{RUN_2026-10-06_22-43-34}
\end{minipage} & \begin{minipage}[b]{\linewidth}\raggedright
confirmation \code{RUN_2026-10-08_05-21-39} = noise repeat
\code{RUN_2026-10-08_21-07-20}
\end{minipage} & \begin{minipage}[b]{\linewidth}\raggedright
wl1500 \code{RUN_2026-10-08_23-12-19}
\end{minipage} \\
\midrule\noalign{}
\endhead
\bottomrule\noalign{}
\endlastfoot
max\_ss slew / cap & 7622 / 140 & 295 / 0 & 423 / 1 & \textbf{102 /
1} \\
nom\_ss & 2648 / 78 & 143 / 0 & 189 / 1 & 21 / 1 \\
min\_ss & 1669 / 33 & 5 / 0 & 16 / 1 & 17 / 1 \\
max\_tt & 2331 / 3 & 209 / 0 & 257 / 1 & 16 / 1 \\
nom\_tt & 1016 / 1 & 64 / 0 & 95 / 1 & 13 / 1 \\
min\_tt & 876 / 0 & 0 / 0 & 12 / 1 & 12 / 1 \\
max\_ff / nom\_ff / min\_ff slew (cap 0) & 1661 / 944 / 753 & 122 / 53 /
0 & 254 / 46 / 11 & 13 / 11 / 0 \\
max\_ss slew ANT / SRAM / other & 1983 / 30 / 5609 & 198 / 16 / 81 & 304
/ 18 / 101 & 54 / 18 / 30 \\
antenna nets / pins & 210 / 305 & 184 / 259 & 205 / 273 & 212 / 253 \\
setup worst slack / violators & +0.579 ns / 0 & +3.276 / 0 & +3.329 / 0
& +3.345 / 0 \\
hold worst slack / violators & +0.279 ns / 0 & +0.275 / 0 & +0.282 / 0 &
+0.274 / 0 \\
Netgen LVS / routing DRC & PASSED / 0 & PASSED / 0 & PASSED / 0 & PASSED
/ 0 \\
Magic DRC (bead \bead{45a}) & 8984 & 8984 & 8984 & 8984 \\
utilisation / instances & 44.84 \% / 315 038 & 45.20 \% / 314 625 &
45.21 \% / 312 320 & 45.23 \% / 312 244 \\
power & 81.0 mW & 80.9 mW & 84.9 mW & 84.9 mW \\
step-37 repair (buffers / nets / resized) & - & - & 34 / 148 / 40 & 200
/ 257 / 128 \\
\end{longtable}
}

}

Result: \textbf{slew improves sharply with no cost elsewhere} (max\_ss
423 -\textgreater{} 102, max\_tt 257 -\textgreater{} 16, max\_ff 254
-\textgreater{} 13; ANT diode pins 304 -\textgreater{} 54; antenna pins
273 -\textgreater{} 253 while nets 205 -\textgreater{} 212;
setup/\allowbreak{}hold/\allowbreak{}LVS/\allowbreak{}DRC unchanged
within a few ps), \textbf{but max-cap stays at 1 at the same six corners
on the same net}: \code{_078722_/X} is 0.646 pF against 0.530 pF at
nom\_tt, slightly worse than before (0.614), and the net is again
\texttt{o22a\_\allowbreak{}4\ -\allowbreak{}>\ macro\ pin} with no
buffer at step 37 (checked in the step-37 ODB). The adoption bar (cap 0
at all nine corners) is \textbf{not met}, so no default changed,
\code{Checker.MaxCapViolations} stays skipped, and \code{config_rccal}
is not promoted. (Single arm for a changed input, with zero repeat
noise; why the resizer does not buffer this particular 2 196 um net
although it buffered 200 others under the cap was not established.
\emph{Inference}: the net\textquotesingle s far end is a hard-macro pin
and repair\_\allowbreak{}design does not rebuffer toward it.)

\end{record}

\subsubsection{6.4 Recommendation and next step}\label{cr:defect:sky130-slew-cap-e45j:64-recommendation-and-next-step}

\begin{record}

\emph{Superseded by section 7: the user chose option (1) below (named
waiver) and adopted rccal on 2026-10-09.}

\begin{itemize}
\tightlist
\item
  \textbf{Do not adopt yet.} Adoption needs cap 0 and it is 1, so the
  decision rule is unchanged.
\item
  The residual is a \textbf{single net of a different kind} than the
  previous 8 623 -\textgreater{} 295 slew story: a cap violation caused
  by one macro input pin characterised at 0.2464 pF. Config-level sweeps
  are unlikely to fix it (two runs: auto-length and 1 500 um wire-length
  both leave it untouched). Options, cheapest first, none executed: (1)
  accept it as a tracked single-net residual and un-skip the cap checker
  only with a documented, named waiver for this net (a waiver is a moved
  target, user\textquotesingle s call); (2) look at the CPU macro: why
  does \code{axi_rdata_i[12]} present 0.2464 pF when its siblings
  present \textasciitilde0.01 pF? An input buffer inside the macro (a
  CPU macro re-harden) would fix it at the source, but that edits a
  macro view and needs explicit approval (cf. bead \bead{ma7}/\code{lxv}
  hard-macro rule); (3) constrain placement of \code{_078722_} next to
  the macro pin (a pin-specific placement constraint, brittle across
  re-synthesis names).
\item
  \texttt{GRT\_\allowbreak{}DESIGN\_\allowbreak{}REPAIR\_\allowbreak{}MAX\_\allowbreak{}WIRE\_\allowbreak{}LENGTH\ =\allowbreak{}\ 1500}
  is independently attractive for \textbf{slew} (max\_ss 423
  -\textgreater{} 102 at no measured cost): worth adopting together with
  rccal when rccal itself is adopted, after one more arm on a different
  RTL vintage to see it is not a one-draw effect. Not adopted now.
\end{itemize}

\end{record}

\subsection{7. Adoption (2026-10-09): rccal is the default; max-cap
gated with a named waiver}\label{cr:defect:sky130-slew-cap-e45j:7-adoption-2026-10-09-rccal-is-the-default-max-cap-gated-with-a-named-waiver}

\begin{record}

\textbf{Decision (user, 2026-10-09).} Adopt the RC-calibrated flow as
the default Sky130 SoC flow, un-skip the max-cap check, waive the one
residual net by name, and do \textbf{not} adopt
\texttt{GRT\_\allowbreak{}DESIGN\_\allowbreak{}REPAIR\_\allowbreak{}MAX\_\allowbreak{}WIRE\_\allowbreak{}LENGTH\ =\allowbreak{}\ 1500}
(it needs one more arm on another RTL vintage). Max-slew stays skipped
(102-423 pins at max\_ss, tracked here).

\textbf{This is a moved target, and it is the user\textquotesingle s to
move.} The cap gate was "0 at nine corners". It is now "0 except the
waived net". The justification is in
\src{pnr/sky130/soc/constraints/max_cap_waivers.json} and in section
6.2: the sole functional load is a CPU-macro input whose Liberty pin
capacitance (0.2464 pF) is \textasciitilde30x the median of the
macro\textquotesingle s inputs and 46 \% of the driver\textquotesingle s
0.5301 pF limit, so no repair on the SoC side can fix it. The real fix
(an input buffer inside the CPU macro) edits a hard-macro view and still
needs approval (cf. beads \code{ma7}/\code{lxv}).

\end{record}

\subsubsection{7.1 Un-skipping the stock checker alone would have gated
nothing (measured)}\label{cr:defect:sky130-slew-cap-e45j:71-un-skipping-the-stock-checker-alone-would-have-gated-nothing-measured}

\begin{record}

Section 5 said the stock checker would fail on any violating corner
matched by
\texttt{TIMING\_\allowbreak{}VIOLATION\_\allowbreak{}CORNERS\ =\allowbreak{}\ ['*']};
that was read from source and is \textbf{wrong}.
\code{Checker.MaxCapViolations} reads \code{MAX_CAP_VIOLATION_CORNERS}
first, and its class sets
\texttt{corner\_\allowbreak{}override\ =\allowbreak{}\ [""]}
(\code{librelane/steps/checker.py}), which matches no corner. Executed
with
\texttt{-\allowbreak{}-\allowbreak{}only\ Checker.\allowbreak{}MaxCapViolations}
on the finished run\textquotesingle s state
(\code{RUN_2026-10-08_05-21-39}, 6 corners with cap = 1):

\begin{itemize}
\tightlist
\item
  default: \textbf{exit 0},
  \texttt{WARNING\ Max\ Cap\ violations\ found\ in\ the\ following\ corners:\allowbreak{}\ .\allowbreak{}.\allowbreak{}.\allowbreak{}\ max\_\allowbreak{}ss\ .\allowbreak{}.\allowbreak{}.\allowbreak{}\ nom\_\allowbreak{}tt}
  then \texttt{VERBOSE\ No\ max\ cap\ violations\ found};
\item
  with
  \texttt{-\allowbreak{}c\ MAX\_\allowbreak{}CAP\_\allowbreak{}VIOLATION\_\allowbreak{}CORNERS=\allowbreak{}["*"]}:
  exit 2 (\texttt{One\ or\ more\ deferred\ errors}), but then there is
  no way to waive one net, because the checker only sees a per-corner
  count.
\end{itemize}

So "un-skip" without more work is a silent no-op, and "un-skip with
corners" fails every run on the known net.

\end{record}

\subsubsection{7.2 What was built}\label{cr:defect:sky130-slew-cap-e45j:72-what-was-built}

\begin{record}

\begin{itemize}
\tightlist
\item
  \src{pnr/sky130/soc/plugin/cvt_maxcap_waiver.py} (pure Python, no
  LibreLane import; also a CLI): parses the \texttt{max\ capacitance}
  table of each corner\textquotesingle s \code{checks.rpt}, resolves
  each violating driver\textquotesingle s net in the final netlist, and
  waives a violation only if that net feeds a \textbf{macro pin} named
  by a waiver (\code{u_cpu*} / \code{axi_rdata_i[12]}) \textbf{and} its
  capacitance is at or below the waiver\textquotesingle s
  \code{max_cap_pf} ceiling (0.70 pF). Keyed on the macro pin, not on
  the synthesis-generated driver name (\code{_078722_}), which goes
  stale at the next re-synthesis. Bus-port bits are taken from the
  netlist\textquotesingle s concatenation order (first element = MSB);
  for this run that gave \code{axi_rdata_i[12]} independently of the
  earlier hand analysis.
\item
  \code{CVT.MaxCapViolations} step in the project plugin, substituted
  for \code{Checker.MaxCapViolations} in \code{config_rccal.json} (and
  \code{config_rccal_wl1500.json}) via \code{meta.substituting_steps},
  waiver file via \code{CVT_MAXCAP_WAIVERS}. It gates \textbf{all}
  corners, raises a deferred error (flow finishes, like the other
  checkers) on any unwaived violation, warns on stale waivers, writes
  \code{max_cap_waivers.rpt} in its step dir and emits
  \texttt{design\_\allowbreak{}\_\allowbreak{}max\_\allowbreak{}cap\_\allowbreak{}violation\_\allowbreak{}\_\allowbreak{}\{waived,\allowbreak{}unwaived\}\_\allowbreak{}\_\allowbreak{}count}.
\item
  Fail-closed inputs (each unit-tested): a report without its
  \texttt{max\ cap\ violation\ count\ N} line, a row count that
  disagrees with the summary line or with the flow\textquotesingle s
  per-corner metric, a corner with a metric but no report, a violation
  whose driver is not in the netlist, and violations with no netlist are
  all errors, never passes. A waiver that waives nothing is reported
  STALE (\code{--strict-stale} / \code{CVT_MAXCAP_STRICT_STALE} makes it
  fail). A waiver without a justification, with a duplicate id, or with
  a non-positive ceiling is rejected.
\item
  \src{pnr/Makefile}: \code{SKY130_SOC_CONFIG} now defaults to
  \code{config_rccal.json}, so \code{librelane-sky130-soc} and
  \code{-noklayout} run the adopted flow. \code{-norccal} /
  \code{-norccal-noklayout} run the pre-adoption \code{config.json}
  (stock max-cap checker skipped there, as it is a no-op anyway).
  \code{SKY130_SOC_CONFIG=} is now honoured by every target
  (\code{-rccal-noklayout} is an alias of \code{-noklayout}; the old
  recipe silently ignored the override).
  \texttt{make\ -\allowbreak{}C\ pnr\ sky130-\allowbreak{}soc-\allowbreak{}maxcap-\allowbreak{}check\ [SKY130\_\allowbreak{}SOC\_\allowbreak{}MAXCAP\_\allowbreak{}RUN=\allowbreak{}<run>]\ [STRICT\_\allowbreak{}STALE=\allowbreak{}1]}
  re-runs the check on any finished run.
\end{itemize}

\end{record}

\subsubsection{7.3 Validation (no new full flow was run; see 7.4)}\label{cr:defect:sky130-slew-cap-e45j:73-validation-no-new-full-flow-was-run-see-74}

\begin{record}

No full flow was run because the finished runs already carry the exact
flow inputs for everything except the final checker step, and the flow
is bit-deterministic here (6.1). The new check was applied to those
artifacts:

{\footnotesize\setlength{\tabcolsep}{3pt}

{
\begin{longtable}[]{@{}
  >{\raggedright\arraybackslash}p{\dimexpr(\linewidth-4\tabcolsep)*355/1000\relax}
  >{\raggedright\arraybackslash}p{\dimexpr(\linewidth-4\tabcolsep)*269/1000\relax}
  >{\raggedright\arraybackslash}p{\dimexpr(\linewidth-4\tabcolsep)*355/1000\relax}@{}}
\toprule\noalign{}
\begin{minipage}[b]{\linewidth}\raggedright
run / case
\end{minipage} & \begin{minipage}[b]{\linewidth}\raggedright
waiver list
\end{minipage} & \begin{minipage}[b]{\linewidth}\raggedright
result
\end{minipage} \\
\midrule\noalign{}
\endhead
\bottomrule\noalign{}
\endlastfoot
\code{RUN_2026-10-08_05-21-39} (rccal, current RTL), CLI &
\code{max_cap_waivers.json}, \code{--strict-stale} & \textbf{PASS}: 9
corners, 6 waived (\code{_078722_/X}, net \code{cpu_bridge_s_rdata[12]},
0.5600-0.6427 pF), 0 unwaived, 0 stale \\
same run, CLI & \textbf{empty} (negative control) & \textbf{FAIL}, exit
1: 6 UNWAIVED \\
\code{RUN_2026-10-08_21-07-20} (noise repeat) &
\code{max_cap_waivers.json} & PASS (identical numbers; bit-identical
flow) \\
\code{RUN_2026-10-08_23-12-19} (wl1500) & \code{max_cap_waivers.json} &
PASS: waived caps 0.5874-0.6761 pF, all under the 0.70 pF ceiling \\
\code{RUN_2026-10-05_06-50-11} (Gate B, no rccal) &
\code{max_cap_waivers.json} & \textbf{FAIL}: 257 unwaived, 1 stale
waiver \\
in-flow,
\texttt{librelane\ -\allowbreak{}-\allowbreak{}only\ CVT.\allowbreak{}MaxCapViolations}
with \code{config_rccal.json} on
\code{RUN_2026-10-08_05-21-39}\textquotesingle s state &
\code{max_cap_waivers.json} & \textbf{exit 0}, 6 waived / 0 unwaived,
metrics \texttt{waived=\allowbreak{}6\ unwaived=\allowbreak{}0} in
\code{state_out.json} \\
same in-flow run,
\texttt{-\allowbreak{}c\ CVT\_\allowbreak{}MAXCAP\_\allowbreak{}WAIVERS=\allowbreak{}<empty\ list>}
& empty & \textbf{exit 2},
\texttt{6\ unwaived\ max-\allowbreak{}cap\ violation(s)\ .\allowbreak{}.\allowbreak{}.\allowbreak{}\ FAIL\ -\allowbreak{}\ deferred} \\
\end{longtable}
}

}

Unit tests: \src{tb/tests/test_cvt_maxcap_waiver.py}, 50 tests written
first (RED: module missing), then green. Two mutations were applied by
hand to prove they bite (reversing the bus-bit order: 9 failures;
removing the ceiling test: 1 failure), then reverted.

\end{record}

\subsubsection{7.4 What is NOT validated}\label{cr:defect:sky130-slew-cap-e45j:74-what-is-not-validated}

\begin{record}

\begin{itemize}
\tightlist
\item
  \sout{No full flow has run with the adopted config.} \textbf{Resolved
  by 7.5} (end-to-end run 2026-10-09: the substituted step ran in
  position and gave the same verdict as the \code{--only} validation).
\item
  All adopted-flow numbers are one RTL vintage
  (\code{3ad22f8}/\code{944d8c5}). On a new vintage the net name
  changes, which is why the waiver is keyed on the macro pin, but the
  cap may move: the ceiling (0.70 pF; 0.2464 pF fixed pin cap + wire)
  fails loudly if the wire grows by \textasciitilde25 \%, and a STALE
  warning appears if the net stops violating.
\item
  Max-slew is still skipped (\code{Checker.MaxSlewViolations}): 102-423
  pins at max\_ss, antenna diodes 54-304 of them.
\item
  Antenna (\code{Odb.CheckDesignAntennaProperties}) and fanout (no
  \code{set_max_fanout} in the SDC, so "0 fanout violations" stays
  vacuous) are unchanged.
\end{itemize}

\end{record}

\subsubsection{\texorpdfstring{7.5 End-to-end confirmation (2026-10-09):
the adopted flow completes with \code{CVT.MaxCapViolations} in
position}{7.5 End-to-end confirmation (2026-10-09): the adopted flow completes with  in position}}\label{cr:defect:sky130-slew-cap-e45j:75-end-to-end-confirmation-2026-10-09-the-adopted-flow-completes-with-cvtmaxcapviolations-in-position}

\begin{record}

One full
\texttt{make\ -\allowbreak{}C\ pnr\ librelane-\allowbreak{}sky130-\allowbreak{}soc-\allowbreak{}noklayout}
(no overrides, so \code{SKY130_SOC_CONFIG=config_rccal.json}) on
\code{origin/main} \code{6190313} (\prref{239} merged), run
\code{RUN_2026-10-09_06-18-39}, 06:18:31 -\textgreater{} 09:44:08 (3 h
26 min wall). RTL is \textbf{unchanged} since the reference: the last
commit touching \src{rtl/} is \code{ce2b1fc} (2026-10-07), before
reference run \code{RUN_2026-10-08_05-21-39}; the commits since
\code{3ad22f8} are tests, docs and the PD flow files only (the
\code{rqvo} UART fix is not on main). Metrics below are the last
\code{state_out.json} of each run (step 63 / step 62).

{\footnotesize\setlength{\tabcolsep}{3pt}

{
\begin{longtable}[]{@{}
  >{\raggedright\arraybackslash}p{\dimexpr(\linewidth-4\tabcolsep)*254/1000\relax}
  >{\raggedright\arraybackslash}p{\dimexpr(\linewidth-4\tabcolsep)*360/1000\relax}
  >{\raggedright\arraybackslash}p{\dimexpr(\linewidth-4\tabcolsep)*366/1000\relax}@{}}
\toprule\noalign{}
\begin{minipage}[b]{\linewidth}\raggedright
metric (all-corner unless stated)
\end{minipage} & \begin{minipage}[b]{\linewidth}\raggedright
reference \code{RUN_2026-10-08_05-21-39}
\end{minipage} & \begin{minipage}[b]{\linewidth}\raggedright
end-to-end \code{RUN_2026-10-09_06-18-39}
\end{minipage} \\
\midrule\noalign{}
\endhead
\bottomrule\noalign{}
\endlastfoot
max-slew, max\_ss / max\_tt / nom\_tt (pins) & 423 / 257 / 95 (all nine:
423 254 257 11 16 12 46 189 95) & \textbf{identical} at all nine
corners \\
max-cap, raw (6 corners tt/ss) & 1 per corner (\code{_078722_}, net
\code{cpu_bridge_s_rdata[12]}) & \textbf{identical} \\
max-cap, waived / unwaived & not computable in-flow (stock checker
skipped) & \textbf{6 waived / 0 unwaived / 0 stale},
\texttt{max-\allowbreak{}cap\ check:\allowbreak{}\ PASS}; waived caps
0.5600-0.6427 pF, all under the 0.70 pF ceiling \\
antenna (nets / pins) & 205 / 273, 10 918 diodes & \textbf{identical} \\
setup worst slack, per corner & +3.3286 ns (max\_ss), +10.0518
(nom\_tt); 0 violators & \textbf{identical}, 0 violators at all nine
corners \\
hold worst slack, per corner & +0.2822 ns (min\_ff), +0.5397 (nom\_tt);
0 violators & \textbf{identical}, 0 violators at all nine corners \\
LVS & 0 errors, PASSED & \textbf{identical} (all \code{design__lvs_*} =
0) \\
routing DRC & 0 (iterations 59021, 21775, 18600, 1095, 67, 2, 0) &
\textbf{identical}, same iteration trail \\
Magic DRC & 8 984 (waived PDK artifact, bead \bead{45a}) & 8 984
(\textbf{identical}) \\
KLayout DRC & skipped & skipped
(\texttt{-\allowbreak{}-\allowbreak{}skip\ KLayout.\allowbreak{}DRC\ -\allowbreak{}-\allowbreak{}skip\ Checker.\allowbreak{}KLayoutDRC},
deferred, not waived) \\
util / instances / power & 45.2134 \% / 312 320 / 84.89 mW &
identical \\
\end{longtable}
}

}

\textbf{Did \code{CVT.MaxCapViolations} run in-flow, and what did it
say?} Yes. \code{flow.log} line 650:
\texttt{Running\ 'CVT.\allowbreak{}MaxCapViolations'} as step
\textbf{62}, in the slot the stock \code{Checker.MaxCapViolations}
occupied (after 61 \code{Checker.HoldViolations}, before 63
\code{Misc.ReportManufacturability}); \code{resolved.json} shows
\code{meta.substituting_steps}
\texttt{Checker.\allowbreak{}MaxCapViolations\ -\allowbreak{}>\ CVT.\allowbreak{}MaxCapViolations}
and \code{CVT_MAXCAP_WAIVERS} pointing at
\code{constraints/max_cap_waivers.json}; the stock step does not appear.
Runtime 4.4 s. Verdict \textbf{PASS} (the six violations, all
\code{_078722_/X} -\textgreater{} \code{cpu_bridge_s_rdata[12]}, are
waived by \code{cpu-axi-rdata-i12-pin-cap}; unwaived count 0 in
\code{design__max_cap_violation__unwaived__count}). The
step\textquotesingle s WARNING is re-emitted by \code{flow.py} and shows
up in the manufacturability summary, as designed. This matches the
\code{--only} validation of 7.3 (6 waived / 0 unwaived, same caps to
four digits), so the 7.4 inference is confirmed.

\textbf{Bit-identical to the reference?} Yes, modulo the worktree path.
After replacing the embedded \code{agent-a<hash>} worktree path (it
appears in sv2v-derived net names) and dropping \code{//} header lines,
the SHA-256 of the final netlist
(\code{48-openroad-fillinsertion/soc_top.nl.v}) and of the final DEF
(\code{49-odb-cellfrequencytables/soc_top.def}) are equal between the
two runs (\code{746cd83c9917a19b}, \code{ed09a01574b539c1}); the SDC is
byte-equal. Of 292 metrics keys in the last state, 289 are equal; the 3
that differ are the two new
\texttt{design\_\allowbreak{}\_\allowbreak{}max\_\allowbreak{}cap\_\allowbreak{}violation\_\allowbreak{}\_\allowbreak{}\{waived,\allowbreak{}unwaived\}\_\allowbreak{}\_\allowbreak{}count}
keys (absent in the reference, expected) and
\code{design_powergrid__drop__average__net:VGND__corner:nom_tt}
(1.63812e-05 vs 1.63834e-05 V, +0.013 \%, an IR-analysis readout; no
timing or physical metric moved). Second confirmation of
6.1\textquotesingle s determinism, now across a third worktree path and
with the new step in the flow.

\textbf{Exit status.} \code{make} returned 2 and LibreLane reported
\texttt{One\ or\ more\ deferred\ errors\ were\ encountered:\allowbreak{}\ 8984\ Magic\ DRC\ errors\ found.\allowbreak{}}
This is the \textbf{same deferred error as the reference run} (its
\code{error.log} ends with the same line): the Magic DRC count is the
waived PDK artifact of bead \bead{45a}, deferred errors let the flow
reach step 63. It is not a new failure and not caused by
\code{CVT.MaxCapViolations}. Anyone gating on the make exit code needs
to know that the adopted Sky130 SoC flow exits non-zero until the Magic
DRC waiver is turned into a checker-level exemption (not done here;
untouched).

\textbf{Setup finding (worktree runs).} The first launch aborted in
\textasciitilde9 s at config load
(\texttt{PermissionError:\allowbreak{}\ .\allowbreak{}.\allowbreak{}.\allowbreak{}\ cpu/\allowbreak{}macro/\allowbreak{}\ rv32i\_\allowbreak{}cpu\_\allowbreak{}top.\allowbreak{}gds\ is\ not\ located\ any\ path\ readable\ to\ LibreLane})
before any run directory existed: the macro \code{*.gds} /
\code{*.spice} views are gitignored
(\texttt{pnr/\allowbreak{}sky130/\allowbreak{}cpu/\allowbreak{}macro/\allowbreak{}rv32i\_\allowbreak{}cpu\_\allowbreak{}top.\allowbreak{}\{gds,\allowbreak{}spice\}},
\texttt{pnr/\allowbreak{}sky130/\allowbreak{}soc/\allowbreak{}macro/\allowbreak{}sky130\_\allowbreak{}sram\_\allowbreak{}4kbyte\_\allowbreak{}1rw1r\_\allowbreak{}32x1024\_\allowbreak{}8.\allowbreak{}\{gds,\allowbreak{}sp\}})
and exist only in the primary checkout. A fresh worktree needs them
copied
(\texttt{rsync\ -\allowbreak{}a\ -\allowbreak{}-\allowbreak{}ignore-\allowbreak{}existing\ <primary>/\allowbreak{}pnr/\allowbreak{}sky130/\allowbreak{}\{cpu,\allowbreak{}soc\}/\allowbreak{}macro/\allowbreak{}\ .\allowbreak{}.\allowbreak{}.\allowbreak{}})
before \code{make} can start. Not a flow defect; the copied files are
the same inputs the primary-tree runs use.

Still open (unchanged by this run): max-slew is skipped (423 at
max\_ss), the waiver is for the one macro-pin net, and a new RTL vintage
may move the cap (ceiling 0.70 pF; STALE warning if the net stops
violating).

\end{record}

\subsection{\texorpdfstring{8. New baseline on the sv2v-rebuilt CPU
macro (2026-10-11, bead \code{dud4}): stale waiver, 3 unwaived ss-corner
nets}{8. New baseline on the sv2v-rebuilt CPU macro (2026-10-11, bead ): stale waiver, 3 unwaived ss-corner nets}}\label{cr:defect:sky130-slew-cap-e45j:8-new-baseline-on-the-sv2v-rebuilt-cpu-macro-2026-10-11-bead-dud4-stale-waiver-3-unwaived-ss-corner-nets}

\begin{record}

\code{RUN_2026-10-11_08-39-41}
(\texttt{make\ -\allowbreak{}C\ pnr\ librelane-\allowbreak{}sky130-\allowbreak{}soc-\allowbreak{}rccal-\allowbreak{}noklayout},
08:39 -\textgreater{} 12:11 WIB, 3 h 32 min) is the first SoC run on the
re-hardened CPU macro (\cref{cr:defect:sky130-cpu-synlig-dud4},
"Re-harden"). \textbf{This is a new baseline, not a single-variable
comparison with \code{RUN_2026-10-09_06-18-39} (§7.5):} the CPU macro
frontend (Synlig -\textgreater{} sv2v), the macro\textquotesingle s
registered APB outputs (old macro predates \code{999e44a}), the
hazard-select RTL fix (\prref{259}), the csr counter fix (\prref{268}) and the
UART/\allowbreak{}I2C/\allowbreak{}GPU RTL changes since all differ.
\code{CVT_SYNTH_GATE=warn} was exported in the shell for this run (see
the PR body: the \code{gc0y} structural check false-positives on 509
sv2v function-inlining temporaries in the SoC); no tracked file carries
it. KLayout DRC skipped as in §7.5.

{\small\setlength{\tabcolsep}{4pt}

{
\begin{longtable}[]{@{}
  >{\raggedright\arraybackslash}p{\dimexpr(\linewidth-4\tabcolsep)*395/1000\relax}
  >{\raggedright\arraybackslash}p{\dimexpr(\linewidth-4\tabcolsep)*272/1000\relax}
  >{\raggedright\arraybackslash}p{\dimexpr(\linewidth-4\tabcolsep)*313/1000\relax}@{}}
\toprule\noalign{}
\begin{minipage}[b]{\linewidth}\raggedright
metric
\end{minipage} & \begin{minipage}[b]{\linewidth}\raggedright
baseline \code{RUN_2026-10-09_06-18-39}
\end{minipage} & \begin{minipage}[b]{\linewidth}\raggedright
new \code{RUN_2026-10-11_08-39-41}
\end{minipage} \\
\midrule\noalign{}
\endhead
\bottomrule\noalign{}
\endlastfoot
LVS / routing DRC / Magic DRC & 0 / 0 / 8 984 & 0 / 0 / 8 984 \\
setup worst slack (violators), all 9 corners & +3.3286 ns max\_ss (0) &
+2.4825 ns max\_ss (0); nom\_tt +10.1818 \\
hold worst slack (violators), all 9 corners & +0.2822 ns min\_ff (0) &
+0.2874 ns min\_ff (0); nom\_tt +0.5395 \\
antenna nets / pins; diodes & 205 / 273; 10 918 & 187 / 265; 12 500 \\
max-slew pins max\_ss / max\_tt / nom\_tt (checker skipped) & 423 / 257
/ 95 & 370 / 230 / 200 \\
max-cap raw violations & 1 per tt/ss corner (6) & max\_ss 3, nom\_ss 1,
all other corners 0 \\
\code{CVT.MaxCapViolations} & PASS (6 waived / 0 unwaived) &
\textbf{FAIL: 0 waived / 4 unwaived / 1 stale waiver} \\
util / instances / power & 45.21 \% / 312 320 / 84.9 mW & 45.22 \% / 313
777 / 84.6 mW \\
\end{longtable}
}

}

\textbf{The waiver is stale.} \code{cpu-axi-rdata-i12-pin-cap} waived
nothing at any corner. Its premise, the anomalous 0.2464 pF Liberty pin
cap, belonged to the old macro. The new Liberty has
\code{axi_rdata_i[12]} at 0.0469 pF (nom\_tt) / 0.0496 pF (max\_ss); the
largest data input is now \code{axi_rdata_i[14]} at 0.0675 / 0.0691 pF
(input median 0.0056 pF). \textbf{The waiver was neither removed nor
changed} (the check still fails, and waiver changes are an owner
decision).

\textbf{Three unwaived nets, all flat SoC logic with no macro load, ss
corners only} (no waiver added):

{
\begin{longtable}[]{@{}llll@{}}
\toprule\noalign{}
net / driver & corner & cap & limit \\
\midrule\noalign{}
\endhead
\bottomrule\noalign{}
\endlastfoot
\code{_114232_/Y} (\code{_060214_}) & max\_ss & 0.3020 pF & 0.2701 pF
(+11.8 \%) \\
\code{_114232_/Y} (\code{_060214_}) & nom\_ss & 0.2862 pF & 0.2701 pF
(+6.0 \%) \\
\code{_111540_/X} (\code{_057530_}) & max\_ss & 0.3278 pF & 0.3250 pF
(+0.9 \%) \\
\code{_112095_/X} (\code{_058083_}) & max\_ss & 0.3345 pF & 0.3342 pF
(+0.1 \%) \\
\end{longtable}
}

Report: \code{62-cvt-maxcapviolations/max_cap_waivers.rpt} of the run.

\textbf{Owner decision 2026-10-11, "waive 2, fix 1" (applied in
\src{pnr/sky130/soc/constraints/max_cap_waivers.json}):}

\begin{itemize}
\tightlist
\item
  The stale \code{cpu-axi-rdata-i12-pin-cap} waiver is \textbf{removed}.
\item
  \code{_111540_/X} (0.3278 vs 0.3250 pF, +0.9 \%) and \code{_112095_/X}
  (0.3345 vs 0.3342 pF, +0.1 \%), both max\_ss only, are \textbf{waived
  by driver name}, each with a justification, restricted to the
  \code{max_ss_100C_1v60} corner, a ceiling 0.5 \% above the measured
  value, and \textbf{pinned to \code{RUN_2026-10-11_08-39-41}}.
  Synthesis-generated names go stale on any re-run, so these two waivers
  apply only to that run, show up as STALE on any other, and must then
  be re-evaluated, not renewed.
\item
  \code{_114232_/Y} (+11.8 \% at max\_ss, +6.0 \% at nom\_ss) stays
  \textbf{unwaived}. It must be closed at the next SoC run; the
  \code{j41m.4} scan-pin re-harden of the CPU macro forces one (recorded
  in bead \bead{e45j}).
\item
  To express a driver-name waiver the checker
  (\code{cvt_maxcap_waiver.py}) gained an optional \code{driver} key
  (exact instance, no globs) that is mutually exclusive with
  \code{load_pin} and \textbf{requires} \code{run_id}; the flow step and
  the CLI pass the run directory name. Load-pin waivers are unchanged. 7
  new unit tests (\src{tb/tests/test_cvt_maxcap_waiver.py}, 59 pass).
\end{itemize}

Offline check against the existing run (checker script only, no flow
re-run):
\texttt{python3\ pnr/\allowbreak{}sky130/\allowbreak{}soc/\allowbreak{}plugin/\allowbreak{}cvt\_\allowbreak{}maxcap\_\allowbreak{}waiver.\allowbreak{}py\ /\allowbreak{}nobackup/\allowbreak{}sky130\_\allowbreak{}soc\_\allowbreak{}runs/\allowbreak{}RUN\_\allowbreak{}2026-\allowbreak{}10-\allowbreak{}11\_\allowbreak{}08-\allowbreak{}39-\allowbreak{}41\ -\allowbreak{}-\allowbreak{}waivers\ pnr/\allowbreak{}sky130/\allowbreak{}soc/\allowbreak{}constraints/\allowbreak{}max\_\allowbreak{}cap\_\allowbreak{}waivers.\allowbreak{}json}
-\textgreater{}
\texttt{2\ waived,\allowbreak{}\ 2\ unwaived\ violation(s),\allowbreak{}\ 0\ stale\ waiver(s)};
the two unwaived are both \code{_114232_/Y} (max\_ss 0.3020 pF, nom\_ss
0.2862 pF vs 0.2701 pF);
\texttt{max-\allowbreak{}cap\ check:\allowbreak{}\ FAIL} (exit 1), as
intended until \code{_114232_} is closed.

Previous owner decision options, now resolved: waive by name (two nets),
accept (none), re-run (\code{_114232_} at the next run). The
setup/\allowbreak{}hold/\allowbreak{}LVS/\allowbreak{}DRC numbers do not
depend on that decision. The SoC closes setup and hold at all nine
corners (typical-corner caveat of
\src{docs/SKY130_REAL_DRC_LVS_EVALUATION.md} (\cref{tm:pd:sky130-real-drc-lvs}) §7 unchanged), whereas the
CPU \textbf{standalone} macro does not close ss reg-to-reg setup
(max\_ss -3.70 ns, 996 violators, icache tag-SRAM half-cycle paths at 75
MHz); the SoC runs at 40 MHz.

\end{record}

\subsection{2026-10-26 addendum: stable hand-instantiated driver waivers
(PR \#330, DFT Stage 2)}\label{cr:defect:sky130-slew-cap-e45j:2026-10-26-addendum-stable-hand-instantiated-driver-waivers-pr-330-dft-stage-2}

\begin{record}

A \code{driver} waiver can now set \texttt{"stable":\allowbreak{}\ true}
and omit \code{run_id} when the driver is a stable hand-instantiated
driver, e.g. the clock mux \code{u_cpu_pll_sub.u_ref_cm.u_mux/X}. The
loader refuses \code{stable} on a tool-generated name; ceiling, corners
and justification bind as before; a stable waiver on a driver that does
not violate is STALE. Synthesis-named drivers keep the mandatory
\code{run_id}. Waiver file state after the scan run
\code{RUN_2026-10-26_09-59-03}: mux\_clk (no run\_id, 0.45 pF),
din0{[}22{]} (0.37 pF, max\_ss + nom\_ss), din0{[}23{]} (0.36),
din0{[}26{]} (0.355); \code{_047147_/Y} stays unwaived and is tracked
here. \code{--strict-stale} on that run: 7 waived, 1 unwaived, 0 stale.
Details: \src{docs/design/DFT_ARCHITECTURE.md} (\cref{tm:ch:dft-architecture}) 14.13.

\end{record}

\subsection{\texorpdfstring{9. Wire-length 1500 adopted as default
(2026-10-28, bead
\code{e45j})}{9. Wire-length 1500 adopted as default (2026-10-28, bead )}}\label{cr:defect:sky130-slew-cap-e45j:9-wire-length-1500-adopted-as-default-2026-10-28-bead-e45j}

\begin{record}

User decision 2026-10-28:
\texttt{GRT\_\allowbreak{}DESIGN\_\allowbreak{}REPAIR\_\allowbreak{}MAX\_\allowbreak{}WIRE\_\allowbreak{}LENGTH\ =\allowbreak{}\ 1500}
is now in \src{pnr/sky130/soc/config_rccal.json}, the default Sky130 SoC
flow. Section 6.4 asked for one more arm on a different RTL vintage
before adopting it; this is that arm. \code{config_rccal_wl1500.json}
was the A/B arm of sections 6 and 9 and is \textbf{removed} (it was a
copy of \code{config_rccal.json} plus this one key, so it is redundant;
nothing else referenced it besides historical text in this file,
\code{DFT_ARCHITECTURE.md} (\cref{tm:ch:dft-architecture}) and \code{SKY130_CDC_MAXSS_ANALYSIS.md} (\cref{tm:pd:sky130-cdc-maxss}),
which are left as history).

\textbf{The arm.} \code{RUN_2026-10-28_05-14-30} against baseline
\code{RUN_2026-10-27_05-40-27} (\prref{333}, scan re-harden, CPU scan
macro, \code{CVTClassicScan}, \code{CVT_SYNTH_GATE=warn}, KLayout DRC
skipped,
\texttt{make\ -\allowbreak{}C\ pnr\ librelane-\allowbreak{}sky130-\allowbreak{}soc-\allowbreak{}noklayout\ SKY130\_\allowbreak{}SOC\_\allowbreak{}CONFIG=\allowbreak{}config\_\allowbreak{}rccal\_\allowbreak{}wl1500.\allowbreak{}json}).
Single variable, measured: the flattened \code{resolved.json} differs in
exactly one key (0 -\textgreater{} 1500), and \code{soc_top.nl.v} of
steps 04, 07 and 11 is sha256-identical to the baseline (worktree path
normalised). RSZ-0065 in step 38 confirms the knob reached the post-GRT
repair; step 29 (post-GPL, \code{DESIGN_REPAIR_MAX_WIRE_LENGTH}) is
unchanged at the automatic 3 205 um. The flow is bit-deterministic here
(section 6.1), so every difference is a response to the key.

{\footnotesize\setlength{\tabcolsep}{3pt}

{
\begin{longtable}[]{@{}
  >{\raggedright\arraybackslash}p{\dimexpr(\linewidth-4\tabcolsep)*323/1000\relax}
  >{\raggedright\arraybackslash}p{\dimexpr(\linewidth-4\tabcolsep)*328/1000\relax}
  >{\raggedright\arraybackslash}p{\dimexpr(\linewidth-4\tabcolsep)*328/1000\relax}@{}}
\toprule\noalign{}
\begin{minipage}[b]{\linewidth}\raggedright
\end{minipage} & \begin{minipage}[b]{\linewidth}\raggedright
wl1500 \code{RUN_2026-10-28_05-14-30}
\end{minipage} & \begin{minipage}[b]{\linewidth}\raggedright
baseline \code{RUN_2026-10-27_05-40-27}
\end{minipage} \\
\midrule\noalign{}
\endhead
\bottomrule\noalign{}
\endlastfoot
max-slew violating pins, max\_ss & \textbf{129} & 1067 \\
nom\_ss / min\_ss & 57 / 33 & 709 / 64 \\
max\_tt / nom\_tt & 1 / 0 & 695 / 145 \\
max\_ff / nom\_ff & 0 / 0 & 596 / 48 \\
max\_ss slew, diode pins & 89 & 930 \\
max-cap violators in the flow STA, max\_ss & 5 & 6 \\
max-cap with the waiver file, \code{--strict-stale} & 8 waived, 1
unwaived, 2 stale (old file); \textbf{9 waived, 0 unwaived, 0 stale}
(file of this section) & 10 unwaived, 3 stale (old file) \\
setup, max\_ss (flow STA) & 5 violators, WNS -0.055, TNS -0.226 ns & 6
violators, WNS -0.153, TNS -0.603 ns \\
\code{core_clk} setup nom\_tt / nom\_ss / max\_ss (replica STA, ns) &
+5.600 / +0.220 / -0.055 & +5.588 / +0.132 / -0.153 \\
hold violators, 9 corners & 0 & 0 \\
\code{core_clk} hold worst max\_ss / nom\_ff / max\_ff (ns) & +0.336 /
+0.245 / +0.247 & +0.231 / +0.280 / +0.268 \\
done -\textgreater{} \code{dft_scan_en} -\textgreater{} SE setup,
max\_ss (ns) & +0.740 & +0.671 \\
Netgen LVS / routing DRC & 0 / 0 & 0 / 0 \\
Magic DRC & 10 962 (nwell.4 10 930 + met4.4a 32) & 10 959 (10 927 +
32) \\
antenna nets / pins (diodes) & 292 / 375 (18 683) & 290 / 382 (20
500) \\
utilisation / power & 46.63 \% / 119.03 mW & 46.64 \% / 119.10 mW \\
instances / timing-repair buffers & 331 309 / 28 879 & 332 973 / 28
726 \\
step-38 repair (buffers / nets / resized) & 240 / 341 / 184 & 45 / 174 /
63 \\
post-route CDC budget check & PASS at all 9 corners & PASS (after the
check-SDC fix of \code{DFT_ARCHITECTURE.md} (\cref{tm:ch:dft-architecture}) 14.14.1) \\
\end{longtable}
}

}

Slew: the 31 diode-heavy nets driven by \code{wire} repair clkbuf\_16
(up to 59 diodes each) are gone, which is where 930 -\textgreater{} 89
diode pins comes from. What is left at max\_ss is 10 nets: 25 pins on
macro/\allowbreak{}port nets, two \code{a22oi_4} nets, three hold
dlygates and a \code{mux2}. The max-slew checker stays skipped in the
flow, so these are measured from the STA reports, not gated. Setup: the
same NPU weight-memory SRAM -\textgreater{} \code{npu_weight_mem} flop
paths (bead \bead{7dhp}, accepted) at max\_ss only; the sixth endpoint
(\code{_169106_}) is now positive. Hold stays positive everywhere but is
lower at the nom\_\allowbreak{}ff/\allowbreak{}max\_\allowbreak{}ff
corners (-0.035 / -0.021 ns). Both runs ended \code{make} with deferred
errors; this run\textquotesingle s were Magic DRC (accepted, bead
\bead{45a}) and one max-cap net, no early exit (all 64 steps ran, LVS
included).

\textbf{Waiver changes with it (user decision 2026-10-28),
\src{pnr/sky130/soc/constraints/max_cap_waivers.json}.}

\begin{itemize}
\tightlist
\item
  \code{hold-dlygate-scd-link} class ceiling 0.115 -\textgreater{}
  \textbf{0.120 pF}: \code{hold18202/X} measured 0.1151 pF at max\_ss
  (0.1053 nom\_ss), 0.1 \% over the old ceiling, same class and cause
  (diode pin capacitance on a scan-link hold buffer); the other class
  members measured 0.1077, 0.1038 and 0.1031 pF. 0.120 is 4.2 \% above
  the worst measured value.
\item
  \code{gen-a21oi-107570-run-2026-10-27} and
  \code{gen-or4-076179-run-2026-10-27} (run-pinned to the baseline)
  \textbf{deleted}: those two nets no longer violate, so they waived
  nothing (STALE under \code{--strict-stale}).
\item
  \code{u-ref-cm-mux-clk-cpu-pll-sub} is unchanged: \code{mux_clk} is
  0.442 / 0.463 / 0.468 pF
  (min/\allowbreak{}nom/\allowbreak{}max\_\allowbreak{}ss) against its
  0.55 pF ceiling (baseline 0.475 / 0.508 / 0.519).
\item
  Proof:
  \texttt{make\ -\allowbreak{}C\ pnr\ sky130-\allowbreak{}soc-\allowbreak{}maxcap-\allowbreak{}check\ SKY130\_\allowbreak{}SOC\_\allowbreak{}MAXCAP\_\allowbreak{}RUN=\allowbreak{}/\allowbreak{}nobackup/\allowbreak{}sky130\_\allowbreak{}soc\_\allowbreak{}runs/\allowbreak{}RUN\_\allowbreak{}2026-\allowbreak{}10-\allowbreak{}28\_\allowbreak{}05-\allowbreak{}14-\allowbreak{}30\ STRICT\_\allowbreak{}STALE=\allowbreak{}1}
  -\textgreater{} \textbf{9 waived, 0 unwaived, 0 stale: PASS}.
  \src{tb/tests/test_cvt_maxcap_waiver.py} 100 pass (its pin on the
  committed file now expects 0.12 and the two run-pinned entries
  absent). The adopted default has not itself been run end to end; this
  run used the identical config through a differently named file, so the
  next full run is the first with \code{config_rccal.json} directly.
\end{itemize}

\textbf{Scan-chain check needs \code{--en-net}.}
\src{tools/dft/check_scan_chains.py} counts the
resizer\textquotesingle s clone of the enable gate
(\texttt{and2\_\allowbreak{}2\ clone13} next to
\texttt{and2\_\allowbreak{}4\ \_\allowbreak{}090536\_\allowbreak{}},
same function and input nets) as the same root only when told the root
net: pass
\textbf{\texttt{-\allowbreak{}-\allowbreak{}en-\allowbreak{}net\ dft\_\allowbreak{}scan\_\allowbreak{}en}}.
Without it the check fails with "SCE pins trace to 2 different roots"
although the netlist is correct (measured both ways on this run). With
it:
\texttt{check\_\allowbreak{}scan\_\allowbreak{}chains.\allowbreak{}py\ <49-\allowbreak{}openroad-\allowbreak{}fillinsertion/\allowbreak{}soc\_\allowbreak{}top.\allowbreak{}nl.\allowbreak{}v>\ -\allowbreak{}-\allowbreak{}en-\allowbreak{}net\ dft\_\allowbreak{}scan\_\allowbreak{}en\ -\allowbreak{}-\allowbreak{}out-\allowbreak{}gate-\allowbreak{}net\ crypto\_\allowbreak{}wipe\_\allowbreak{}done\ -\allowbreak{}-\allowbreak{}expect-\allowbreak{}chains\ 6\ -\allowbreak{}-\allowbreak{}expect-\allowbreak{}unscanned\ 922\ -\allowbreak{}-\allowbreak{}expect-\allowbreak{}length\ 3565,\allowbreak{}3565,\allowbreak{}3565,\allowbreak{}3565,\allowbreak{}3562,\allowbreak{}981\ -\allowbreak{}-\allowbreak{}explain-\allowbreak{}qrx\ 'u\_\allowbreak{}crypto\textbackslash{}.\allowbreak{}(g\_\allowbreak{}aes\textbackslash{}.\allowbreak{}u\_\allowbreak{}aes\textbackslash{}.\allowbreak{}rk\_\allowbreak{}q|key\_\allowbreak{}q)'\ -\allowbreak{}-\allowbreak{}explain-\allowbreak{}clk\ tck\_\allowbreak{}i,\allowbreak{}scan\_\allowbreak{}clk\_\allowbreak{}i}
-\textgreater{} exit 0, 18 803 scan flops, 922 unscanned (664
\code{tck_i}, 2 \code{scan_clk_i}, 256 crypto keys), SCE roots
\texttt{\{dft\_\allowbreak{}scan\_\allowbreak{}en:\allowbreak{}\ 18803\}}.

Limits: one arm on this vintage (plus the section 6 arm on an older one,
both in the same direction); KLayout DRC skipped (bead \bead{00ef}); the
stock max-slew checker still skipped; hold margin at the ff corners is a
few tens of ps lower and was not investigated further.

\end{record}

% Converted from docs/SYNLIG_STRUCT_SELECT_AINF.md by tools/docs/md2tex.py; edit this file, not the original.
\section{\texorpdfstring{Synlig struct-member part-select defect: rule,
site table, equivalence (bead
\code{ainf})}{Synlig struct-member part-select defect: rule, site table, equivalence (bead )}}\label{cr:defect:synlig-struct-select-ainf}

\begin{record}

Date 2026-10-10. Scope: does the Synlig frontend silently mis-elaborate
the struct-member part-selects that \prref{259} (\code{dud4}) did not
rewrite? Everything below is MEASURED unless marked INFERRED. Tools:
yosys 0.46 + Synlig plugin
(\code{/nix/store/7vsbsrhbj6kx5lfdc8qs9wlccxfb8f7h-yosys-synlig-sv},
Surelog 1.82) vs sv2v 0.0.13.1 + yosys; comparison = miter +
\texttt{sat\ -\allowbreak{}verify\ -\allowbreak{}prove-\allowbreak{}asserts\ -\allowbreak{}seq\ 4\ -\allowbreak{}set-\allowbreak{}init-\allowbreak{}zero}
(flow from \prref{264}\textquotesingle s
\src{tools/verif/synth_gate/equiv_probe.sh}, re-implemented in
\src{tools/verif/struct_select/run_probes.sh}; every job under
\texttt{systemd-\allowbreak{}run\ .\allowbreak{}.\allowbreak{}.\allowbreak{}\ MemoryMax=\allowbreak{}4G}).

\end{record}

\subsection{\texorpdfstring{Step 1: the defect rule (probe evidence:
\code{tools/verif/struct_select/results/round1..7.txt}, 370
probes)}{Step 1: the defect rule (probe evidence: , 370 probes)}}\label{cr:defect:synlig-struct-select-ainf:step-1-the-defect-rule-probe-evidence-toolsverifstruct-selectresultsround17txt-370-probes}

\begin{record}

Two independent Synlig defects. Neither is triggered by plain reads of a
member of an ordinary struct.

\end{record}

\subsubsection{\texorpdfstring{D1 (read side): a
packed/\allowbreak{}unpacked struct with EXACTLY TWO members declared
with a packed
range}{D1 (read side): a packed/unpacked struct with EXACTLY TWO members declared with a packed range}}\label{cr:defect:synlig-struct-select-ainf:d1-read-side-a-packedunpacked-struct-with-exactly-two-members-declared-with-a-packed-range}

\begin{record}

"Ranged member" = declared \texttt{logic\ [n:\allowbreak{}0]\ m;} (a
\code{[0:0]} range counts; a scalar \texttt{logic\ m;}, a struct-typed
member and an enum-typed member do NOT count). If a struct type has
exactly two ranged members, then for ANY member \code{m} of it

\begin{lstlisting}[language={}]
s.m[hi:lo]   (or s.m[i], constant or runtime)   is evaluated by Synlig as   s[ hi*W + (hi-lo) : hi*W ]
\end{lstlisting}

where \code{W} is the width of the LAST-declared (LSB-most) ranged
member. \code{lo} and the member\textquotesingle s own offset are
ignored; the select behaves like "element \code{hi} of an array of W-bit
elements", truncated to the select width.

\begin{itemize}
\tightlist
\item
  \code{if_id_reg_t} =
  \texttt{\{pc[31:\allowbreak{}0],\allowbreak{}\ instruction[31:\allowbreak{}0],\allowbreak{}\ valid\}}
  has two ranged members (valid is a scalar) and W = 32:
  \code{instruction[19:15]} -\textgreater{}
  \texttt{[19*32+\allowbreak{}4\ :\allowbreak{}\ 19*32]\ =\allowbreak{}\ [612:\allowbreak{}608]}
  (the logged \code{[639:608]} is the whole 32-bit element, which is
  then truncated); \code{[24:20]} -\textgreater{} \code{[768+4:768]}.
  This reproduces the dud4 numbers exactly.
\item
  \texttt{n\_\allowbreak{}t\ =\allowbreak{}\ \{pad[7:\allowbreak{}0],\allowbreak{}\ in(struct),\allowbreak{}\ tail[4:\allowbreak{}0]\}}
  (two ranged members, W = 5): \code{n.tail[3:2]} -\textgreater{}
  \texttt{[3*5+\allowbreak{}1:\allowbreak{}3*5]\ =\allowbreak{}\ [16:\allowbreak{}15]},
  \code{n.pad[7:4]} -\textgreater{}
  \texttt{[7*5+\allowbreak{}3:\allowbreak{}7*5]\ =\allowbreak{}\ [38:\allowbreak{}35]}
  (both measured, in bounds, no warning, wrong).
\item
  Fits all 27 failing two-ranged-member probes of round 3 (random widths
  1..32) without exception.
\item
  Counter-evidence that this is not "number of members": a 3-member
  struct with one scalar (2 ranged) fails, a 3-member struct with
  \texttt{logic\ [0:\allowbreak{}0]} instead of the scalar (3 ranged) is
  correct (\code{s01} vs \code{s02}, round 5); 1, 3, 4, 5 ranged members
  are correct in all 120+ probes (rounds 2, 3, 6, 7).
\end{itemize}

Outcome classes (a) correct:
\texttt{hi\ =\allowbreak{}=\allowbreak{}\ 0} and the member is the LSB
member and \texttt{lo\ =\allowbreak{}=\allowbreak{}\ 0}; (b)
\texttt{hi*W\ +\allowbreak{}\ width\ >\ struct\ width}: warning
\texttt{Range\ select\ [.\allowbreak{}.\allowbreak{}]\ out\ of\ bounds\ on\ signal\ .\allowbreak{}.\allowbreak{}.\allowbreak{}}
and \code{x} (this is dud4); (c) in bounds but wrong, NO warning, NO
\code{x}: possible whenever
\texttt{hi*W\ +\allowbreak{}\ (hi-\allowbreak{}lo)\ <\ total\ width};
measured instances: round 7 \code{X2_rhs_hi1} (\code{s.b[1:0]} read from
\code{s[33:32]} instead of \code{s[1:0]}), \code{X2_rhs_bit},
\code{X2_rhs_concat}, \code{X2_rhs_reg}, \code{X2_rhs_local},
\code{i03}, \code{i04}, round 5 \code{s*}. (c) IS possible, and for
small \code{hi} it is the normal case, not the exception. Split by
\src{tools/verif/struct_select/classify.py} (round 3, 27 failing of 150:
17 warn+x, 2 warn+wrong, \textbf{8 fully silent}; round 7: 21 fully
silent, 17 warn+wrong, 5 warn+x). "warn+wrong" = the
\texttt{Range\ select\ .\allowbreak{}.\allowbreak{}.\allowbreak{}\ out\ of\ bounds}
warning IS printed (the 32-bit "element" overruns the struct) but the
surviving bits are wrong data, not \code{x}; the \prref{264} warning
scanner sees those. "Fully silent" = the element fits inside the struct:
no warning, no \code{x}, wrong bits (e.g. \code{s.m[0]} of a non-LSB
member of a two-ranged struct, or \code{s.b[1:0]} with W large); only a
miter / simulation sees them. D2 is always fully silent.

Where D1 shows: port-typed struct, local wire/reg struct, concatenation
operands, \code{case} selectors, bit selects, runtime bit select
\code{s.m[i]}, port-connection expressions, unpacked structs with two
ranged members. Where it does NOT show:
\texttt{s.\allowbreak{}m[b\ +\allowbreak{}:\allowbreak{}\ n]} /
\texttt{s.\allowbreak{}m[b\ -\allowbreak{}:\allowbreak{}\ n]} (indexed
part-select, correct), selects of a full-width wire copy
(\texttt{w\ =\allowbreak{}\ s.\allowbreak{}m;\ w[7:\allowbreak{}4]}), an
array-of-struct element base (\code{arr[1].m[1:0]} was correct in the
one probe, \code{X2_rhs_arr}; not proven in general). It does not depend
on the base signal being a port, local, register or array element beyond
that.

\end{record}

\subsubsection{D2 (write side): part-select of a struct member on the
LEFT-HAND side}\label{cr:defect:synlig-struct-select-ainf:d2-write-side-part-select-of-a-struct-member-on-the-left-hand-side}

\begin{record}

\texttt{y.\allowbreak{}m[hi:\allowbreak{}lo]\ =\allowbreak{}\ v} with
\texttt{hi\ !=\allowbreak{}\ lo} is elaborated as the WHOLE-member
assignment \texttt{y.\allowbreak{}m\ =\allowbreak{}\ v} (zero-extended);
the bits of \code{m} outside the select are lost and the select position
is ignored. No warning, no \code{x}. Independent of the number of ranged
members (measured R = 2, 3, 4: \code{lhs_copy}, \code{lhs_zero},
\code{lhs_ff}, \code{lhs_assign} all NOT equivalent). The single-bit LHS
\texttt{y.\allowbreak{}m[i]\ =\allowbreak{}\ v} is correct for R
\textgreater= 3 and wrong for R = 2 (\code{X3_lhs_bit},
\code{X4_lhs_bit} equivalent; \code{X2_lhs_bit} not).

\end{record}

\subsection{Step 1 reading for this repo}\label{cr:defect:synlig-struct-select-ainf:step-1-reading-for-this-repo}

\begin{record}

The rule is a function of the struct TYPE (count of ranged members), not
of the member position, so every site can be predicted from its struct
definition (step 2).

\end{record}

\subsection{\texorpdfstring{Step 2a: mechanical site table
(\src{tools/verif/struct_select/sitetable.py}, output
\code{results/sitetable.txt})}{Step 2a: mechanical site table (, output )}}\label{cr:defect:synlig-struct-select-ainf:step-2a-mechanical-site-table-toolsverifstruct-selectsitetablepy-output-resultssitetabletxt}

\begin{record}

Struct types in \src{rtl/} (R = members declared with a packed range,
i.e. the D1 trigger when R == 2):

{\footnotesize\setlength{\tabcolsep}{3pt}

{
\begin{longtable}[]{@{}
  >{\raggedright\arraybackslash}p{\dimexpr(\linewidth-6\tabcolsep)*401/1000\relax}
  >{\centering\arraybackslash}p{\dimexpr(\linewidth-6\tabcolsep)*48/1000\relax}
  >{\raggedright\arraybackslash}p{\dimexpr(\linewidth-6\tabcolsep)*129/1000\relax}
  >{\raggedright\arraybackslash}p{\dimexpr(\linewidth-6\tabcolsep)*401/1000\relax}@{}}
\toprule\noalign{}
\begin{minipage}[b]{\linewidth}\raggedright
type
\end{minipage} & \begin{minipage}[b]{\linewidth}\centering
R
\end{minipage} & \begin{minipage}[b]{\linewidth}\raggedright
where
\end{minipage} & \begin{minipage}[b]{\linewidth}\raggedright
note
\end{minipage} \\
\midrule\noalign{}
\endhead
\bottomrule\noalign{}
\endlastfoot
\code{if_id_reg_t} & \textbf{2} &
rv32i\_\allowbreak{}pipeline\_\allowbreak{}pkg & D1 HOT (\code{pc},
\code{instruction} ranged; \code{valid} scalar). The dud4 type. \\
\code{warp_state_t} & \textbf{2} & gpu\_pkg & D1 HOT, but no member
select exists on it \\
\code{div_entry_t} & \textbf{2} & gpu\_pkg & D1 HOT, but no member
select exists on it \\
\code{desc_t} & 3 & dma\_\allowbreak{}engine & one member-removal from
R=2; no member select \\
\code{mem_wb_reg_t} 7, \code{ex_mem_reg_t} 12, \code{ex1_ex2_reg_t} 12,
\code{id_ex_reg_t} 14, \code{ex1a_ex1b_t} 19 & \textgreater= 7 &
pipeline pkg & far from 2 \\
\code{if_id_t} 4, \code{id_ex_t} 9, \code{ex_wb_t} 6
(gpu\_\allowbreak{}compute\_\allowbreak{}unit), \code{kernel_desc_t} 8,
\code{gpu_[rib]_instr_t} 5-6 & \textgreater= 4 & gpu & \\
\end{longtable}
}

}

READ selects \code{x.m[..]} in the whole \src{rtl/} tree: 23 (24 incl.
the two on one gpu\_top line), all on types with R \textgreater= 7, so
\textbf{all predicted correct}:
rv32i\_\allowbreak{}pipeline\_\allowbreak{}ex:109
(\code{id_ex_reg_t.csr_op[2]}),
rv32i\_\allowbreak{}pipeline\_\allowbreak{}mem:103,104,111,112
(\code{ex_mem_reg_t.alu_result[0]}, \code{[1:0]}),
rv32i\_\allowbreak{}pipeline\_\allowbreak{}wb:75
(\code{mem_wb_reg_t.trap_cause[3:0]}),
rv32i\_\allowbreak{}pipeline\_\allowbreak{}ex1c:122-150,197
(\code{ex1a_ex1b_t.alu_result[1:0]/[1]/[31:1]},
\code{.fwd_store[7:0]/[15:0]}), gpu\_\allowbreak{}top:253
(\code{kernel_desc_t.block_x[9:3]}, \code{[2:0]}; the line was :250 in
the dud4 sweep),
gpu\_\allowbreak{}compute\_\allowbreak{}unit:377,378,386,388
(\code{id_ex_t.rs1_data[l]}, \code{.rs2_data[l]}, \code{.imm[11]}). None
of these is on an R == 2 type, so D1 does not apply. \textbf{No select
at all exists on any of the three R == 2 types}
(\code{if_id_reg_t}\textquotesingle s two were removed by \prref{259}). Not
covered by probes: \code{id_ex_t.rs1_data[l]}, a member that is itself a
2-D packed array \code{[N_LANES-1:0][31:0]} selected by a loop variable
-\/- predicted correct by the type rule but verified separately in step
2b.

WRITE (LHS) selects
\texttt{y.\allowbreak{}m[hi:\allowbreak{}lo]\ =\allowbreak{}\ .\allowbreak{}.\allowbreak{}}
/
\texttt{y.\allowbreak{}m[i]\ =\allowbreak{}\ .\allowbreak{}.\allowbreak{}}
(D2): \textbf{none} in \src{rtl/} (regex over single-line statements,
plain, \code{assign}, and concatenation-LHS forms; 0 hits). D2 is
therefore not triggered anywhere today.

\end{record}

\subsection{\texorpdfstring{Step 2b: Synlig-vs-sv2v equivalence on the
real modules (\code{modeq.sh}, \code{run_modeq.sh},
\code{mkcone.py})}{Step 2b: Synlig-vs-sv2v equivalence on the real modules (, , )}}\label{cr:defect:synlig-struct-select-ainf:step-2b-synlig-vs-sv2v-equivalence-on-the-real-modules-modeqsh-run-modeqsh-mkconepy}

\begin{record}

Method: elaborate the module with Synlig and with sv2v+yosys,
\code{flatten}, expose flops
(\texttt{expose\ -\allowbreak{}evert-\allowbreak{}dff}, so the check is
combinational with matched state), miter +
\texttt{sat\ -\allowbreak{}verify\ -\allowbreak{}prove-\allowbreak{}asserts}.
Memory cap
\texttt{systemd-\allowbreak{}run\ -\allowbreak{}p\ MemoryMax=\allowbreak{}4G\ -\allowbreak{}p\ MemorySwapMax=\allowbreak{}0},
timeout 900 s. Negative control: sv2v arm fed a copy of
\code{rv32i_pipeline_wb.sv} with \code{trap_cause[3:0]} changed to
\code{[4:1]} -\textgreater{} NOT-EQUIV (model found), as required. SAT
sizes are non-trivial (e.g. 10 884 variables / 29 627 clauses for
\code{rv32i_pipeline_ex}).

{\small\setlength{\tabcolsep}{4pt}

{
\begin{longtable}[]{@{}
  >{\raggedright\arraybackslash}p{\dimexpr(\linewidth-6\tabcolsep)*377/1000\relax}
  >{\raggedright\arraybackslash}p{\dimexpr(\linewidth-6\tabcolsep)*256/1000\relax}
  >{\raggedright\arraybackslash}p{\dimexpr(\linewidth-6\tabcolsep)*128/1000\relax}
  >{\raggedright\arraybackslash}p{\dimexpr(\linewidth-6\tabcolsep)*219/1000\relax}@{}}
\toprule\noalign{}
\begin{minipage}[b]{\linewidth}\raggedright
module / cone
\end{minipage} & \begin{minipage}[b]{\linewidth}\raggedright
sites
\end{minipage} & \begin{minipage}[b]{\linewidth}\raggedright
prediction
\end{minipage} & \begin{minipage}[b]{\linewidth}\raggedright
result
\end{minipage} \\
\midrule\noalign{}
\endhead
\bottomrule\noalign{}
\endlastfoot
\code{rv32i_pipeline_wb} & wb:75 & correct & \textbf{EQUIV} \\
\code{rv32i_pipeline_ex1c} & ex1c:122-150,197 (byte/\allowbreak{}half
store lanes) & correct & \textbf{EQUIV} \\
\code{rv32i_pipeline_mem} & mem:103,104,111,112 (misalign) & correct &
\textbf{EQUIV} \\
\code{rv32i_pipeline_ex} (with alu, branch\_\allowbreak{}comp) & ex:109
& correct & \textbf{EQUIV} \\
\code{gpu_top} cone (\code{n_warps_raw}, text extracted verbatim) &
gpu\_\allowbreak{}top:253 & correct & \textbf{EQUIV} \\
\code{gpu_compute_unit} cone (always\_\allowbreak{}comb at :370-390
extracted verbatim, \code{id_ex_t} from the source) & 377,378,386,388 &
\textbf{correct (type rule)} & \textbf{NOT EQUIVALENT -\/- prediction
wrong} \\
\end{longtable}
}

}

The three D1/D2 sites classes behave as predicted. The GPU result is a
\textbf{third, independent defect (D3)} that the type rule does not
cover:

\end{record}

\subsubsection{D3: element select of a struct member that is a
multi-dimensional packed array}\label{cr:defect:synlig-struct-select-ainf:d3-element-select-of-a-struct-member-that-is-a-multi-dimensional-packed-array}

\begin{record}

\code{id_ex_t} has
\texttt{logic\ [N\_\allowbreak{}LANES-\allowbreak{}1:\allowbreak{}0][31:\allowbreak{}0]\ rs1\_\allowbreak{}data,\allowbreak{}\ rs2\_\allowbreak{}data;}.
\code{id_ex_q.rs1_data[l]} (32-bit element \code{l}) is elaborated by
Synlig as the single BIT \code{rs1_data[l]} (zero-extended), at the
correct member base. Evidence in the cone netlist: Synlig
\texttt{mem\_\allowbreak{}addr\_\allowbreak{}o[31:\allowbreak{}0]\ =\allowbreak{}\ id\_\allowbreak{}ex\_\allowbreak{}q[320]\ +\allowbreak{}\ sext(imm)}
versus sv2v
\texttt{id\_\allowbreak{}ex\_\allowbreak{}q[351:\allowbreak{}320]\ +\allowbreak{}\ sext(imm)};
lane 1 reads bit 321 instead of \code{[383:352]}, etc. Probes (round 8,
\code{results/round8_d3.txt}; independent of R,
constant/\allowbreak{}runtime/\allowbreak{}loop index,
port/\allowbreak{}local base, always fully silent: \code{oob=0}):
\code{s.r[2]} -\textgreater{} \code{s[10]} instead of \code{s[31:24]},
\code{s.r[i]}, \code{s.r[l]} in a loop, a one-ranged-member struct too.
Correct forms: \texttt{w\ =\allowbreak{}\ s.\allowbreak{}r;\ w[l]}, a
plain (non-struct) 2-D vector \code{r[l]}, and a select of bits INSIDE
an element (\code{s.r[2][5:2]}, equivalent).

Consequence in \code{gpu_compute_unit}: the memory address
(\code{mem_addr_o[l]}, \code{shmem_addr_o[l]}) and store data
(\code{mem_wdata_o[l]}, \code{shmem_wdata_o[l]}) of
VLD/\allowbreak{}VST/\allowbreak{}VLDS/\allowbreak{}VSTS are computed
from ONE bit of the rs1/rs2 vector register instead of the 32-bit value.
Not proven on whole-module level (flop names differ between the
frontends after aliasing, which defeats the name-based flop exposure;
the cone is verbatim source text with the same typedef). No CPU module
contains this idiom (\code{sitetable.txt}: the only
multi-dimensional-array-member element selects in \src{rtl/} are these
four GPU lines).

\end{record}

\subsection{Step 4: verdict, recommendation, limits}\label{cr:defect:synlig-struct-select-ainf:step-4-verdict-recommendation-limits}

\begin{record}

Verdict: see the dated paragraph in \cref{cr:defect:asap7-cpu-synlig-ma7}.
CPU: remaining sites correct under Synlig (type rule + EQUIV of
ex/\allowbreak{}mem/\allowbreak{}wb/\allowbreak{}ex1c; step 3 gate-level
load/\allowbreak{}store programs not needed, no CPU module was predicted
or measured wrong). GPU: second live miscompile (D3) in
\code{gpu_compute_unit}; \code{gpu_top} cone is fine.

Recommendation (no \src{rtl/} edits made here): (1) rewrite the
remaining sites through named signals and broaden
\cref{um:guide:coding-guidelines} to "no
part/\allowbreak{}bit/\allowbreak{}element select of a struct member
anywhere"; (2) GPU: wire-first
(\texttt{logic\ [N\_\allowbreak{}LANES-\allowbreak{}1:\allowbreak{}0][31:\allowbreak{}0]\ rs1\_\allowbreak{}w\ =\allowbreak{}\ id\_\allowbreak{}ex\_\allowbreak{}q.\allowbreak{}rs1\_\allowbreak{}data;})
or build the GPU with sv2v only (the \code{ma7} config change already
did) and treat the committed Synlig GPU views as suspect; (3) CI:
\src{tools/verif/check_struct_member_select.py} (source-idiom tripwire,
baseline \src{tools/verif/struct_member_select_baseline.txt} =
today\textquotesingle s 23 select sites + 3 R == 2 types, unit tests
\src{tb/tests/test_check_struct_member_select.py}); it catches D1, D2
and D3 because it flags every member select, but its baseline key has no
line numbers, so a second select of the same \code{base.member} in the
same file is not flagged. A stronger check is a per-module
Synlig-vs-sv2v miter (\code{modeq.sh}) in the PD flow.

Not verified: whole-module equivalence of
\code{gpu_compute_unit}/\code{gpu_top}/\code{rv32i_core} (flop-name
aliasing between the frontends defeats the name-based flop exposure;
cones were cut verbatim instead); the GPU netlist behaviour in
simulation; the committed GPU macro views; any select form hidden from
the single-line regex (multi-line expressions, macros); Sky130 netlists;
step 3 (gate-level ld/st programs) was skipped as agreed.

\end{record}

% Converted from docs/SKY130_CPU_SS_SETUP_R51P.md by tools/docs/md2tex.py; edit this file, not the original.
\section{Sky130 CPU macro: ss reg-to-reg setup failure (bead r51p) -
path analysis}\label{cr:defect:sky130-cpu-ss-setup-r51p}

\begin{record}

Status: cache-side RTL fix landed on \code{fix/icache-ss-setup-r51p} (\prref{318}, stacked on \#313). User decision 2026-10-19: no cycle-count
change and no pipeline restructuring; the remaining ss residual stays
open in r51p and is measured at the batched CPU re-harden. Source data:
\code{pnr/sky130/cpu/runs/RUN_2026-10-11_06-00-12/57-openroad-stapostpnr/max_ss_100C_1v60/max.rpt}
(top 1000 setup paths, 75 MHz = 13.333 ns, uncertainty 0.5 ns).
Measured, not inferred, unless marked.

\end{record}

\subsection{1. The premise in the bead is not what the reports show}\label{cr:defect:sky130-cpu-ss-setup-r51p:1-the-premise-in-the-bead-is-not-what-the-reports-show}

\begin{record}

The bead and the user decision describe the worst paths as "into the
icache tag SRAM\textquotesingle s negedge-sampled inputs (half-cycle
paths)". The reports say otherwise.

\begin{itemize}
\tightlist
\item
  \textbf{The SRAM input arcs are full-cycle.} In
  \src{pnr/sky130/cpu/pdk_lib_patched/sky130_sram_1kbyte_1rw1r_32x256_8_TT_1p8V_25C.lib}
  every input pin (\code{csb0}, \code{web0}, \code{addr0}, \code{din0},
  \code{wmask0}) has
  \texttt{timing\_\allowbreak{}type\ :\allowbreak{}\ setup\_\allowbreak{}rising}
  against \code{clk0}. The worst path\textquotesingle s required time is
  \texttt{13.\allowbreak{}333\ (rise\ edge)\ +\allowbreak{}\ clk0\ arrival\ -\allowbreak{}\ 0.\allowbreak{}5\ uncertainty\ -\allowbreak{}\ 0.\allowbreak{}103\ setup}:
  a capture at the next rising edge, one full period. STA prints the
  endpoint as "falling edge-triggered flip-flop" only because the macro
  also has a \code{falling_edge} arc on \code{dout0}.
\item
  \textbf{The half-cycle paths are on the SRAM outputs} (\code{dout0}
  launches on the falling edge). They are already cut by the
  \code{CS_SRAM_LATCH} register stage (\code{tag_dout_r} /
  \code{data_dout_r}). None of the 1000 reported worst paths starts at
  an SRAM.
\item
  \textbf{The SRAM endpoints are 5 to 50 of 1000, not the population.}
  947 of the 1000 reported endpoints are ordinary flops. The worst SRAM
  endpoint (\code{u_tag_sram/csb0}, -3.704 ns) is 0.02 ns worse than the
  worst flop endpoint (\code{pc_if_from_core}, -3.684 ns). Retiming only
  the SRAM input pins cannot move WNS.
\end{itemize}

\end{record}

\subsection{\texorpdfstring{2. What the failing cone is (measured, from
the run\textquotesingle s own netlist
\code{54-openroad-fillinsertion/rv32i_cpu_top.nl.v})}{2. What the failing cone is (measured, from the run\textquotesingle s own netlist )}}\label{cr:defect:sky130-cpu-ss-setup-r51p:2-what-the-failing-cone-is-measured-from-the-runs-own-netlist-54-openroad-fillinsertionrv32i-cpu-topnlv}

\begin{record}

Start points of the 1000 paths (flop Q-net name):
\code{u_core.u_dcache.state_q} 407 (worst -3.704),
\code{u_core.ex_mem_reg} 329 (-3.363), \code{u_core.ex1b_ex2_reg_q} 116
(-2.907), plus a few others. End points: \code{u_icache.refill_buf_q}
128, \code{u_dcache.wb_buf_q} 82,
\code{u_csr.mcycle_q}/\code{minstret_q} 74, \code{if_id_reg} 33,
\code{ic_addr_q} 32, \code{pc_if_from_core} 30, debug-register write
data \code{dbg_*} \textasciitilde230, SRAM csb0 5.

Three combinational cells (\code{_22517_} o21ai, \code{_22618_} a21oi,
\code{wire1348}) lie on 58 \% of all 1000 paths. Their fan-in cone
(netlist traversal) is exactly: 8 bits of \code{ex_mem_reg}, 5 bits of
\code{u_dcache.state_q}, \code{u_dcache.hit_bank_q},
\code{mem_trap_entry_r}, \code{ex_pc_redirect_r},
\code{ex1c_ex1b_reg_q}. That is the \textbf{global stall / redirect
network}:

\begin{lstlisting}[language={}]
dcache.state_q, ex_mem_reg  ->  dc_stall_o -> mem_cache_stall (rv32i_pipeline_mem.sv:151)
                            ->  hazard-unit priority chain (rv32i_hazard_unit.sv:362-427: mem_cache_stall beats ex_pc_redirect)
                            ->  pc_redirect / jal_redirect -> ic_valid_o (rv32i_pipeline_if.sv:116)
                            ->  icache: addr_mismatch / hit / ic_stall_o / SRAM csb0 / refill_buf_q enable
                            ->  if_cache_stall -> stall_pc -> pc_q, if_id_reg, ic_addr_q enables
\end{lstlisting}

Further down the same cone (cell \code{_22767_} onward) sit
\code{mem_wb_reg} (72 leaves) and \code{dbg_bp0_addr} /
\code{dbg_bp1_addr} (59 / 57 leaves): the combinational breakpoint
compare
\texttt{bp0\_\allowbreak{}hit\ =\allowbreak{}\ ctrl[0]\ \&\&\ (commit\_\allowbreak{}pc\_\allowbreak{}o\ =\allowbreak{}=\allowbreak{}\ dbg\_\allowbreak{}bp0\_\allowbreak{}addr)\ \&\&\ commit\_\allowbreak{}valid\_\allowbreak{}o}
(\code{rv32i_cpu_top.sv:794-797}) feeds \code{dbg_halt_req} (line 696),
which gates \code{ic_valid_o}. It is in the cone but is not the
latest-arriving input: it launches from flops and settles before the
stall network does.

\end{record}

\subsection{\texorpdfstring{3. Where the time goes (worst path,
\code{u_dcache.state_q[3]} -\textgreater{}
\code{u_icache.u_tag_sram/csb0}, 15.84 ns data
path)}{3. Where the time goes (worst path,  -\textgreater{} , 15.84 ns data path)}}\label{cr:defect:sky130-cpu-ss-setup-r51p:3-where-the-time-goes-worst-path-u-dcachestate-q3-u-icacheu-tag-sramcsb0-1584-ns-data-path}

\begin{record}

{\small\setlength{\tabcolsep}{4pt}

{
\begin{longtable}[]{@{}
  >{\raggedright\arraybackslash}p{\dimexpr(\linewidth-2\tabcolsep)*894/1000\relax}
  >{\raggedright\arraybackslash}p{\dimexpr(\linewidth-2\tabcolsep)*86/1000\relax}@{}}
\toprule\noalign{}
\begin{minipage}[b]{\linewidth}\raggedright
segment
\end{minipage} & \begin{minipage}[b]{\linewidth}\raggedright
ns
\end{minipage} \\
\midrule\noalign{}
\endhead
\bottomrule\noalign{}
\endlastfoot
clk-to-Q 0.64 + 4 fanout buffers (state\_\allowbreak{}q{[}3{]} fans out
to \textasciitilde65 sinks) & 2.0 \\
\code{dc_stall_o} state decode: and4bb + a31o + a211o (4-bit binary
state, 3 levels) & 2.5 \\
trunk: o21ai + a21oi + buffer (\code{mem_cache_stall} and its priority
gating) & 1.6 \\
gating into \code{tag_csb0}: a2111oi (1.59 ns, slew 1.5 ns) + 11
repeater buffers over a long route (\textasciitilde3.1 ns) + and4 + 4
buffers + a21oi + 4 buffers & \textasciitilde9.7 \\
\end{longtable}
}

}

Pre-PNR (no wires, ideal clocks) the same netlist already gives nom\_ss
WNS -1.78 ns / TNS -79 ns
(\code{12-openroad-staprepnr/nom_ss_100C_1v60/wns.max.rpt}). So roughly
half of the post-route -3.2 (nom\_ss) is logic depth and half is
wire/\allowbreak{}buffer/\allowbreak{}skew. The cure for the logic half
is RTL; the wire half is the PD agent\textquotesingle s.

For flop endpoints starting at \code{ex_mem_reg[0]} (-3.363): Q 1.21 +
buffers 1.1 + \code{need_mem_access}
o211ai/\allowbreak{}nor2/\allowbreak{}o21ai (3 levels, 3 repeaters)
reach the trunk at \textasciitilde7.1 ns after launch, then a 14-gate
tail (hazard priority, \code{ic_valid_o}, the icache
\code{valid_array[latch_index]} 256:1 read tree, \code{addr_mismatch},
\code{ic_stall_o}, enable muxes) reaches the flop at 22.8 ns.

\end{record}

\subsection{4. Consequences for the fix}\label{cr:defect:sky130-cpu-ss-setup-r51p:4-consequences-for-the-fix}

\begin{record}

\begin{itemize}
\tightlist
\item
  The user-prescribed change (retime icache tag-SRAM inputs for a full
  cycle) is a no-op in STA terms: they already get one.
\item
  What can be done in RTL, cycle-exactly, and attacks the measured cone:

  \begin{enumerate}
  \def\labelenumi{\arabic{enumi}.}
  \tightlist
  \item
    (landed) Cache FSM state encoding: put every
    unconditionally-stalling state in the \code{1xxx} half so
    \code{dc_stall_o} / \code{ic_stall_o} need \code{state_q[3]} instead
    of a 4-bit decode (removes \textasciitilde2 of the 2.5 ns decode
    segment). Pure encoding.
  \item
    (landed, Sky130 arm only) Remove \code{ic_valid_i} (and thereby the
    whole redirect/\allowbreak{}halt/\allowbreak{}stall cone) from
    \code{tag_csb0} / \code{data_csb0}: a read issued in \code{CS_IDLE}
    with
    \texttt{ic\_\allowbreak{}valid\_\allowbreak{}i\ =\allowbreak{}\ 0}
    is never consumed (dout is sampled only in \code{CS_SRAM_LATCH},
    entered only with
    \texttt{ic\_\allowbreak{}valid\_\allowbreak{}i\ =\allowbreak{}\ 1}),
    so the SRAM enable can depend on
    \texttt{state\_\allowbreak{}q\ =\allowbreak{}=\allowbreak{}\ CS\_\allowbreak{}IDLE}
    alone. This is the legitimate form of the user\textquotesingle s
    intent for the five SRAM pins.
  \end{enumerate}
\item
  Not expected to close ss at 75 MHz on its own. Remaining violators
  after 1 and 2 are the \code{ex_mem_reg} -\textgreater{}
  \code{need_mem_access} entrance and the 14-gate tail; both need
  pipeline-structure changes, which the user ruled out for this PR
  (section 6).
\item
  Context for a human decision (not acted on): the SoC runs this macro
  at 40 MHz (25 ns); the 75 MHz standalone ss failure does not affect
  SoC closure. The user decision for r51p is: RTL fix, do not relax the
  period, do not use SDC.
\end{itemize}

\end{record}

\subsection{5. What changed (RTL) and why it is cycle-exact}\label{cr:defect:sky130-cpu-ss-setup-r51p:5-what-changed-rtl-and-why-it-is-cycle-exact}

\begin{record}

\begin{itemize}
\tightlist
\item
  \src{rtl/mem/rv32i_cache_pkg.sv}: new \code{cache_state_t} encoding.
  \texttt{state\_\allowbreak{}q[3]\ =\allowbreak{}\ 1} for LATCH,
  HIT\_\allowbreak{}PENDING, REFILL, WRITEBACK,
  UNCACHED\_\allowbreak{}RD/\allowbreak{}WR, FLUSH\_\allowbreak{}SCAN,
  INVAL\_\allowbreak{}SCAN; IDLE (0000, reset), TAG\_\allowbreak{}CHECK,
  DONE have it clear. All arms share the package.
\item
  \code{rv32i_dcache.sv} / \code{rv32i_icache.sv}: \code{dc_stall_o} /
  \code{ic_stall_o} =
  \texttt{state\_\allowbreak{}q[3]\ ||\ (IDLE\ \&\&\ valid)\ ||\ (TAG\_\allowbreak{}CHECK\ \&\&\ !hit)}.
  Same truth table over the reachable states (icache never enters the
  D\$-only states).
\item
  \code{rv32i_icache.sv}, Sky130 arm only
  (\texttt{ifdef\ SRAM\_\allowbreak{}SKY130}):
  \texttt{idle\_\allowbreak{}read\_\allowbreak{}en\ =\allowbreak{}\ 1}
  so tag/data \code{csb0} in \code{CS_IDLE} do not depend on
  \code{ic_valid_i} / \code{ic_invalidate_i}. Why only Sky130: on ASAP7
  \code{csb0} feeds the SRAM clock gate (extra enabled cycles are extra
  dynamic power on a flow that was signed off as-is), and the FreePDK45
  / FPGA arms have no timing problem here. Why it is safe:
  \code{tag_dout_r} / \code{data_dout_r} capture \code{dout0} only in
  \code{CS_SRAM_LATCH}, which is entered only from a \code{CS_IDLE}
  cycle with
  \texttt{ic\_\allowbreak{}valid\_\allowbreak{}i\ =\allowbreak{}\ 1} and
  \texttt{ic\_\allowbreak{}invalidate\_\allowbreak{}i\ =\allowbreak{}\ 0};
  a read in any other IDLE cycle is never consumed and has no side
  effect. Writes (\code{tag_we_q} / \code{data_we_q}) still override.
\item
  No \code{ma7} / \code{ainf} exposure: no struct-member selects, no
  runtime-indexed muxes added (\code{check_struct_member_select.py} exit
  0).
\end{itemize}

Verification (all on this branch, measured):

{\footnotesize\setlength{\tabcolsep}{3pt}

{
\begin{longtable}[]{@{}
  >{\raggedright\arraybackslash}p{\dimexpr(\linewidth-2\tabcolsep)*490/1000\relax}
  >{\raggedright\arraybackslash}p{\dimexpr(\linewidth-2\tabcolsep)*490/1000\relax}@{}}
\toprule\noalign{}
\begin{minipage}[b]{\linewidth}\raggedright
gate
\end{minipage} & \begin{minipage}[b]{\linewidth}\raggedright
result
\end{minipage} \\
\midrule\noalign{}
\endhead
\bottomrule\noalign{}
\endlastfoot
\src{tb/tests/test_cache_stall_encoding.py} (new; RED before RTL, 3
failed / 2 passed) & 5/5 \\
icache / dcache / cache\_\allowbreak{}integration, Sky130 arm & 9/9,
8/8, 5/5 \\
icache / dcache / cache\_\allowbreak{}integration, default arm & 9/9,
8/8, 5/5 \\
\texttt{make\ -\allowbreak{}C\ sim\ rollup} (34 suites) & 265 tests, 0
fail \\
\code{test_mbist_cpu_macro} (collars still work) & 7/7 \\
RTL-vs-RTL differential, base (\#313) vs this branch, six
\src{tools/verif/gls/sky130} programs & 6/6 IDENTICAL, same total cycles
(466/\allowbreak{}566/\allowbreak{}966/\allowbreak{}4666/\allowbreak{}4666/\allowbreak{}4666) \\
Verilator lint \code{lint}, \code{lint_soc}, \code{lint_soc_sky130} & 0
errors; only existing warning classes (UNUSEDPARAM, DECLFILENAME,
MODDUP, PINCONNECTEMPTY on \code{dout1}) \\
\code{dft-scan-check}, \code{-asap7}, \code{-cpu-macro},
\code{-cpu-macro-asap7} & 0 violations each \\
\code{soc_all_ci} plain and \code{FPGA=1} & left to PR CI (PD flow
occupies the host) \\
mutation campaign, 15 mutants on the changed logic & 13 killed, 2
survived (both provably equivalent, below) \\
\end{longtable}
}

}

Mutants: dcache/\allowbreak{}icache stall (bit 3 dropped, idle ungated,
hit inverted) x6, encoding bit-3 or collision x4, idle read never
(Sky130 / tag / data) x3 all KILLED (by the structural test or by
\code{icache} / \code{dcache}). Survivors: M14
\texttt{idle\_\allowbreak{}read\_\allowbreak{}en\ =\allowbreak{}\ !ic\_\allowbreak{}invalidate\_\allowbreak{}i}
and M15 (functional arm drops \code{ic_valid_i}): both only add reads
that are never consumed, so they are equivalent mutants by the argument
above (M15 sits in code not compiled in the Sky130-arm suites).
\code{PASS_FLOOR} unchanged: no cocotb test was added or removed.

\end{record}

\subsection{6. PD handoff: expected versus residual (for the batched CPU
re-harden)}\label{cr:defect:sky130-cpu-ss-setup-r51p:6-pd-handoff-expected-versus-residual-for-the-batched-cpu-re-harden}

\begin{record}

Expected effect (INFERRED from the RUN\_\allowbreak{}2026-10-11 path
structure; not measured, no LibreLane run was done):

{\small\setlength{\tabcolsep}{4pt}

{
\begin{longtable}[]{@{}
  >{\raggedright\arraybackslash}p{\dimexpr(\linewidth-6\tabcolsep)*405/1000\relax}
  >{\raggedright\arraybackslash}p{\dimexpr(\linewidth-6\tabcolsep)*49/1000\relax}
  >{\raggedright\arraybackslash}p{\dimexpr(\linewidth-6\tabcolsep)*108/1000\relax}
  >{\raggedright\arraybackslash}p{\dimexpr(\linewidth-6\tabcolsep)*418/1000\relax}@{}}
\toprule\noalign{}
\begin{minipage}[b]{\linewidth}\raggedright
path class (RUN\_\allowbreak{}2026-10-11, max\_ss)
\end{minipage} & \begin{minipage}[b]{\linewidth}\raggedright
paths
\end{minipage} & \begin{minipage}[b]{\linewidth}\raggedright
worst before
\end{minipage} & \begin{minipage}[b]{\linewidth}\raggedright
expected after
\end{minipage} \\
\midrule\noalign{}
\endhead
\bottomrule\noalign{}
\endlastfoot
launched at \code{u_dcache.state_q} (decode head of the cone) & 407 &
-3.704 & roughly +1.3 to +2.0 ns: and4bb+a31o+a211o (2.46 ns) replaced
by one OR level \\
SRAM \code{csb0} endpoints (tag + 4 data) & 5 & -3.704 .. -2.3 & cone is
now icache \code{state_q} + \code{tag_we_q} only: expected to clear \\
launched at \code{ex_mem_reg} (via \code{need_mem_access}) & 329 &
-3.363 & UNCHANGED \\
launched at \code{ex1b_ex2_reg_q} (load-use) & 116 & -2.907 &
UNCHANGED \\
other (\code{mem_wb_reg}, \code{last_commit_insn_q}, ...) &
\textasciitilde140 & -2.98 .. -0.4 & unchanged \\
\end{longtable}
}

}

So the expected residual WNS is about -3.4 ns, set by \code{ex_mem_reg}
and \code{ex1b_ex2_reg_q} entrances, not by the state decode. TNS should
improve; WNS barely. Pre-PNR (no wires) the same netlist was already
-1.78 ns at nom\_ss, so about half of the residual is logic depth (RTL)
and half is wire/\allowbreak{}buffer/\allowbreak{}skew (placement; the
\code{csb0} path alone carried an 11-repeater, 3.1 ns chain).

Re-check after the re-harden:

\begin{enumerate}
\def\labelenumi{\arabic{enumi}.}
\tightlist
\item
  \code{max_ss} and \code{nom_ss}
  WNS/\allowbreak{}TNS/\allowbreak{}violator count versus -3.704 /
  -1677.6 / 996 (RUN\_\allowbreak{}2026-10-11). Expect WNS about -3.4,
  startpoint class \code{ex_mem_reg} / \code{ex1b_ex2_reg_q}.
\item
  Re-run \code{classify.py}-style startpoint/\allowbreak{}endpoint
  grouping: \code{u_dcache.state_q} should have dropped out of the top
  paths; SRAM \code{csb0} endpoints should be positive. The collars from
  \#313 add a \code{mem_collar} mux in series with each pin (about one
  gate); that is on the full-cycle input arc and is not expected to
  bind.
\item
  Check the icache/\allowbreak{}dcache SRAM \code{csb0} activity did not
  raise the Sky130 power/IR figures noticeably (Sky130 arm now reads in
  every IDLE cycle).
\item
  Hold: no new input-to-flop paths were added; the usual boundary-port
  hold set is unchanged in structure.
\item
  \code{check_cpu_netlist_census} / gate-vs-RTL
  (\code{run_netlist_proof.sh}) on the new netlist as for any re-harden.
\end{enumerate}

Not done in this PR (user decision: no pipeline restructuring; remains
open in r51p): registering \code{need_mem_access} /
\code{mem_misaligned} as precomputed flags in \code{rv32i_pipeline_ex2}
(D-side is a pure function of \code{ex1b_ex2_reg_q}, shallow), the
\code{ic_valid_o} / \code{dbg_halt_req} breakpoint-compare coupling
(\code{rv32i_cpu_top.sv:794-797}), and the load-use priority chain
behind \code{ex1b_ex2_reg_q}. SoC context unchanged: the SoC runs this
macro at 40 MHz and closes all 9 corners.

\end{record}


\section{PMU mutation campaign (\code{mutate_pmu_iso_ack.py})}\label{cr:defect:pmu-mutation-campaign}

\begin{record}

Date 2026-11-14. First recorded run of the campaign. Tool: \src{tools/verif/mutate_pmu_iso_ack.py}, Verilator 5.048 in the nix simulation shell, one run on one host, 27 mutants. A mutant is killed when any test of its suite fails. The suites are \code{pmu} for the \code{pmu.sv} mutants and \code{soc_pmu_iso_ack} for the \code{soc_top.sv} mutants. All statements are measured.

\end{record}

\begin{record}

Result: 25 killed, 2 survived (both marked expected, both equivalent), 0 unexpected survivors, 0 invalid. Expected survivors: \code{t06_counter_not_reset} (the state is on in reset, so the counter is cleared at the first edge after the reset and nothing reads it outside a wait) and \code{s05_ack_after_cpu_gating} (the only source of the acknowledge synchroniser is cleared by \code{cpu_core_rst_n}, so the acknowledge reads 0 in that window with either reset).

\end{record}

\begin{record}

New mutant \code{r01_cpu_gate_without_reset_term} removes \code{pmu_cpu_rst_n_cpu_sync} from \code{cpu_gated_clk_en}. It is killed by \code{test_cpu_clock_never_ticks_while_pmu_reset_asserted} and by \code{test_cpu_clock_never_ticks_while_forced_pmu_reset_and_clock_enabled}, and by no other test (2 of 13 tests fail). Before these tests, no test failed for this mutant. The first test measures a power cycle that the PMU sequences. The second asserts the PMU reset with the clock enable high, so the reset term is the only guard and the test does not depend on the order of the PMU actions.

\end{record}

\begin{record}

Mutant \code{t04_timeout_one_long}: at the 16 of \code{tb_pmu} the constant of the mutant truncates to 0 in the 4-bit counter, so the wait times out after one cycle. This agrees with the 9 of 25 tests that fail. At the default of 255 the mutant is one cycle late. The description in the script is corrected.

\end{record}

\begin{longtable}[]{@{}ll@{}}
\toprule
Mutant group & Tests that fail \\
\midrule
\endhead
\code{p01}..\code{p05} & 7, 20, 11, 4, 7 of 25 \\
\code{t01}..\code{t05} & 3, 3, 3, 9, 4 of 25 \\
\code{t07}..\code{t12} & 4, 3, 2, 1, 4, 21 of 25 \\
\code{g01}, \code{g02} & 3, 5 of 25 \\
\code{s01}..\code{s04} & 4, 6, 2, 1 of 13 \\
\code{c01}, \code{c02} & 1, 1 of 13 \\
\code{r01} & 2 of 13 \\
\bottomrule
\end{longtable}


\section{PMU order: clock stop before the clamp (\code{zftj})}\label{cr:defect:pmu-iso-order}

\begin{record}

Date 2026-11-14 to 2026-11-15. Beads \bead{zftj.7}, \bead{zftj.8} and \bead{zftj.10}. GH \#379, \#380, \#382 and \#386, fixed by PR \#408. The tests of PR \#388 confirmed the three defects of the PMU in CI before the change. \#379: after a timed-out power-up the acknowledge stayed high, and the next power-down left \code{DOM_PD_ISO} after one cycle with no handshake, and \code{STATUS} bit 3 was cleared at the start of that power-down. \#380: \code{cpu_gated_clk} had edges while the ISO\_CPU clamp was active, 4 in the power-down and 2 in the power-up at 1:1. \#382: a reset of the PMU in \code{DOM_OFF} released the reset, the clamp and the clock enable on one CPU clock edge. \#386: \code{ADDR_W} had no lower bound.

\end{record}

\begin{record}

Owner decisions of 2026-11-14. Option A of the design proposal for both domains: the power-down is save, clock stop, isolation with a wait for the acknowledge, reset; the power-up is reset release, restore, release of the isolation with a wait for the release, clock start. The CPU clamp control becomes \code{cpu_iso_clamp}, the clamp request \code{pmu_cpu_iso_en_sync} ANDed with the inverse of \code{cpu_gated_clk_en}, and its flop copy \code{cpu_iso_clamp_q} is the source of the acknowledge. Both waits accept only a fresh acknowledge (\code{ack_low_seen_q} for the power-down, the mirror \code{ack_high_seen_q} for the power-up). \code{STATUS} bit 3 becomes sticky, and a new write-1-to-clear register \code{STATUS_CLR} at 0x00C clears it. A guard rejects \code{ADDR_W} below 4, and the target \code{pmu_guard_neg} tests it and the guard of \code{ISO_ACK_TIMEOUT}. For \#382 the decision is to accept the behaviour and document it. The old rule DS-PMU-100, clamp before reset and clock stop, is rewritten. Option B, a second acknowledge for a stopped clock, was rejected.

\end{record}

\begin{record}

Correction. Earlier text in the books and in source comments called the place of the clock stop between the isolation and the reset an order of the power intent (UPF-mandated). The power intent \src{pnr/constraints/phase5_soc.upf} orders only save, isolate, reset and switch off, and switch on, reset release, restore and release isolation. It does not mention the clock gate. The place of the clock stop was a choice of \code{pmu.sv}. The record of the multi-clock work in the status history keeps the old words, and this record corrects them.

\end{record}

\begin{record}

Results in CI on the branch of PR \#408, period of \code{clk_i} 7 ns, measured. Suites: \code{pmu} 35 of 35, \code{soc_pmu_iso_ack} 13 of 13, \code{soc_pmu_multiclock} 12 of 12, \code{soc_integration} 9 of 9. The tests \code{test_p1_power_down_after_timed_out_power_up_needs_fresh_ack} and \code{test_p2_no_cpu_handshake_while_clamp_active} lost the mark \code{expect_fail} and pass. The two \code{test_p4_*} tests of a reset of the PMU in \code{DOM_OFF} stay \code{expect_fail} as evidence of the accepted limit. The sweep of \code{test_soc_pmu_multiclock} found no gated edge with the clamp active at the CPU clock periods 7, 3, 21 and 56 ns, at 13, 14, 10 and 8 offsets. The raw clamp request also overlapped no gated edge in the sweep, so only the injection tests discriminate the qualifier. A request forced with the gate open overlapped 8 to 11 gated edges at 7 ns and 3 ns as a raw request and none as a qualified clamp. Length of the waits \code{DOM_PD_ISO} and \code{DOM_PU_ISO}: 6 and 6 cycles at 1:1, 11 and 12 at 21 ns, 23 and 24 at 56 ns, 4 and 4 at 3 ns. Release monitor after a sequenced power-up, CPU edges of the clamp release, the reset release, the clock enable and the first gated edge: 9, 7, 15, 16 at 7 ns; 12, 8, 22, 23 at 3 ns; 6, 5, 10, 11 at 21 ns; 5, 5, 8, 9 at 56 ns. After a reset of the PMU in \code{DOM_OFF} the three releases share one edge at each period, and the first gated edge is one edge later.

\end{record}

\begin{record}

Mutation campaign of \src{tools/verif/mutate_pmu_iso_ack.py} on the new RTL, as reported in the comment of PR \#408: 39 mutants, 35 killed, 4 expected survivors, 0 unexplained survivors, 0 invalid. Expected survivors: \code{t06_counter_not_reset} (equivalent, the state is reset), \code{s05_ack_after_cpu_gating} (reset domain of the acknowledge synchroniser, structure only), \code{s06_ack_from_unregistered_clamp} (one cycle less latency, the glitch safety is structural) and \code{s07_ack_means_request_seen} (equivalent in a sequenced cycle, because the gate is already closed when the request arrives). The new mutants \code{q11_raw_clamp} and \code{q12_clamp_without_reset_term} are killed only by the injection tests. The mutant \code{r01_cpu_gate_without_reset_term} is killed only by \code{test_cpu_clock_never_ticks_while_forced_pmu_reset_and_clock_enabled} (1 of 13 tests fail). The sequenced test no longer kills it, because the clock now stops before the reset. Count note: the mutant table of the script has 40 entries, and the per-suite counts of the comment (28, 7 and 2 killed) add up to 37. This record keeps the totals of the comment and does not resolve the difference.

\end{record}

\begin{record}

Correction of the mutation counts above, 2026-11-15, measured. The campaign was run again in full on the head of the branch of PR \#408 (commit 049e7ae3), Verilator 5.048 in the nix simulation shell, one run on one host. The script \src{tools/verif/mutate_pmu_iso_ack.py} has 40 mutants. Result: 40 mutants, 36 killed, 4 expected survivors, 0 unexplained survivors, 0 invalid. By suite: \code{pmu} 28 mutants, 27 killed, 1 expected survivor (\code{t06_counter_not_reset}); \code{soc_pmu_iso_ack} 10 mutants, 7 killed (\code{s01} to \code{s04}, \code{c01}, \code{c02}, \code{r01_cpu_gate_without_reset_term}), 3 expected survivors (\code{s05}, \code{s06}, \code{s07}); \code{soc_pmu_multiclock} 2 mutants, 2 killed (\code{q11_raw_clamp}, \code{q12_clamp_without_reset_term}). The figures of the comment of the PR, 39 mutants, 35 killed and 28, 7 and 2 killed per suite, were wrong: the figure 28 was the number of mutants of the suite \code{pmu} and not the number killed, and the totals were one too low. The four expected survivors and their reasons are unchanged, with one correction: \code{s06_ack_from_unregistered_clamp} is not covered by \code{check_registered_outputs.py}, as the script said until this correction. That check reads only the outputs of \code{pmu} and not \code{cpu_iso_clamp_q} in \code{soc_top}. Nothing checks that the source of the acknowledge synchroniser is a flop. The reason is glitch safety only, as in the Design Specification and in the Technical Manual. A new comment on the PR replaces the earlier table.

\end{record}

\begin{record}

Reference firmware and the two software rules, 2026-11-15. Bead \bead{zftj.25}. GH \#380, PR \#412. PR \#408 added the rules \reqref{DS-PMU-182} (stop the bus traffic before a CPU power-down) and \reqref{DS-PMU-187} (clear \code{STATUS} bit 3 by a write to \code{STATUS_CLR}). The reference firmware \src{tb/cocotb/soc/pmu_cycle_fw/gen_pmu_cycle_hex.py} did not follow them. After it started the DMA copy of the power-down request, it stored 0 to the DMA control register, and it ran a loop that read, wrote and flushed the data cache in each pass. No firmware in the repository wrote \code{STATUS_CLR}. Measured, one run on one host: with the old firmware, a monitor of the CPU handshakes between the DMA start and \code{DOM_OFF} counted one write address, one write data and one write response, which belong to the store that cleared the DMA control bit. The monitor is in \code{test_pmu_live_power_cycle} and in its new companion. Both failed on the old firmware.

\end{record}

\begin{record}

Change of the same date. The firmware flushes the data cache and programs the DMA engine, stores 1 to the DMA control register in a line of the instruction cache that also holds its wait loop, and then runs only register instructions. After the wait loop it reads \code{STATUS}, records and clears a set bit 3, reads \code{STATUS} again and stores the result words. The counter of the live test is now the register \code{x9}, because the loop makes no memory access. The new test \code{test_pmu_live_power_cycle_clears_timeout_flag} holds the acknowledge input at 1 so that both waits time out. The new host test file \src{tb/tests/test_pmu_cycle_fw.py} has 5 tests. The test \code{test_pmu_live_power_cycle_monitor_sees_violating_firmware} runs an old-style image of the firmware as a control for the monitor. It also restores the one hit of the \code{load_use_ex1c} branch of \code{rv32i_hazard_unit}, which the SoC line coverage gate needs: the old loop gave it when the clock stopped with a load in flight, and the new firmware has no such state. Measured with the new firmware: \code{soc_pmu_multiclock} 14 of 14, \code{test_pmu_cycle_fw.py} and \code{test_pmu_confirm_fw.py} 10 of 10. A build of the firmware without the write to \code{STATUS_CLR} fails the new test on the value of \code{STATUS} bit 3. \code{PASS_FLOOR} in the two workflows goes from 925 to 927. Limit: in one build of the firmware with the write to \code{STATUS_CLR} removed and so the code after it moved by one word, the store of the end marker did not reach the SRAM in the run without a timeout. The cause is not found. A build with the code moved by three words did not show it. The committed firmware does not show it.

\end{record}

\section{Scan mode while the PMU holds the CPU off (\code{soc_top})}\label{cr:defect:scan-with-cpu-off}

\begin{record}

Date 2026-11-24. Bead \bead{zftj.24}. The SoC tests of scan mode with the CPU domain off checked only the domain reset. The clock gate, the clamp and the exit from scan mode had no test in that state. No RTL defect was found, and the RTL did not change.

\end{record}

\begin{record}

Tests, measured, local, one host. \code{test_soc_dft_scan.py} has three new tests, and the suite passes 16 of 16. Three temporary edits of \code{soc_top.sv} were run against the whole suite and then removed. With the test enable of \code{u_cpu_cg} tied to 0, three tests fail: \code{test_cpu_clock_gate_forced_open} and the first two new tests. With the test enable tied to 1, two tests fail: the first new test, at its check that the gate is closed in functional mode, and the scan-exit test. With a clamp that scan mode removes, the clamp test and the scan-exit test fail. A test without the wake write fails the scan-exit test at its wake check. No RTL edit was found that breaks only the wake. \code{soc_all_ci} has no failure: 994 tests pass. \code{PASS_FLOOR} goes from 955 to 958.

\end{record}

\section{APB bridge: destination reset races and their fix (\code{apb_cdc_bridge})}\label{cr:defect:apb-cdc-bridge-reset-race}

\begin{record}

Date 2026-11-15. Issues GH \#374, \#376 and \#378, beads \bead{zftj.13}, \bead{zftj.14} and \bead{zftj.15}, PR \#410. Issue \#375 (no timeout when the destination clock stops) is not part of this change.

\end{record}

\begin{record}

Root cause, one for the three issues. The destination kept its whole response state in flops that \code{m_rst_n_i} clears: \code{ack_q}, \code{resp_prdata_dq}, \code{resp_pslverr_dq} and the request synchroniser \code{u_req_sync}. The same reset was the only abort signal of the source, and that signal (\code{u_m_rst_status_sync}, a \code{cdc_sync} of the raw pin) was a sampled level. It had no order against \code{ack_s}, and a pulse between two edges of \code{s_clk_i} was not seen. Results: an error after a completed write (\#374, \code{test_b1_no_error_completion_after_completed_write}, 30 points), a repeated transfer after a short pulse (\#376, \code{test_b3_short_reset_pulse_makes_one_transfer}, 30 of 30 points), and an okay read with data 0 (\#378, \code{test_b5_read_never_completes_okay_with_wrong_data}, 33 points).

\end{record}

\begin{record}

Decision of the owner, 2026-11-14: Candidate A of the design proposal. The completion record (\code{ack_q}, \code{resp_prdata_dq}, \code{resp_pslverr_dq}) and \code{u_req_sync} reset from the source reset only, through the new \code{u_comp_rst_sync}. The destination reset status is \code{u_m_alive_sync} (\code{cdc_reset_sync}, asynchronous clear from \code{m_rst_sync_n}) and then \code{u_m_alive_dly_sync} (\code{cdc_sync}, \code{SYNC_STAGES_TO_S}+1 ranks, reset value ``not alive''). The abort thus comes one edge after the acknowledge of a completed transfer. The new \code{armed_q} lets the destination start only a request that arrived after it saw \code{req_m} low. The state \code{D_ACK} is removed. The source also leaves \code{S_DROP} when the destination is in reset or the request was not forwarded, and a request that finds a stale \code{ack_s} completes locally with an error. A write that completed before a destination reset returns an okay response. Rejected: Candidate B (record kept, no status delay, no arm gate, no escape from \code{S_DROP}), which loses to one cycle of synchroniser skew and can run an abandoned request after its error. Also rejected: a gray-coded destination state, and an okay response after every destination reset.

\end{record}

\begin{record}

Requirements. DS-CDC-103 goes from \code{SYNC_STAGES_TO_S}+2 back to \code{SYNC_STAGES_TO_S}+3. The bound of +2 came from \#377 and PR \#396, which matched the measured latency of the old RTL. The owner reversed it on purpose: the extra edge is the margin that lets a completed answer win against one cycle of skew. The measured count with the new RTL is 5 with \code{SYNC_STAGES_TO_S} of 2. DS-CDC-109 (minimum width of the destination reset) is retired, because the status no longer samples the pin. New: DS-CDC-117 (a completed transfer survives a destination reset), DS-CDC-118 (at most one transfer for a reset pulse of any width), DS-CDC-119 (no wait for an acknowledge after a destination reset), DS-CDC-120 (a stale acknowledge gives a local error). Changed: DS-CDC-100, DS-CDC-101, DS-CDC-104, DS-CDC-106, DS-CDC-110, DS-CDC-112 and DS-CDC-113. The local exception to DS-CDC-005 (a false ``not in reset'' status after a source reset) is gone.

\end{record}

\begin{record}

Tests, measured in the PR. \code{apb_cdc_bridge}: 42 tests, before 36. \code{test_b1_*}, \code{test_b3_*} and \code{test_b5_*} pass without \code{expect_fail}. Only \code{test_b2_stalled_destination_clock_completes_within_bound} (\#375) stays \code{expect_fail}. New suite \code{apb_cdc_bridge_skew} (settling-delay model on \code{u_ack_sync} and \code{u_m_alive_dly_sync}): 3 of 3 pass. \src{tb/tests/test_apb_cdc_bridge_reset_structure.py}: 7 pass. \code{soc_multiclock}, \code{soc_multiclock_reset}, \code{soc_pll_multiclock} and \code{soc_jtag}: 2, 14, 1 and 4 pass. The old ack-race test aligned on the fall of the acknowledge and was vacuous. It was rewritten. Not run: \code{soc_all_ci}, the SpyGlass flow of \code{dqc}, and timing on a routed netlist.

\end{record}

\begin{longtable}[]{@{}ll@{}}
\toprule
Mutant & Tests that fail \\
\midrule
\endhead
\code{ack_q} cleared by the destination reset & 10 \\
\code{u_req_sync} on the destination reset & 9 \\
No \code{armed_q} gate & 17 \\
One rank fewer in the status chain & 3, skew suite only \\
No escape from \code{S_DROP} & 2 \\
No check of a stale acknowledge & 1 \\
Abort before the acknowledge & 3, skew suite only \\
Status from the raw pin & not run, fails lint before simulation \\
\bottomrule
\end{longtable}

\begin{record}

Constraints. \src{pnr/constraints/phase5_soc_multiclock_check.sdc}: the matcher of the asynchronous pins also matches \code{m_alive_sync}, so these pins do not fall into the domain-wide maximum-delay fallback. \src{pnr/sky130/soc/constraints/sky130_soc_cdc_check.sdc}: the budget of the removed \code{u_m_rst_status_sync} is deleted, and section 10a already waives the new reset crossing. \src{tools/cdc/cdc_config.yml} and the signed ASAP7 constraint file are unchanged. The binding is proven by name-binding tests with negative controls only. The net \code{cpu_domain_rst_n} is an AND of reset-tree signals, and its term \code{cpu_core_rst_n} comes from the reset port and the PLL lock through a scan mux and not from a flop.

\end{record}

\section{Bus fabric protocol corners and stale \code{pslverr} waivers}\label{cr:defect:fabric-protocol-corners}

\begin{record}

Date 2026-11-24. Bead \bead{zftj.30}, GH \#403. The issue listed protocol corners of the bus fabric that no test covered, and two coverage waivers with a false reason. No RTL was changed.

\end{record}

\begin{record}

Bridge \code{axi4_to_axilite}, measured, local, one host, at the unit level (no wrapper joins the crossbar and the bridge). Twelve new tests: four pass (WLAST on the AWLEN beat, WLAST never asserted, an INCR burst, the strobes of a narrow write) and eight fail as expected. Measured failures: ``bridge hung after WLAST on beat 2 of 4: no B after 200 cycles''; ``bvalid at edge 8 before all 3 W beats were accepted''; FIXED write and read addresses 0x20004040, 0x20004044, 0x20004048, 0x2000404c; WRAP addresses 0x4008, 0x400c, 0x4010, 0x4014 where 0x4008, 0x400c, 0x4000, 0x4004 is correct; AxSIZE 0 words 0x4000, 0x4004, 0x4008, 0x400c where 0x4000, 0x4000, 0x4000, 0x4004 is correct, and the same fault for AxSIZE 1. Each test was also run without the mark and failed for the stated reason.

\end{record}

\begin{record}

Crossbar, measured. Nine new tests in \code{test_crossbar_errors.py}, all pass: a reset while each of W\_AW, W\_DATA, W\_RESP, R\_AR, R\_DATA, DW\_DATA, DW\_RESP and DR\_DATA is busy releases every engine; the slave IDs equal the master index for each master; a burst across a region boundary goes to the slave of its first beat; the DECERR engine ends on WLAST; an early WLAST closes the write; a missing WLAST holds back B; an early RLAST ends the read; a missing RLAST locks only that read engine. Mutants of the reset, \code{s_awid} and the DECERR WLAST logic made the new tests fail. \code{soc_all_ci}: exit 0, 1004 tests pass, 12 fail as expected, no failure. \code{PASS_FLOOR} goes from 955 to 968.

\end{record}

\begin{record}

Waivers. Three waivers in \code{coverage_waivers.txt} said that every APB slave ties \code{pslverr} to 0. The APB slots of the PLL and PLL2 are \code{apb_cdc_bridge}, which can answer PSLVERR. On the PLL slot the destination reset also holds \code{core_rst_n} low, so the bit cannot toggle. On the PLL2 slot the destination reset is \code{cpu_rst_n_i}, so a DMA or GPU access while it is low gets PSLVERR. The array waivers of \code{soc_top} and \code{soc_bus} now exclude element 6, and their reason text is correct. The 1-bit waiver of \code{soc_bus} \code{apb_m_pslverr} stays, with a correct reason, as an exception until a SoC test toggles it. \code{coverage_report.py --check-waivers}: 36 anchored, 0 stale. The instrumented coverage run was not done.

\end{record}

\begin{record}

Limits and open decisions. No test joins the crossbar and the bridge. The slave model cannot send a wrong BID or RID. Items 2 to 4 of the issue (no timeout in the fabric, the skid buffer of \code{u_gif_adapter} not reset with the GPU domain, ready outputs high in reset) are decisions of the owner and are not tested here.

\end{record}

\begin{record}

Date 2026-11-24. Owner decisions on GH \#403. The bridge \code{axi4_to_axilite} gets full support of FIXED, WRAP and narrow bursts, and a write burst ends on WLAST. Items 2 to 4 become new beads: \bead{zftj.34} (reset of the skid buffer of \code{u_gif_adapter} with the GPU domain), \bead{zftj.35} (ready outputs low in reset) and \bead{zftj.36} (a fabric timeout that answers SLVERR, design note first).

\end{record}

\begin{record}

RTL change in \code{axi4_to_axilite}. A shared function \code{next_addr} gives the next beat address for FIXED, INCR (aligned to the beat size, then plus the beat size) and WRAP (with a mask latched at AW or AR). AxSIZE above 2 is 4 bytes, and AxBURST 3 is INCR. A write ends on the first forwarded beat with WLAST. When the AWLEN count ends with no WLAST and the master still offers a beat, a new state \code{WS_SINK} accepts the extra beats without forwarding them and sends B after WLAST. When no beat is offered, B follows the count, as before. The latency of an INCR burst with WLAST on the AWLEN beat did not change. On the coarse netlist \code{\$shiftx} stays 0 and \code{\$pmux} goes from 7 to 11 (the case statements on size and burst, and the new state). Three coverage waivers had line numbers of the module in their reason text, and these were updated.

\end{record}

\begin{record}

Tests, measured, local, one host. The eight tests that were \code{expect_fail} now pass without the mark, with the same names. Seven new tests: early WLAST at five lengths, an early WLAST that keeps SLVERR or DECERR of a forwarded beat, late beats that are not forwarded, WRAP at 2, 4, 8 and 16 beats with AxSIZE 0 to 2 and three start offsets, a narrow INCR burst across a word, AxBURST 3 and AxSIZE 3 and 7. The suite has 29 tests and all pass. Each RTL mutant was caught: FIXED as INCR (3 tests fail), WRAP with no mask (4), a wrap window one beat short or long (4 each), a narrow step of 4 (5), no alignment after the first beat (3), WLAST ignored (5), a sink that forwards (2) and no sink (2). \code{soc_all_ci}: exit 0, 1019 tests pass by the count of the CI workflow, no failure. \code{PASS_FLOOR} goes from 968 to 983.

\end{record}

\begin{record}

Limits. No test joins the crossbar and the bridge. The case of a late beat that the master offers after the bridge looks at WVALID is not tested, and in that case B follows the count. The SoC masters use only the INCR path, so the SoC regression does not exercise FIXED, WRAP or narrow bursts.

\end{record}

\section{Bus access to a powered-off domain answers SLVERR (\code{axi_lite_interconnect}, \code{apb_cdc_bridge})}\label{cr:defect:off-domain-slverr}

\begin{record}

Date 2026-11-17. Bead \bead{zftj.26}. It fixes the off case of GH \#381 (bead \bead{zftj.11}: an access to the GPU while the GPU domain is off never ends) and of GH \#375 (bead \bead{zftj.16}: the APB bridge has no timeout when the destination clock stops). Before the change, the control slave of the GPU took the address of an access while it was held in reset and never answered, and the CPU waited for ever. A debug access to the CPU domain in flight at a power-down waited for the domain reset.

\end{record}

\begin{record}

Owner decisions of 2026-11-15, recorded on bead \bead{zftj.22}. An access to an off domain ends with SLVERR, and a read returns \code{OFF_RDATA}, 0xDEADBEEF. The bridge gets no timeout: a destination clock that stops while the domain is on still stalls a transfer, and this is a stated limit. The off levels come from \code{soc_top}, from registered outputs of the PMU. The answer comes from \code{axi_lite_interconnect} (input \code{slv_off_i}) and from \code{apb_cdc_bridge} (input \code{dest_off_i}). The PMU does not change. The test \code{test_p3_control_slave_ready_low_while_in_reset} stays \code{expect_fail} with a written waiver in \reqref{DS-PMU-185}, and rule 6 of \src{AGENTS.md} (``a reset value must not advertise availability'') is read for this case as ``no request is accepted and lost''. An error of the bridge from a reset alone keeps the read data 0 (\reqref{DS-CDC-116}). Only the off case returns \code{OFF_RDATA}.

\end{record}

\begin{record}

RTL change. \code{axi_lite_interconnect} has a new input \code{slv_off_i} with one bit for each slot and a parameter \code{OFF_RDATA}. Each engine latches the bit of the selected slot at the decode into a sticky flag (\code{w_off_q}, \code{r_off_q}) and sets it again in each later state. The effective level is the flag or the present bit. While it is high the engine accepts the address, sinks the write data and answers SLVERR itself, with \code{OFF_RDATA} for a read. It does not forward the access. A response of the slave that is already shown to the master is delivered once, unchanged. A fall of the level during a transaction does not resume forwarding. DECERR for an unmapped address has priority. \code{soc_bus} passes the new input \code{axil_slv_off_i}. \code{apb_cdc_bridge} has a new input \code{dest_off_i} and a parameter \code{OFF_RDATA}. The status \code{m_rst_n_s} is now \code{m_rst_n_s_raw} with \code{dest_off_i} cleared in, and \code{off_err_q} latches \code{dest_off_i} at SETUP. A request while the destination is off, or a rise of \code{dest_off_i} before the acknowledge, ends with PSLVERR and \code{OFF_RDATA}. A real acknowledge wins. In \code{soc_top}, \code{gpu_dom_off} is \code{~pmu_gpu_clk_en | pmu_gpu_iso_en | ~pmu_gpu_rst_n} and drives the GPU slot. \code{cpu_dom_off} is \code{~pmu_cpu_clk_en | pmu_cpu_iso_en | ~pmu_cpu_rst_n | pmu_cpu_iso_ack} and drives \code{dest_off_i} of \code{u_apb_dbg_cdc}. The PLL bridges tie the input to 0. Two later commits of the same bead keep the write response stable: the flag is also set in \code{W_RESP}, and a forwarded write response that is shown with BREADY low is locked.

\end{record}

\begin{record}

Tests, measured, local, one host. \code{soc_all_ci} exits 0. New tests: \code{test_axil_interconnect.py} has new tests of the off slot (off write and read, a rise of the level in each phase, a fall, a pulse, other slots, priority of DECERR, return to forwarding), \code{test_apb_cdc_bridge.py} has 46 tests, \code{apb_cdc_bridge_skew} has 4, and \code{test_soc_pmu_iso_ack.py} has the new \code{test_debug_access_while_cpu_off_completes_with_error}. \code{test_p3_gpu_access_while_off_completes_within_bound} of \code{test_soc_pmu_multiclock} lost the mark \code{expect_fail} and asserts SLVERR and 0xDEADBEEF. \code{PASS_FLOOR} in the two workflows goes from 927 to at least 947. Negative controls: with the off inputs tied to 0, the new tests fail. On the RTL of \code{origin/main} the P3 test fails with ``6 of 6 access(es) did not end within 4000 fabric cycles''.

\end{record}

\begin{record}

Limits. The bridge has no timeout, so a destination clock that stops while the domain is on and not in reset still stalls a debug access (\code{test_b2_stalled_destination_clock_completes_within_bound} stays \code{expect_fail} by decision). Only the debug APB path into the CPU domain has the guard. The CPU AXI bridge \code{u_cpu_axi_cdc} has no guard and no drain step (\reqref{DS-PMU-182}). A transaction that the interconnect cuts short can leave the slave with a request that is not answered, and only the domain reset that the PMU applies clears it. An error from the bridge for an access in flight at a power-down does not prove that the access had no effect in the CPU domain. New and changed requirements: \reqref{DS-PMU-111} to \reqref{DS-PMU-114}, \reqref{DS-CDC-121}, \reqref{DS-CDC-122}; \reqref{DS-PMU-182}, \reqref{DS-PMU-185} and \reqref{DS-CDC-116} changed.

\end{record}

\section{Ready outputs of the data fabric high in reset}\label{cr:defect:fabric-ready-in-reset}

\begin{record}

Date 2026-11-25. Bead \bead{zftj.35}, item 4 of GH \#403. Rule 6 of \src{AGENTS.md} requires a ready to be low while its domain is in reset. The owner approved an RTL fix on 2026-11-24. A survey of the ready outputs found readies high in reset in \code{axi4_crossbar} (AWREADY and ARREADY through the early accept and the DECERR path), \code{axi4_to_axilite}, \code{axil_to_apb}, \code{sram_controller}, \code{boot_rom} and \code{axilite_to_axi4} (ARREADY follows the crossbar, and BREADY was a constant 1). \code{axi_lite_interconnect} already conformed. \code{apb_interconnect} has no reset, and its PREADY while PSEL is low was not changed.

\end{record}

\begin{record}

Measured before the fix. Each unit test held reset with every VALID high for 8 cycles, and each listed ready was high in 8 of 8 cycles. At SoC level, with a fabric clock of 13 ns and a CPU clock of 3 ns and the fabric held in reset while the CPU ran, the readies of the boot ROM, the SRAM, the peripheral bridge, the APB bridge and the BREADY of the GPU adapter were high in 60 of 60 sampled cycles.

\end{record}

\begin{record}

RTL change: each listed ready is the AND of \code{rst_n} and its old term, with no added cycle. In \code{axilite_to_axi4} the request term \code{ar_pending} is gated with the reset of the skid buffer, so ARVALID to the crossbar and ARREADY to the GPU are both low in the fabric reset and in the GPU reset. This prevents a handshake on the crossbar that the GPU does not see. \code{m_axi_bready} is \code{rst_n}. New requirement \reqref{DS-FAB-005}.

\end{record}

\begin{record}

Tests, measured, local, one host. Eight tests that were \code{expect_fail} now pass, each with a control after the release and a passing guard for \code{axi_lite_interconnect}. New tests: a crossbar reset in the middle of a write burst, a read and a DECERR response; and two adapter tests that check one AR for each fetch after a reset with a held request. Each mutant that removes the gate of one module is caught, and so is an adapter that gates ARREADY but not ARVALID. \code{soc_pmu_iso_ack}: 32 of 32 pass. \code{soc_all_ci}: exit 0, 1082 pass. \code{PASS_FLOOR} goes from 1021 to 1043.

\end{record}

\begin{record}

Limits. The direct AND puts the reset on a combinational output path in the modules with an asynchronous reset, and no RDC tool ran on this host. The crossbar WREADY, BVALID and RVALID can be high in the first cycle of a reset during a transfer. The mutant of \code{sram_controller} changed one of its two variants only.

\end{record}

\section{GPU compute unit loses a fetch that is accepted after its issue cycle (\code{gpu_compute_unit})}\label{cr:defect:gpu-fetch-lost}

\begin{record}

Date 2026-11-25. Bead \bead{zftj.37}, found by a control test of bead \bead{zftj.34}. The compute unit recorded a fetch only when \code{warp_issue_i} and \code{if_arready_i} were high in the same cycle, and ARVALID was a pulse of one cycle. When the AR waited in the skid buffer of \code{u_gif_adapter}, the issue had no record, the warp stayed busy and the kernel did not finish. Measured before the fix: ``the GPU kernel did not finish within 3000 cycles after the held AR was accepted''. The CPU and the GPU share the SRAM, so contention on the crossbar can cause this.

\end{record}

\begin{record}

RTL change in \code{gpu_compute_unit}: the fetch is recorded on \code{fetch_go} (issue and no pipeline stall) whatever ARREADY is. A new flop \code{ar_pend_q} holds ARVALID and the recorded address until the handshake, and clears on the handshake, on \code{soft_clr_i} and on reset. The adapter is unchanged except for its header text. The check \code{if_ar_stuck} in \code{gpu_top} already assumed a held request. New requirement \reqref{DS-FAB-004}.

\end{record}

\begin{record}

Tests, measured, local, one host. New suite \code{tb/cocotb/gpu/test_fetch_handshake.py} on \code{gpu_top} (target \code{make -C sim gpu_fetch}, in \code{gpu_all}), five tests: a long ARREADY stall, periodic ARREADY with three latencies, 40 random kernels with random ARREADY and latency, eight warps under back-pressure, and a GPU soft reset while an AR waits. Each checks that ARVALID and ARADDR hold until the handshake, that AR handshakes, R beats and recorded fetches are equal, and the results against \code{gpu_ref_model.py}. In \code{test_soc_pmu_iso_ack.py} the test of the lost fetch passes without \code{expect_fail}, a control now counts only the R beat of the held AR, and a new test runs a kernel while the CPU loads the same SRAM (11 cycles of AR wait measured). Mutants in the compute unit, each failing all five GPU tests: the old pulse, \code{ar_pend_q} never set, the address not held, and \code{ar_pend_q} not cleared. \code{soc_all_ci}: 1060 pass, 0 fail. \code{gpu_all}: 91 pass. \code{PASS_FLOOR} 1019 to 1021, \code{CPU_PASS_FLOOR} 262 to 267.

\end{record}

\begin{record}

Limits. The mutants ran only against the GPU suite. No test has DMA contention. The clear of \code{ar_pend_q} on \code{soft_clr_i} is not reachable, because \code{if_ar_stuck} holds off the halt while an AR waits. A first version of the SoC contention test seemed to hang on a CPU read of the GPU status after the kernel ended (bead \bead{zftj.38}). The read completed. The harness sampled the one-cycle commit of the CPU at the fabric clock, which is slower than the CPU clock, so it missed the commit of the pass address. The harness now records each commit at the CPU clock, and the test polls the GPU status and checks that the read reaches the GPU slave.

\end{record}

\section{GPU fetch adapter keeps a stale fetch across a GPU power cycle (\code{axilite_to_axi4})}\label{cr:defect:gif-stale-fetch}

\begin{record}

Date 2026-11-25. Bead \bead{zftj.34}, from item 3 of GH \#403. The adapter \code{u_gif_adapter} of the GPU instruction fetch ran on \code{core_rst_n} only. Owner decisions of 2026-11-24: reset the AR skid buffer with the GPU domain reset, and also drain a read response that is in flight at a power-down, in one change. The adapter gets a second reset input. A separate defect that the tests found, a fetch that the GPU loses when the adapter accepts it later than the issue cycle, became bead \bead{zftj.37}.

\end{record}

\begin{record}

Measured on the RTL before the fix, SoC level, real PMU power cycle, CPU clock 3 ns and fabric clock 7 ns. PMU sequence of the GPU: clock gate, clamp, and the reset one cycle after the clamp. A held request: ``\code{m_axi_arvalid} still high in 5 cycles with \code{gpu_domain_rst_n} low''; the crossbar accepted it four cycles after the reset; its response reached the new GPU. A response in flight: ``crossbar SRAM read engine busy (state 2) in 18 of 18 cycles of the GPU-off window'', and the response reached the new GPU after the clamp released. A connection of the existing reset input to the GPU domain reset fixed only the three tests of the held request.

\end{record}

\begin{record}

RTL change. \code{axilite_to_axi4} has the inputs \code{gpu_rst_i} and \code{gpu_off_i}, the flag \code{rd_out_q} on \code{core_rst_n} and the flag \code{stale_q}. The skid buffer resets with the GPU reset. While a read is in flight and the GPU is off, or the read became stale, the adapter forces RREADY and masks RVALID to the GPU. \code{soc_top} drives the inputs from \code{gpu_domain_rst_n} and \code{gpu_dom_off}. New requirements \reqref{DS-PMU-115} and \reqref{DS-PMU-116}.

\end{record}

\begin{record}

Tests, measured, local, one host. \code{test_soc_pmu_iso_ack.py} has 31 tests: the five tests of the defect now pass without \code{expect_fail}; three controls; eight new tests (a new mode during a power-down, a response after the GPU is on again, a response while only the clock is gated, a second kernel after each case and after two power cycles, and both cases with the CPU clock slower than the fabric clock). Each RTL mutant was caught: skid reset on \code{rst_n} only, \code{rd_out_q} reset by the GPU reset, no RVALID mask, no RREADY force, no \code{stale_q}, and a drain on the reset only. The last two are caught only by the new tests of an early and a late response. \code{soc_all_ci} before the rebase on the crossbar changes: exit 0, 1007 tests pass. \code{PASS_FLOOR} goes from 1006 to 1019. The test of the lost fetch stays \code{expect_fail}.

\end{record}

\begin{record}

Limits. The drain starts at the clock gate. The PMU has no state that gates the GPU clock without a reset, and a deposit of a one-cycle gate pulse showed that such a state would lose the fetch. The tests force the read engine of the crossbar busy through a backdoor, because no SoC workload keeps it busy at the exact cycle.

\end{record}

\section{\code{axi4_crossbar} arbitration has no fairness rule}\label{cr:defect:crossbar-fixed-priority}

\begin{record}

Date 2026-11-18. Bead \bead{zftj.28}, GH \#401. The arbitration of \code{axi4_crossbar} is fixed priority: the lowest-index master wins. No test measured the grant order or the hold-off of a lower master. No RTL was changed. A fairness rule is a decision of the owner.

\end{record}

\begin{record}

Tests, measured, local, one host, on the three-master wrapper \code{tb_axi4_crossbar}. The new suite \code{test_crossbar_arbitration.py} has ten tests, five for reads and five for writes, and all pass. Requests in the same cycle are granted in ascending index for every subset of masters, both slaves, and AR and AW. The grant rule: while the engine of a slave is idle, the lowest-index master with xVALID high in that cycle wins, with the ready in the same cycle; the engine is idle for one cycle after RLAST or B; it keeps no record of arrival order. With master 0 re-requesting with no idle cycle in a window of 400 cycles, master 0 made 102 reads (82 writes), and masters 1 and 2 waited 402 and 406 cycles (404 and 409 for writes), until master 0 stopped. With one or two idle cycles between the transfers of master 0, the order after masters 1 and 2 arrive is 1, 0, 2. A master that holds one slave does not delay a master on the other slave. Negative controls on scratch copies of the RTL, not committed: highest index wins fails 8 of 10 tests, round robin fails 6 of 10, a block of slave 1 while slave 0 is busy fails the 2 cross-slave tests, and an ignored master 2 fails all 10.

\end{record}

\begin{record}

Limits. The wrapper has three masters, so master 2 stands for the lowest-priority master of the SoC. The tests use their own drivers, not the masters of the SoC. The existing \code{test_crossbar} tests have no timeout, so the mutants ran against the new suite only.

\end{record}

\begin{record}

Date 2026-11-24. Owner decision on GH \#401: the crossbar gets round-robin arbitration. RTL change in \code{axi4_crossbar}: each slave has a one-hot pointer for writes (\code{wr_ptr}) and one for reads (\code{rd_ptr}). Reset sets the pointer to master 0. The engine loads the pointer with the master one above the granted master when it takes a request in the idle state. The grant function \code{rr_grant} searches from the pointer and wraps to master 0. The early accept of the master face uses the same grant vectors as the engine. On the coarse netlist the counts of \code{\$shiftx} (12) and \code{\$pmux} (22) did not change. The latency of a single transfer did not change. New requirement \reqref{DS-FAB-003}.

\end{record}

\begin{record}

Tests, measured, local, one host, on \code{tb_axi4_crossbar}. The suite \code{test_crossbar_arbitration.py} now has twelve tests. Six tests that asserted fixed priority now compare each grant with a Python model of the round-robin arbiter. Two were renamed: \code{test_m0_back_to_back_*_hold_off_lower_masters} is now \code{test_m0_back_to_back_*_lower_masters_served_within_bound}, and \code{test_one_idle_cycle_lets_lower_master_in_*} is now \code{test_m0_idle_gap_does_not_reorder_grants_*}. New: \code{test_three_hogs_share_reads_evenly_within_bound} and \code{test_three_hogs_share_writes_evenly_within_bound}. Every test also checks that the order of grants on the master faces equals the order on the slave faces, and that no master waiting with xVALID high sees more than one grant of each other master. With master 0 re-requesting with no idle cycle, the worst wait of masters 1 and 2 was 8 cycles for reads and 10 cycles for writes, equal to the bound $(N-1) \times T$ with $T$ of 4 and 5 cycles. Before the change it was 402 to 409 cycles. Negative controls, each a temporary edit of the RTL: a pointer that does not move fails 8 of 34 crossbar tests, a pointer set to the granted master fails 10, a fixed-priority early accept fails 4 for each direction, a pointer that moves on each cycle fails 8, one pointer for both slaves fails 8, a pointer set two above the grant fails 8, and a search with no wrap deadlocks 17. \code{soc_all_ci} exits 0 with 1011 tests passed by the count of the CI workflow. \code{PASS_FLOOR} goes from 973 to 975.

\end{record}

\begin{record}

Limits. The arbiter counts grants, not cycles, so a master with long bursts gets more time on a slave than a master with short bursts. The bound was measured with three masters on the wrapper, not with the masters of the SoC. The SoC regression passed, so no SoC test depended on fixed priority, but the regression does not measure the grant order in the SoC.

\end{record}

\section{\code{axil_to_apb} drives PSTRB low on a read}\label{cr:defect:axil-to-apb-read-pstrb}

\begin{record}

Date 2026-11-17. Bead \bead{zftj.29}, GH \#402. The bridge \code{axil_to_apb} drove the strobes of the last write on \code{pstrb} during an APB read. AMBA APB4 requires \code{pstrb} low in a read transfer. No test checked \code{pstrb} in a read. Measured before the fix, one run on one host: after a write with \code{wstrb} 4'b0110, the SETUP and ACCESS phases of the next read had \code{pstrb} 4'b0110.

\end{record}

\begin{record}

RTL change. \code{pstrb} is \code{wstrb_q} when \code{is_write_q} is high and 0 otherwise. The output stays combinational. \code{is_write_q} goes low when the bridge captures a read address, so both phases of a read have \code{pstrb} at 0. A write still forwards \code{wstrb}.

\end{record}

\begin{record}

Tests, measured, local, one host. \code{test_axil_to_apb.py} has five new tests: write \code{pstrb} equals \code{wstrb} for five strobe values, a monitor control that sees the read phases, a read after reset, a read after a write, and a write and a read issued together. The last two failed on the old RTL (\code{pstrb} 4'b0110 in 2 of 2 read cycles) and pass on the new RTL. \code{soc_all_ci} has no failure: 991 tests pass and 4 fail as expected. \code{PASS_FLOOR} goes from 950 to 955.

\end{record}

\begin{record}

Limits. The bridge has no \code{pprot} output and ignores AxPROT. No slave in the APB tree uses a protection attribute. This is a stated limit by owner decision.

\end{record}

\section{\code{axi4_crossbar} loses a second request of one master}\label{cr:defect:crossbar-second-request}

\begin{record}

Date 2026-11-18. Bead \bead{zftj.27}, GH \#400. Two suspected defects of \code{axi4_crossbar} were examined by simulation. First: while a write or a read of one master waits for a slow slave, the crossbar accepts a second AW or AR of that master in the cycle that the slave raises its ready signal for the first. The slave gets only the first request. Second: with two reads or two writes of one master in flight to two different slaves, the master face, which has no ID, cannot steer the R beats or the W beats. No RTL was changed. An RTL fix is a decision of the owner.

\end{record}

\begin{record}

Tests, measured, local, one host. The new suite \code{test_crossbar_outstanding.py} has four tests marked \code{expect_fail} and four controls. Each defect test failed with a real mismatch: ``2 AW accepted at edges [1, 6] but 1 B seen'' for writes to one slave, no R beats for the second AR to one slave, ROM beats 2 and 3 lost behind the SRAM beats for two reads to two slaves, and the second slave given the W data of the first burst for two writes to two slaves. The controls send the same traffic with the next request after the B or RLAST of the last one, and they pass. The crossbar suite has 22 tests and no failure. \code{soc_all_ci} on the tree before the merge of \code{zftj.29}: exit 0, 990 tests pass, 8 fail as expected. \code{PASS_FLOOR} goes from 955 to 959.

\end{record}

\begin{record}

GPU masters. The data path (\code{memory_coalescer}) has one state machine that goes on only after R or B, so it has at most one transfer in flight. A passive monitor in \code{test_cpu_gpu_irq_pass} measured on crossbar master 2: 8 reads, 8 writes, a maximum of 1 in flight. The instruction fetch on master 1: 17 reads, a maximum of 1 in flight; the RTL reason is inferred from the stall logic of \code{gpu_compute_unit}. New rules: \reqref{DS-FAB-001} and \reqref{DS-FAB-002}.

\end{record}

\begin{record}

Limits. The monitor covers only the workload of \code{test_cpu_gpu_irq_pass}, and no negative control shows that it fires. The CPU cache arbiter and the DMA engine were not examined against \reqref{DS-FAB-001}. The tests drive only master 0.

\end{record}

\begin{record}

Date 2026-11-24. The owner approved an RTL fix: the crossbar holds the ready of a master face while that master has a transfer in flight. Each master has a write busy flag and a read busy flag. A flag is set by the request handshake of the master and cleared by its B handshake or by its R handshake with RLAST high. The flag masks the master valid in the slave arbiters, in the winner decode, in the early accept and in the unmapped decode. The arms of the slave engines that passed the ready of the slave to the master in the address phase are removed, because the early accept already completes the handshake. \reqref{DS-FAB-001} is now a rule of the crossbar, not of the master.

\end{record}

\begin{record}

Tests, measured, local, one host. The four defect tests of \code{test_crossbar_outstanding.py} pass on the new RTL. They are renamed to hold tests, without \code{expect_fail}, and they check that the ready stays low from the first request handshake up to the response. The 10 arbitration tests and the corner tests pass with the same simulation times as on the old RTL, so no cycle is added. Mutations: without the write busy mask or the read busy mask, the cross-slave hold test of that direction fails. With a flag cleared on the response cycle, the same test fails. With the old ready arm of a slave engine restored, both hold tests of that direction fail. The old RTL fails all four hold tests. \code{soc_all_ci} has no failure: 999 tests pass. \code{PASS_FLOOR} goes from 959 to 963.

\end{record}
